% Séries statistiques regroupées en classes — Mathématiques 1re Tech — Cours signé OnBuch+
% Programme officiel MINESEC. Compiler avec : tectonic N11-series-statistiques-regroupees-classes.tex
\newcommand{\DOCMATIERE}{Mathématiques}
\newcommand{\DOCNIVEAU}{1re Tech}
\newcommand{\DOCMODULE}{Mathématiques – classe de Première technique}
\newcommand{\DOCLECON}{N11}
\newcommand{\DOCTITRE}{Séries statistiques regroupées en classes}
% OnBuch+ — préambule « cours signés » ENRICHI (compilé avec tectonic / XeLaTeX)
% Système visuel « Pop » d'OnBuch+ : fond crème (jamais blanc), bordures encre
% épaisses, ombres « dures » décalées, titres Archivo Black en capitales.
% Version MATHÉMATIQUES : graphes (pgfplots/tikz), boîtes pédagogiques
% (définition, propriété, méthode, attention, le savais-tu, exercice résolu,
% exercice, TP, bilan), page de garde et signature OnBuch+ ancrée en bas.
%
% Macros attendues dans main.tex AVANT \input{preamble} :
%   \DOCMATIERE  \DOCNIVEAU  \DOCMODULE  \DOCLECON (numéro)  \DOCTITRE
\documentclass[11pt]{article}

\usepackage{fontspec}
\usepackage[french]{babel}
\usepackage[a4paper,margin=2cm,top=3.4cm,bottom=2.8cm,headheight=1.4cm,footskip=1.3cm]{geometry}
\usepackage{amsmath,amssymb,amsfonts}
\usepackage{mathtools} % \xmapsto, \coloneqq... souvent utilisés par les modèles
\usepackage{cancel} % simplifications barrées
% \abs{z} (module, valeur absolue) et \norm{u} (norme), utilisés par les modèles
\providecommand{\abs}{}\let\abs\relax\DeclarePairedDelimiter{\abs}{\lvert}{\rvert}
\providecommand{\norm}{}\let\norm\relax\DeclarePairedDelimiter{\norm}{\lVert}{\rVert}
\usepackage[table]{xcolor} % [table] : fournit aussi \rowcolors (alternance de lignes)
\usepackage{pagecolor}
\usepackage{tikz}
\usetikzlibrary{arrows.meta,positioning,calc,shapes.geometric,shapes.misc,shapes.arrows,decorations.pathmorphing,decorations.pathreplacing,decorations.markings,patterns,shadows,mindmap,trees,fit,backgrounds,matrix,intersections,angles,quotes}
\usepackage{pgfplots}
\usepackage{pgf-pie} % diagrammes circulaires \pie (statistiques)
\pgfplotsset{compat=1.18}
\usepgfplotslibrary{fillbetween,groupplots} % aires entre courbes (calcul intégral)
\usepackage{enumitem}
\usepackage{fancyhdr}
% [most] charge toutes les bibliothèques tcolorbox (listings, minted, raster,
% documentation...) et gonfle inutilement la mémoire TeX sur un document long
% (~300 boîtes) : on ne charge que ce qui est réellement utilisé.
\usepackage[skins,breakable]{tcolorbox}
\usepackage{graphicx}
\usepackage{colortbl}
\usepackage{booktabs}
\usepackage{tabularx}
\usepackage{diagbox} % case d'en-tête diagonale des tableaux à double entrée
% Colonnes extensibles centrée / gauche / droite (C, L, R), souvent utilisées par les modèles
\newcolumntype{C}{>{\centering\arraybackslash}X}
\newcolumntype{L}{>{\raggedright\arraybackslash}X}
\newcolumntype{R}{>{\raggedleft\arraybackslash}X}
\usepackage{array}
\usepackage{multicol}
\usepackage{siunitx}
% parse-numbers=false : accepte aussi \SI{2,5 \times 10^{-3}}{...}
\sisetup{locale=FR,output-decimal-marker={,},group-minimum-digits=4,parse-numbers=false}
\DeclareSIUnit\ohm{\ensuremath{\symup{\Omega}}} % U+2126 absent de la police
\usepackage{qrcode}
% \degree : commande fréquemment utilisée par les modèles pour un angle ou une
% température en dehors de \SI{}{} (non fournie par les paquets chargés).
\providecommand{\degree}{^{\circ}}
% \qed (fin de démonstration), utilisé par les modèles sans amsthm
\providecommand{\qed}{\hfill$\square$}
% \barre : double barre des valeurs interdites dans un tableau de variations (tkz-tab)
\providecommand{\barre}{\ensuremath{\|}}
% environnement proof (amsthm non chargé) utilisé par les modèles
\ifdefined\proof\else\newenvironment{proof}[1][Démonstration]{\par\noindent\textit{#1.}\ }{\qed\par}\fi
\usepackage{lastpage}
\usepackage{hyperref}
\hypersetup{hidelinks}

% ============================================================
% Palette « Pop » OnBuch+ — mêmes hex que la classe P (plus_ui.dart)
% ============================================================
\definecolor{poporange}{HTML}{F59321}
\definecolor{popdark}{HTML}{E07A0C}
\definecolor{popcream}{HTML}{FFF6E8}
\definecolor{popink}{HTML}{1C1714}
\definecolor{poppurple}{HTML}{7A5AE0}
\definecolor{popgreen}{HTML}{2E9E6A}
\definecolor{poppink}{HTML}{E0457B}
\definecolor{popblue}{HTML}{2F7DE1}
\definecolor{popgold}{HTML}{C98A12}
\definecolor{popmuted}{HTML}{8A7F76}
\definecolor{popline}{HTML}{EADFD2}
\definecolor{popgray}{HTML}{8A7F76}
\colorlet{popgrey}{popgray}
% Alias tolérants (le modèle nomme parfois les couleurs différemment)
\colorlet{popwhite}{white}
\colorlet{popblack}{popink}
\colorlet{popred}{poppink}
\colorlet{popyellow}{popgold}
% Teintes claires pour fonds de tableaux / zones
\colorlet{popblueL}{popblue!10!white}
\colorlet{poporangeL}{poporange!14!white}
\colorlet{popgreenL}{popgreen!12!white}
\colorlet{poppurpleL}{poppurple!12!white}
\colorlet{poppinkL}{poppink!12!white}
\colorlet{popgoldL}{popgold!14!white}

\pagecolor{popcream}

% ============================================================
% Typographie
% ============================================================
\defaultfontfeatures{Ligatures=TeX}
\setmainfont{PlusJakartaSans-Medium.ttf}[Path=../../../fonts/,BoldFont=PlusJakartaSans-Bold.ttf]
% Maths en sans-serif (Fira Math) avec chiffres et lettres droites de Plus
% Jakarta, pour que les formules se fondent dans le texte.
\usepackage[math-style=ISO,bold-style=ISO]{unicode-math}
\setmathfont{FiraMath-Regular.otf}[Path=../../../fonts/]
\setmathfont{PlusJakartaSans-Medium.ttf}[Path=../../../fonts/,range={up/num,up/{Latin,latin},\mathup}]
\usepackage{icomma} % virgule décimale sans espace en mode math
% Caractères mathématiques Unicode absents des polices de texte (Plus Jakarta,
% Archivo Black) : rendus par leur équivalent mathématique, en texte comme en
% formule — sinon ils s'affichent en carré vide (ex. « Arithmétique dans ℤ »).
\usepackage{newunicodechar}
\newunicodechar{ℝ}{\ensuremath{\mathbb{R}}}
\newunicodechar{ℤ}{\ensuremath{\mathbb{Z}}}
\newunicodechar{ℂ}{\ensuremath{\mathbb{C}}}
\newunicodechar{ℕ}{\ensuremath{\mathbb{N}}}
\newunicodechar{ℚ}{\ensuremath{\mathbb{Q}}}
\newunicodechar{𝔻}{\ensuremath{\mathbb{D}}}
\newunicodechar{α}{\ensuremath{\alpha}}
\newunicodechar{β}{\ensuremath{\beta}}
\newunicodechar{γ}{\ensuremath{\gamma}}
\newunicodechar{δ}{\ensuremath{\delta}}
\newunicodechar{ε}{\ensuremath{\varepsilon}}
\newunicodechar{θ}{\ensuremath{\theta}}
\newunicodechar{λ}{\ensuremath{\lambda}}
\newunicodechar{μ}{\ensuremath{\mu}}
\newunicodechar{σ}{\ensuremath{\sigma}}
\newunicodechar{τ}{\ensuremath{\tau}}
\newunicodechar{φ}{\ensuremath{\varphi}}
\newunicodechar{ψ}{\ensuremath{\psi}}
\newunicodechar{ω}{\ensuremath{\omega}}
\newunicodechar{Σ}{\ensuremath{\Sigma}}
% majuscules produites par \MakeUppercase dans les titres (α-aminés -> Α)
\newunicodechar{Α}{\ensuremath{\alpha}}
\newunicodechar{Β}{\ensuremath{\beta}}
\newunicodechar{Γ}{\ensuremath{\Gamma}}
\newunicodechar{Δ}{\ensuremath{\Delta}}
\newunicodechar{Ε}{\ensuremath{\varepsilon}}
\newunicodechar{Θ}{\ensuremath{\Theta}}
\newunicodechar{Λ}{\ensuremath{\Lambda}}
\newunicodechar{Μ}{\ensuremath{\mu}}
\newunicodechar{Τ}{\ensuremath{\tau}}
\newunicodechar{Φ}{\ensuremath{\Phi}}
\newunicodechar{Ψ}{\ensuremath{\Psi}}
\newunicodechar{Ω}{\ensuremath{\Omega}}
\newunicodechar{↦}{\ensuremath{\mapsto}}
\newunicodechar{∈}{\ensuremath{\in}}
\newunicodechar{∉}{\ensuremath{\notin}}
\newunicodechar{∧}{\ensuremath{\wedge}}
\newunicodechar{⇔}{\ensuremath{\Leftrightarrow}}
\newunicodechar{⟺}{\ensuremath{\Longleftrightarrow}}
\newunicodechar{⇒}{\ensuremath{\Rightarrow}}
\newunicodechar{⟹}{\ensuremath{\Longrightarrow}}
\newunicodechar{∘}{\ensuremath{\circ}}
\newunicodechar{∪}{\ensuremath{\cup}}
\newunicodechar{∩}{\ensuremath{\cap}}
\newunicodechar{≡}{\ensuremath{\equiv}}
\newunicodechar{⊂}{\ensuremath{\subset}}
\newunicodechar{⊥}{\ensuremath{\perp}}
\newunicodechar{∃}{\ensuremath{\exists}}
\newunicodechar{∀}{\ensuremath{\forall}}
\newunicodechar{∛}{\ensuremath{\sqrt[3]{\,}}}
\newunicodechar{∜}{\ensuremath{\sqrt[4]{\,}}}
\newunicodechar{⃗}{\ensuremath{{}^{\to}}}
\newunicodechar{ⁿ}{\ifmmode^{n}\else\textsuperscript{\ensuremath{n}}\fi}
\newunicodechar{ˣ}{\ifmmode^{x}\else\textsuperscript{\ensuremath{x}}\fi}
\newunicodechar{ᵇ}{\ifmmode^{b}\else\textsuperscript{\ensuremath{b}}\fi}
\newunicodechar{ʳ}{\ifmmode^{r}\else\textsuperscript{\ensuremath{r}}\fi}
\newunicodechar{ᵘ}{\ifmmode^{u}\else\textsuperscript{\ensuremath{u}}\fi}
\newunicodechar{ʸ}{\ifmmode^{y}\else\textsuperscript{\ensuremath{y}}\fi}
\newunicodechar{ᵃ}{\ifmmode^{a}\else\textsuperscript{\ensuremath{a}}\fi}
\newunicodechar{ᵉ}{\ifmmode^{e}\else\textsuperscript{\ensuremath{e}}\fi}
\newunicodechar{ᵏ}{\ifmmode^{k}\else\textsuperscript{\ensuremath{k}}\fi}
\newunicodechar{ᵖ}{\ifmmode^{p}\else\textsuperscript{\ensuremath{p}}\fi}
\newunicodechar{ᴺ}{\ifmmode^{N}\else\textsuperscript{\ensuremath{N}}\fi}
\newunicodechar{ᶜ}{\ifmmode^{c}\else\textsuperscript{\ensuremath{c}}\fi}
\newunicodechar{ʰ}{\ifmmode^{h}\else\textsuperscript{\ensuremath{h}}\fi}
\newunicodechar{ᵠ}{\ifmmode^{q}\else\textsuperscript{\ensuremath{q}}\fi}
\newunicodechar{ₙ}{\ifmmode_{n}\else\textsubscript{\ensuremath{n}}\fi}
\newunicodechar{ₐ}{\ifmmode_{a}\else\textsubscript{\ensuremath{a}}\fi}
\newunicodechar{ᵢ}{\ifmmode_{i}\else\textsubscript{\ensuremath{i}}\fi}
\newunicodechar{ᵣ}{\ifmmode_{r}\else\textsubscript{\ensuremath{r}}\fi}
\newunicodechar{ₖ}{\ifmmode_{k}\else\textsubscript{\ensuremath{k}}\fi}
\newunicodechar{ᵦ}{\ifmmode_{\beta}\else\textsubscript{\ensuremath{\beta}}\fi}
\newunicodechar{ₓ}{\ifmmode_{x}\else\textsubscript{\ensuremath{x}}\fi}
\newfontfamily\popdisplay{ArchivoBlack-Regular.ttf}[Path=../../../fonts/]
\newfontfamily\poplogo{SpaceGrotesk-Bold.ttf}[Path=../../../fonts/]
\newfontfamily\popsemi{PlusJakartaSans-SemiBold.ttf}[Path=../../../fonts/]
\newfontfamily\popxbold{PlusJakartaSans-ExtraBold.ttf}[Path=../../../fonts/]
\newfontfamily\popxboldls{PlusJakartaSans-ExtraBold.ttf}[Path=../../../fonts/,LetterSpace=3.0]

\setlist[enumerate]{leftmargin=1.7em,itemsep=5pt,topsep=4pt}
\setlist[itemize]{leftmargin=1.7em,itemsep=4pt,topsep=4pt,label={\color{poporange}\textbullet}}
\setlength{\parindent}{0pt}
\setlength{\parskip}{5pt}
\renewcommand{\arraystretch}{1.25}

% pgfplots : style Pop par défaut
\pgfplotsset{
  every axis/.append style={axis line style={popink,line width=1pt},tick style={popink},
    grid style={popline,line width=0.5pt},label style={font=\small},tick label style={font=\footnotesize}},
  popcurve/.style={poporange,line width=1.8pt,no markers,smooth},
}
% Même style disponible dans un \draw TikZ simple (hors pgfplots)
\tikzset{popcurve/.style={poporange,line width=1.8pt}}

% ============================================================
% Pill / squiggle
% ============================================================
\newcommand{\poppill}[3]{%
  \tikz[baseline=-0.55ex]\node[draw=popink,line width=1.2pt,rounded corners=2.6pt,%
    fill=#2,inner xsep=6.5pt,inner ysep=3pt]{%
    \popxboldls\fontsize{7}{7}\selectfont\color{#3}\MakeUppercase{#1}};%
}
\newcommand{\popsquiggle}[1]{%
  \tikz[baseline=-0.4ex]{
    \draw[#1,line width=2.6pt,line cap=round]
      plot [smooth] coordinates {(0,0.09) (0.34,0.01) (0.68,0.21) (1.02,0.01) (1.36,0.10)};
  }%
}
% Mot-clé surligné (marqueur orange sous le texte)
\newcommand{\cle}[1]{{\popxbold\color{popdark}#1}}

% ============================================================
% En-tête / pied
% ============================================================
\pagestyle{fancy}
\fancyhf{}
\renewcommand{\headrulewidth}{0pt}
\renewcommand{\footrulewidth}{0pt}
\lhead{\raisebox{-0.15ex}{\poplogo\fontsize{13}{13}\selectfont\color{popink}OnBuch\color{poporange}+}}
\rhead{\poppill{\DOCMATIERE\ \textbullet\ \DOCNIVEAU\ \textbullet\ Leçon \DOCLECON}{white}{popink}}
\lfoot{\popxbold\fontsize{7.6}{7.6}\selectfont\color{popmuted}\MakeUppercase{Cours signé OnBuch+}\ \textbullet\ {\popsemi\DOCTITRE}}
\rfoot{\tikz[baseline=-0.6ex]\node[rounded corners=8.5pt,fill=popink,minimum height=17pt,minimum width=17pt,inner xsep=5pt,inner sep=0]{\popxbold\fontsize{8}{8}\selectfont\color{popcream}\thepage\,/\,\pageref*{LastPage}};}

% ============================================================
% Titres
% ============================================================
\newcommand{\chaptitre}[2]{%
  \begin{tcolorbox}[enhanced,colback=poporange,colframe=popink,boxrule=2.6pt,arc=11pt,
      drop shadow=popink,left=15pt,right=15pt,top=12pt,bottom=11pt,before skip=3pt,after skip=18pt]
    \poppill{#2}{popink}{popcream}\par\vspace{9pt}
    {\raggedright\hyphenpenalty=10000\exhyphenpenalty=10000\popdisplay\fontsize{22}{24}\selectfont\color{popcream}\MakeUppercase{#1}\par}
  \end{tcolorbox}
}
% Section numérotée : pastille orange avec le numéro + titre Archivo
\newcounter{popsec}
\newcommand{\coursec}[1]{%
  \stepcounter{popsec}\setcounter{popsub}{0}%
  \par\vspace{16pt}\noindent
  \begin{minipage}{\linewidth}
  \tikz[baseline=-0.7ex]\node[draw=popink,line width=1.6pt,fill=poporange,rounded corners=4pt,minimum size=22pt,inner sep=0]{\popdisplay\fontsize{12}{12}\selectfont\color{popcream}\thepopsec};%
  \hspace{8pt}{\popdisplay\fontsize{15}{17}\selectfont\color{popink}\MakeUppercase{#1}}\par
  \vspace{4pt}\noindent{\color{poporange}\rule{36pt}{3.6pt}}
  \end{minipage}\par\nopagebreak\vspace{8pt}
}
\newcounter{popsub}
\newcommand{\courssub}[1]{%
  \stepcounter{popsub}%
  \par\vspace{9pt}\noindent{\popxbold\fontsize{11.5}{13}\selectfont\color{popink}{\color{poporange}\thepopsec.\thepopsub}\hspace{6pt}#1}\par\nopagebreak\vspace{4pt}
}

% ============================================================
% Boîtes pédagogiques — même recette : fond blanc, bordure épaisse colorée,
% ombre dure encre, badge plein. Titre optionnel [#1].
% ============================================================
\tcbset{popbase/.style={breakable,enhanced,colback=white,boxrule=2.3pt,arc=9pt,
  shadow={3pt}{-3pt}{0pt}{popink},left=14pt,right=14pt,top=11pt,bottom=11pt,before skip=11pt,after skip=15pt,
  fontupper=\popsemi\fontsize{10.6}{14.5}\selectfont\color{popink},
  fontlower=\popsemi\fontsize{10.6}{14.5}\selectfont\color{popink}}}
% Coche dessinée (le glyphe ✓ n'existe pas dans les polices du document)
\newcommand{\popcheck}[1]{\tikz[baseline=-0.5ex]\draw[#1,line width=1.6pt,line cap=round,line join=round](-0.13,0.0)--(-0.04,-0.09)--(0.13,0.1);}
\newcommand{\popboxhead}[3]{\poppill{#1}{#2}{white}\ifx\relax#3\relax\else\hspace{6pt}{\popxbold\fontsize{10.6}{12}\selectfont\color{popink}#3}\fi\par\vspace{7pt}}

\newtcolorbox{aretenir}[1][]{popbase,colframe=popgreen,before upper={\popboxhead{À retenir}{popgreen}{#1}}}
\newtcolorbox{exemplebox}[1][]{popbase,colframe=poporange,before upper={\popboxhead{Exemple}{poporange}{#1}}}
\newtcolorbox{definition}[1][]{popbase,colframe=popblue,before upper={\popboxhead{Définition}{popblue}{#1}}}
\newtcolorbox{propriete}[1][]{popbase,colframe=poppurple,before upper={\popboxhead{Propriété}{poppurple}{#1}}}
\newtcolorbox{methode}[1][]{popbase,colframe=popgold,before upper={\popboxhead{Méthode}{popgold}{#1}}}
\newtcolorbox{attention}[1][]{popbase,colframe=poppink,before upper={\popboxhead{Attention}{poppink}{#1}}}
\newtcolorbox{experience}[1][]{popbase,colframe=popblue,colback=popblueL,before upper={\popboxhead{Expérience}{popblue}{#1}}}
% « Le savais-tu ? » : carte violette pleine (contexte, culture, Cameroun)
\newtcolorbox{savaistu}[1][]{popbase,colback=poppurple,colframe=popink,
  fontupper=\popsemi\fontsize{10.4}{14.2}\selectfont\color{white},
  before upper={\poppill{Le savais-tu\,?}{white}{poppurple}\ifx\relax#1\relax\else\hspace{6pt}{\popxbold\color{white}#1}\fi\par\vspace{7pt}}}
% Exercice résolu : énoncé en haut, \tcblower puis la solution sur fond crème
\newcounter{popexo}\newcounter{popexr}
\newtcolorbox{exoresolu}[1][]{popbase,colframe=poporange,colbacklower=popcream,
  segmentation style={popink,line width=1.2pt,dashed},
  before upper={\stepcounter{popexr}\popboxhead{Exercice résolu \thepopexr}{poporange}{#1}},
  before lower={\poppill{Solution}{popink}{popcream}\par\vspace{6pt}}}
% Exercice (non résolu en place) — la correction est dans la partie Corrigés
\newtcolorbox{exercice}[1][]{popbase,colframe=popink,
  before upper={\stepcounter{popexo}\popboxhead{Exercice \thepopexo}{popink}{#1}}}
% Corrigé
\newtcolorbox{corrige}[1][]{popbase,colframe=popgreen,colback=popgreenL,
  before upper={\popboxhead{Corrigé}{popgreen}{#1}}}
% Objectifs / prérequis / bilan
\newtcolorbox{objectifs}{popbase,colframe=popink,colback=poporangeL,
  before upper={\popboxhead{Objectifs de la leçon}{popink}{}}}
\newtcolorbox{prerequis}{popbase,colframe=popink,colback=popblueL,
  before upper={\popboxhead{Prérequis}{popblue}{}}}
\newtcolorbox{bilan}{popbase,colframe=popink,colback=white,boxrule=2.8pt,
  before upper={\popboxhead{Fiche bilan}{poporange}{L'essentiel à connaître pour l'examen}}}
% Figure encadrée avec légende
\newtcolorbox{popfigure}[1][]{enhanced,breakable=false,colback=white,colframe=popline,boxrule=1.4pt,arc=8pt,
  left=10pt,right=10pt,top=10pt,bottom=8pt,before skip=10pt,after skip=12pt,halign=center,
  lower separated=false,halign lower=center,
  fontlower=\popsemi\fontsize{9}{11.5}\selectfont\color{popmuted},
  #1}
\newcommand{\legende}[1]{\par\vspace{6pt}{\popsemi\fontsize{9}{11.5}\selectfont\color{popmuted}#1}}

% ============================================================
% Signature OnBuch+
% ============================================================
\newcommand{\poplogoinline}[1]{{\poplogo\fontsize{#1}{#1}\selectfont\color{popink}OnBuch\color{poporange}+}}
\newcommand{\signatureonbuch}{%
  \begin{tcolorbox}[enhanced,colback=popink,colframe=popink,boxrule=2.2pt,arc=10pt,
      drop shadow=poporange,left=14pt,right=14pt,top=10pt,bottom=10pt,before skip=0pt,after skip=0pt]
    \tikz[baseline=-0.7ex]\node[circle,fill=popgreen,draw=popcream,line width=1.3pt,minimum size=22pt,inner sep=0]{\popcheck{white}};%
    \hspace{9pt}\begin{minipage}[c]{0.62\linewidth}
      {\popxbold\fontsize{12}{13}\selectfont\color{popcream}Signé\ \textbullet\ {\poplogo OnBuch\color{poporange}+}}\par\vspace{2pt}
      {\raggedright\popsemi\fontsize{8}{10.5}\selectfont\color{popline}\MakeUppercase{Équipe pédagogique OnBuch+}\\\MakeUppercase{Conforme au programme officiel MINESEC}\par}
    \end{minipage}\hfill
    {\poplogo\fontsize{22}{22}\selectfont\color{popcream}OnBuch\color{poporange}+}
  \end{tcolorbox}%
}

% ============================================================
% Page de garde
% #1 titre · #2 matière · #3 niveau · #4 module · #5 n° leçon · #6 durée ·
% #7 description / objectifs (texte ou itemize)
% ============================================================
\newcommand{\pagegarde}[7]{%
  \thispagestyle{empty}
  \newgeometry{a4paper,margin=1.7cm,top=1.6cm,bottom=1.4cm}%
  \begin{tikzpicture}[remember picture,overlay]
    \fill[poporange!18] ([xshift=-3.2cm,yshift=-3.6cm]current page.north east) circle (4.6cm);
    \fill[poppurple!14] ([xshift=2.4cm,yshift=7.6cm]current page.south west) circle (3.8cm);
  \end{tikzpicture}%
  {\poplogo\fontsize{30}{30}\selectfont\color{popink}OnBuch\color{poporange}+}\hfill
  \raisebox{0.6ex}{\poppill{Cours signé \textbullet\ Premium}{popink}{popcream}}\par
  \vspace{1pt}\popsquiggle{poppurple}\par
  \vspace{8pt}
  \begin{tcolorbox}[enhanced,colback=poporange,colframe=popink,boxrule=2.8pt,arc=13pt,
      drop shadow=popink,left=18pt,right=18pt,top=12pt,bottom=13pt,after skip=10pt]
    \poppill{#2 \textbullet\ #3}{popink}{popcream}\hspace{5pt}\poppill{#4}{white}{popink}\par\vspace{8pt}
    {\popdisplay\fontsize{14}{14}\selectfont\color{popink}LEÇON #5}\par\vspace{4pt}
    {\raggedright\hyphenpenalty=10000\exhyphenpenalty=10000\popdisplay\fontsize{25}{27}\selectfont\color{popcream}\MakeUppercase{#1}\par}
    \vspace{8pt}
    {\popxbold\fontsize{9.5}{9.5}\selectfont\color{popink}\MakeUppercase{Durée conseillée : #6}}
  \end{tcolorbox}
  \begin{tcolorbox}[enhanced,colback=white,colframe=popink,boxrule=2.2pt,arc=10pt,
      drop shadow=popink,left=16pt,right=16pt,top=10pt,bottom=10pt,after skip=8pt]
    \poppill{Dans cette leçon}{poppurple}{white}\par\vspace{5pt}
    \popsemi\fontsize{9.3}{12.6}\selectfont\color{popink}#7
  \end{tcolorbox}
  \noindent\begin{minipage}[t]{0.487\linewidth}
    \begin{tcolorbox}[enhanced,colback=popgreen,colframe=popink,boxrule=2.2pt,arc=10pt,
        drop shadow=popink,left=11pt,right=11pt,top=10pt,bottom=10pt,height=3.1cm,valign=center]
      \tcbox[colback=white,colframe=popink,boxrule=1.3pt,arc=5pt,boxsep=0pt,nobeforeafter,
        left=3pt,right=3pt,top=3pt,bottom=3pt,valign=center]{\qrcode[height=1.9cm]{https://play.google.com/store/apps/details?id=com.luvvix.onbuch&hl=fr}}%
      \hspace{8pt}\begin{minipage}[c]{0.5\linewidth}
        {\popdisplay\fontsize{11}{12.5}\selectfont\color{white}\MakeUppercase{Télécharge OnBuch}}\par\vspace{4pt}
        {\raggedright\popsemi\fontsize{8.4}{11}\selectfont\color{white}Gratuit \textbullet\ Google Play\\Tuteur IA, annales, concours, résultats\par}
      \end{minipage}
    \end{tcolorbox}
  \end{minipage}\hfill
  \begin{minipage}[t]{0.487\linewidth}
    \begin{tcolorbox}[enhanced,colback=poppurple,colframe=popink,boxrule=2.2pt,arc=10pt,
        drop shadow=popink,left=11pt,right=11pt,top=10pt,bottom=10pt,height=3.1cm,valign=center]
      \tikz[baseline=-0.7ex]\node[circle,fill=white,draw=popink,line width=1.3pt,minimum size=1.6cm,inner sep=0]{\popdisplay\fontsize{24}{24}\selectfont\color{poppurple}+};%
      \hspace{8pt}\begin{minipage}[c]{0.6\linewidth}
        {\popdisplay\fontsize{11}{12.5}\selectfont\color{white}\MakeUppercase{Passe à OnBuch+}}\par\vspace{4pt}
        {\raggedright\popsemi\fontsize{8.4}{11}\selectfont\color{white}Tous les cours signés, Tuteur IA illimité, fiches PDF — onglet OnBuch+ de l'app\par}
      \end{minipage}
    \end{tcolorbox}
  \end{minipage}
  \vspace{8pt}
  \popcontenu
  \vfill
  \signatureonbuch
  \restoregeometry
  \newpage
}

% Rangée « Ce cours contient » (page de garde)
\newcommand{\poptile}[3]{% #1 couleur #2 gros texte #3 légende
  \begin{tcolorbox}[enhanced,colback=#1,colframe=popink,boxrule=2pt,arc=9pt,shadow={2.5pt}{-2.5pt}{0pt}{popink},
      left=6pt,right=6pt,top=9pt,bottom=9pt,halign=center,valign=center,height=1.75cm,width=\linewidth,nobeforeafter]
    {\popdisplay\fontsize{12.5}{13}\selectfont\color{white}\MakeUppercase{#2}}\par\vspace{3pt}
    {\centering\popsemi\fontsize{7.8}{9.5}\selectfont\color{white}#3\par}
  \end{tcolorbox}}
\newcommand{\popcontenu}{%
  \par\noindent{\popxboldls\fontsize{7.5}{7.5}\selectfont\color{popmuted}CE COURS SIGNÉ CONTIENT}\par\vspace{7pt}
  \noindent\begin{minipage}[t]{0.235\linewidth}\poptile{popblue}{Cours}{complet, illustré, pas à pas}\end{minipage}\hfill
  \begin{minipage}[t]{0.235\linewidth}\poptile{poporange}{Méthodes}{et exercices résolus}\end{minipage}\hfill
  \begin{minipage}[t]{0.235\linewidth}\poptile{poppink}{Exercices}{type examen + corrigés}\end{minipage}\hfill
  \begin{minipage}[t]{0.235\linewidth}\poptile{popgreen}{Fiche bilan}{l'essentiel à retenir}\end{minipage}\par}

% Sommaire
\newtcolorbox{sommaire}{enhanced,colback=white,colframe=popink,boxrule=2.2pt,arc=10pt,
  drop shadow=popink,left=18pt,right=18pt,top=13pt,bottom=13pt,before skip=6pt,after skip=20pt,
  before upper={\poppill{Sommaire}{poppurple}{white}\par\vspace{9pt}\popsemi\fontsize{10.6}{14.5}\selectfont\color{popink}}}

% Signature de fin de cours — poussée tout en bas de la dernière page
\newcommand{\findecours}{%
  \par\vfill\nopagebreak
  \begin{center}\popsquiggle{poporange}\end{center}\vspace{4pt}
  \signatureonbuch
}
\AtBeginDocument{\def\llbracket{\mathopen{[\mkern-3.5mu[}}\def\rrbracket{\mathclose{]\mkern-3.5mu]}}} % intervalles d'entiers (glyphe ⟦ absent de la police)
% Glyphes absents de Fira Math : redéfinis à partir de glyphes disponibles
\AtBeginDocument{%
  \def\setminus{\mathbin{\backslash}}\let\smallsetminus\setminus
  \def\vdots{\mathord{\vbox{\baselineskip3.2pt\lineskiplimit0pt\kern4pt\hbox{.}\hbox{.}\hbox{.}}}}%
  \def\ddots{\mathinner{\mkern1mu\raise6.5pt\hbox{.}\mkern2mu\raise3.6pt\hbox{.}\mkern2mu\raise0.7pt\hbox{.}\mkern1mu}}%
  \def\triangle{\mathord{\scalebox{0.9}{$\symup{\Delta}$}}}%
  \def\nabla{\mathord{\raisebox{\depth}{\scalebox{1}[-1]{$\symup{\Delta}$}}}}%
  \def\checkmark{\popcheck{popdark}}%
  \def\star{\mathbin{\ast}}\def\bigstar{\mathord{\ast}}%
  \def\diamond{\mathbin{\scalebox{0.6}{$\Diamond$}}}%
  \def\bigcup{\mathop{\mathchoice{\raisebox{-0.3em}{\scalebox{1.7}{$\cup$}}}{\scalebox{1.25}{$\cup$}}{\cup}{\cup}}\displaylimits}%
  \def\bigcap{\mathop{\mathchoice{\raisebox{-0.3em}{\scalebox{1.7}{$\cap$}}}{\scalebox{1.25}{$\cap$}}{\cap}{\cap}}\displaylimits}%
  \def\lesseqqgtr{\mathrel{\lessgtr}}\def\bigoplus{\mathop{\oplus}\displaylimits}%
}
\providecommand{\ssi}{si et seulement si} % abréviation fréquente
\AtBeginDocument{\providecommand{\wideparen}{\overparen}} % arc de cercle

%% FIGFIX-BEGIN v4 — correctifs globaux des figures (docs/onbuch-generation/tools/figfix)
\usepackage{adjustbox}\usepackage{etoolbox}
% toute figure TikZ/pgfplots plus large que la ligne est réduite à la largeur disponible
\BeforeBeginEnvironment{tikzpicture}{\begin{adjustbox}{max width=\linewidth}}
\AfterEndEnvironment{tikzpicture}{\end{adjustbox}}
% légendes pgfplots lisibles par-dessus les courbes
\pgfplotsset{every axis/.append style={legend style={fill=white,fill opacity=0.92,text opacity=1,draw=black!15,inner sep=3pt}}}
% étiquettes d'arêtes (syntaxe edge["..."]) avec fond blanc
\tikzset{every edge quotes/.append style={fill=white,inner sep=1.2pt}}
%% FIGFIX-END
\begin{document}

\pagegarde{Séries statistiques regroupées en classes}{Mathématiques}{1re Tech}{Mathématiques – classe de Première technique}{N11}{2 séances (fiche de progression harmonisée)}{Cette leçon apprend à organiser une série statistique volumineuse en classes d'amplitude donnée, à calculer effectifs et fréquences cumulés croissants ou décroissants, et à représenter graphiquement la série (histogramme, polygones et courbes cumulatives) afin d'interpréter des données concrètes de la vie courante.
\vspace{4pt}
\begin{multicols}{2}\small
\begin{itemize}
\item À la fin de cette leçon, je sais regrouper une série statistique brute en classes d'amplitude donnée.
\item À la fin de cette leçon, je sais calculer les effectifs et les fréquences de chaque classe.
\item À la fin de cette leçon, je sais construire les effectifs cumulés croissants et décroissants et les fréquences cumulées correspondantes.
\item À la fin de cette leçon, je sais représenter une série groupée en classes par un histogramme.
\item À la fin de cette leçon, je sais tracer le polygone des effectifs et les courbes cumulatives croissante et décroissante.
\item À la fin de cette leçon, je sais lire et interpréter un tableau ou un graphique statistique pour répondre à une question concrète (combien d'individus dépassent un seuil, etc.).
\end{itemize}\end{multicols}}

\begin{sommaire}
\begin{multicols}{2}
\begin{enumerate}
\item Regroupement en classes
\item Effectifs et fréquences cumulés
\item Histogramme d'une série groupée en classes
\item Courbes cumulatives
\item Étude complète d'une situation concrète
\item Activité d'intégration
\item Méthodes et exercices résolus
\item Exercices
\item Corrigés des exercices
\item Fiche bilan
\end{enumerate}
\end{multicols}
\end{sommaire}

\begin{objectifs}
\begin{itemize}
\item À la fin de cette leçon, je sais regrouper une série statistique brute en classes d'amplitude donnée.
\item À la fin de cette leçon, je sais calculer les effectifs et les fréquences de chaque classe.
\item À la fin de cette leçon, je sais construire les effectifs cumulés croissants et décroissants et les fréquences cumulées correspondantes.
\item À la fin de cette leçon, je sais représenter une série groupée en classes par un histogramme.
\item À la fin de cette leçon, je sais tracer le polygone des effectifs et les courbes cumulatives croissante et décroissante.
\item À la fin de cette leçon, je sais lire et interpréter un tableau ou un graphique statistique pour répondre à une question concrète (combien d'individus dépassent un seuil, etc.).
\item À la fin de cette leçon, je sais rédiger et justifier une démarche statistique complète, du tableau brut à la conclusion.
\item À la fin de cette leçon, je sais déterminer le centre d'une classe et l'utiliser pour préparer le calcul de la moyenne (leçon suivante).
\end{itemize}
\end{objectifs}

\begin{prerequis}
\begin{itemize}
\item Notions de statistique de 4e/3e : effectif, effectif total, fréquence, pourcentage
\item Lecture et construction de tableaux et de diagrammes (bâtons, circulaires)
\item Calculs sur les fractions, les proportions et les pourcentages
\item Repérage dans le plan et lecture de coordonnées
\item Notion d'intervalle : bornes, amplitude, appartenance d'un nombre à un intervalle
\end{itemize}
\end{prerequis}

\newpage

\coursec{Regroupement en classes}

Au lycée technique de Nkongsamba, le proviseur vient de recevoir les notes de mathématiques des $120$ élèves de Première technique au premier trimestre. Devant lui : une longue liste de $120$ nombres compris entre $0$ et $20$, écrits les uns à la suite des autres. Impossible, à l'œil nu, de répondre aux questions qui l'intéressent : combien d'élèves ont la moyenne ? Les notes sont-elles plutôt faibles ou plutôt bonnes ? Faut-il organiser un soutien ? Les données sont trop nombreuses pour être exploitées telles quelles.

Il prend alors une décision simple et efficace : \textbf{regrouper les notes en classes}, c'est-à-dire en intervalles de $4$ points : de $0$ à $4$, de $4$ à $8$, de $8$ à $12$, etc. En quelques minutes, la liste informe devient un tableau clair, sur lequel on pourra calculer des effectifs, des fréquences, puis des effectifs cumulés.

\textbf{Problématique :} comment organiser une série statistique volumineuse pour en tirer des informations utiles à la décision ? C'est tout l'objet de cette leçon, et cette première section pose les fondations : le regroupement en classes.

\courssub{Rappels : population, individu, caractère statistique}

Avant de regrouper, rappelons le vocabulaire de base de la statistique descriptive, déjà rencontré dans les classes précédentes.

\begin{definition}[Population, individu, caractère]
\begin{itemize}
\item La \cle{population} est l'ensemble sur lequel porte l'étude statistique.
\item Un \cle{individu} est un élément de cette population.
\item Le \cle{caractère} (ou variable statistique) est la propriété que l'on observe sur chaque individu.
\end{itemize}
On distingue trois types de caractères :
\begin{itemize}
\item \cle{qualitatif} : le caractère n'est pas mesurable par un nombre (couleur d'un produit, filière d'un élève, marque d'un moteur, type de défaut) ;
\item \cle{quantitatif discret} : le caractère prend des valeurs numériques isolées, dénombrables (nombre d'enfants par famille, nombre de pièces défectueuses par lot, nombre de pannes par mois) ;
\item \cle{quantitatif continu} : le caractère peut prendre n'importe quelle valeur dans un intervalle (taille, masse, durée, tension électrique, débit, résistance, revenu, note sur 20 avec demi-points).
\end{itemize}
\end{definition}

\begin{exemplebox}[Dans notre situation]
La population est l'ensemble des $120$ élèves de Première technique du lycée de Nkongsamba. Un individu est un élève. Le caractère étudié est la note de mathématiques du premier trimestre : c'est un caractère quantitatif, que l'on traitera comme continu (les notes, avec les demi-points et quarts de points, peuvent prendre de très nombreuses valeurs entre $0$ et $20$).
\end{exemplebox}

\courssub{Pourquoi regrouper en classes ?}

Quand le caractère est discret avec peu de valeurs (par exemple le nombre d'enfants par famille, de $0$ à $6$), on présente la série valeur par valeur. Mais dans deux situations, cette présentation devient inutilisable :

\begin{itemize}
\item \textbf{Les valeurs sont très nombreuses.} Avec $120$ notes, on peut avoir $40$ ou $50$ valeurs distinctes : un tableau « valeur / effectif » aurait des dizaines de lignes et ne dégagerait aucune tendance.
\item \textbf{Le caractère est continu.} Une taille de \SI{1,72}{m}, une tension de \SI{220,4}{V}, un salaire de \SI{152300}{FCFA} : deux individus n'ont presque jamais exactement la même valeur. Compter « valeur par valeur » n'a alors aucun sens.
\end{itemize}

\begin{aretenir}[Principe du regroupement en classes]
Regrouper une série en \cle{classes}, c'est partager l'intervalle des valeurs du caractère en un petit nombre d'intervalles (les classes), puis compter combien d'individus tombent dans chaque intervalle. On perd un peu de précision (on ne connaît plus les valeurs exactes), mais on gagne un tableau lisible et directement exploitable.
\end{aretenir}

\begin{savaistu}[Les recensements au Cameroun]
Lors des recensements de la population camerounaise, le BUCREP (Bureau Central des Recensements et des Études de Population) ne publie pas l'âge exact de chaque habitant : il présente les âges \textbf{regroupés en classes de $5$ ans} : $[0\,;5[$, $[5\,;10[$, $[10\,;15[$, etc. C'est exactement la méthode que tu apprends dans cette leçon, appliquée à des millions d'individus ! Ces tableaux permettent à l'État de prévoir le nombre d'écoles, d'hôpitaux ou d'emplois à créer.
\end{savaistu}

\courssub{Vocabulaire : classe, bornes, amplitude, centre}

\begin{definition}[Classe, bornes, amplitude, centre]
Une \cle{classe} est un intervalle $[a\,;b[$ (ou $[a\,;b]$ pour la dernière) dans lequel on regroupe les valeurs du caractère.
\begin{itemize}
\item Les nombres $a$ et $b$ sont les \cle{bornes} de la classe : $a$ est la borne inférieure, $b$ la borne supérieure.
\item L'\cle{amplitude} de la classe est la différence $b - a$.
\item Le \cle{centre} de la classe (ou point milieu) est le milieu de l'intervalle :
\[ c = \frac{a+b}{2}. \]
\end{itemize}
\end{definition}

\begin{exemplebox}[Lecture sur une classe de notes]
La classe $[8\,;12[$ a pour borne inférieure $8$, pour borne supérieure $12$, pour amplitude $12-8=4$ et pour centre :
\[ c = \frac{8+12}{2} = 10. \]
Le centre jouera un rôle essentiel dans la suite de l'unité : pour calculer une moyenne sur une série groupée, on fera « comme si » toutes les valeurs de la classe étaient égales au centre.
\end{exemplebox}

\begin{attention}[Convention d'écriture des classes : le piège des bornes]
Une note située exactement sur une borne ne doit être comptée \textbf{qu'une seule fois}. Pour cela, on adopte la convention des intervalles \textbf{fermés à gauche, ouverts à droite} :
\[ [0\,;4[\,,\ [4\,;8[\,,\ [8\,;12[\,,\ [12\,;16[\,,\ [16\,;20]. \]
Ainsi la note $8$ appartient à la classe $[8\,;12[$ et \emph{pas} à la classe $[4\,;8[$. Seule la \textbf{dernière classe} est fermée des deux côtés, $[16\,;20]$, pour que la note maximale $20$ soit bien comptée. Erreur fréquente au Probatoire : écrire des classes du type $[0\,;4]$, $[4\,;8]$, ... et compter deux fois les valeurs $4$, $8$, $12$... Le total ne vaut alors plus $N$ et toute la suite est fausse.
\end{attention}

\courssub{Choix du nombre de classes et de l'amplitude}

Il n'existe pas de règle absolue, mais des habitudes solides :

\begin{itemize}
\item On choisit en pratique \textbf{entre $5$ et $10$ classes} : trop peu de classes écrasent l'information, trop de classes rendent le tableau aussi illisible que la liste de départ.
\item Au Probatoire, l'\textbf{amplitude est constante} : toutes les classes ont la même largeur (sauf mention contraire de l'énoncé, qui la donne alors).
\item On calcule l'amplitude en divisant l'\textbf{étendue} (différence entre la plus grande et la plus petite valeur) par le nombre de classes souhaité, puis on arrondit à une valeur simple (multiple de $1$, $2$, $5$, $10$, $0,1$, $0,2$, $0,5$ selon l'échelle).
\end{itemize}

\begin{exemplebox}[Choix de l'amplitude pour les notes de Nkongsamba]
Les notes s'étendent de $0$ à $20$, l'étendue vaut donc $20$. En visant $5$ classes, on obtient $\dfrac{20}{5} = 4$ : une amplitude de $4$ points, simple et pratique. On peut aussi retenir $6$ classes si l'on veut isoler la dernière tranche ; ici l'énoncé du proviseur fixe des intervalles de $4$ points, ce qui donne les classes $[0\,;4[$, $[4\,;8[$, $[8\,;12[$, $[12\,;16[$, $[16\,;20]$.
\end{exemplebox}

La figure suivante visualise ce découpage sur une droite graduée.

\begin{popfigure}
\begin{center}
\begin{tikzpicture}[scale=0.62,>=Stealth]
\draw[->,thick,popdark] (-0.5,0) -- (21.5,0);
\foreach \x in {0,4,8,12,16,20}{
  \draw[thick,popdark] (\x,0.15) -- (\x,-0.15) node[below]{\x};
}
% classes
\foreach \a/\b/\c in {0/4/popblue,4/8/popgreen,8/12/poporange,12/16/poppurple,16/20/poppink}{
  \draw[very thick,\c] (\a,0.55) -- (\b,0.55);
  \fill[\c] (\a,0.55) circle (2.2pt);
  \draw[\c,fill=white] (\b,0.55) circle (2.2pt);
}
\fill[poppink] (20,0.55) circle (2.2pt);
% labels classes
\node[popblue,above] at (2,0.65) {\small $[0\,;4[$};
\node[popgreen,above] at (6,0.65) {\small $[4\,;8[$};
\node[poporange,above] at (10,0.65) {\small $[8\,;12[$};
\node[poppurple,above] at (14,0.65) {\small $[12\,;16[$};
\node[poppink,above] at (18,0.65) {\small $[16\,;20]$};
% centres
\foreach \x/\c in {2/popblue,6/popgreen,10/poporange,14/poppurple,18/poppink}{
  \fill[\c] (\x,-1.1) circle (2.2pt);
  \node[\c,below] at (\x,-1.15) {\small centre $\x$};
}
% amplitude
\draw[<->,thick,popgold] (8,-2.3) -- (12,-2.3) node[midway,below]{\small amplitude $= 4$};
\end{tikzpicture}
\end{center}
\legende{Découpage de l'intervalle $[0\,;20]$ en $5$ classes d'amplitude $4$. Les points pleins marquent les bornes incluses, les points creux les bornes exclues ; sous chaque classe figure son centre.}
\end{popfigure}

\courssub{Effectifs et fréquences d'une classe}

\begin{definition}[Effectif, effectif total, fréquence]
\begin{itemize}
\item L'\cle{effectif} d'une classe, noté $n_i$, est le nombre d'individus dont la valeur appartient à cette classe.
\item L'\cle{effectif total}, noté $N$, est le nombre total d'individus :
\[ N = n_1 + n_2 + \dots + n_k = \sum_{i=1}^{k} n_i \quad (k = \text{nombre de classes}). \]
\item La \cle{fréquence} de la classe numéro $i$ est le quotient
\[ f_i = \frac{n_i}{N}. \]
Elle est comprise entre $0$ et $1$. Multipliée par $100$, elle donne la \cle{fréquence en pourcentage} : $f_i \times 100\,\%$.
\end{itemize}
\end{definition}

\begin{propriete}[Somme des fréquences]
La somme des fréquences de toutes les classes vaut toujours $1$ (soit $100\,\%$) :
\[ f_1 + f_2 + \dots + f_k = \frac{n_1 + n_2 + \dots + n_k}{N} = \frac{N}{N} = 1. \]
Cette propriété sert de \textbf{test de vérification} : si la somme de tes fréquences ne vaut pas $1$ (ou $100\,\%$, aux arrondis près), il y a une erreur de comptage ou de calcul.
\end{propriete}

\begin{definition}[Classe modale]
La \cle{classe modale} est la classe qui possède le plus grand effectif (ou la plus grande fréquence). Elle correspond à la « valeur la plus fréquente » de la série groupée.
\end{definition}

\courssub{Méthode de construction du tableau}

\begin{methode}[Construire un tableau de regroupement en classes en 5 étapes]
\begin{enumerate}
\item \textbf{Choisir l'amplitude} (ou le nombre de classes) : calculer l'étendue, la diviser par le nombre de classes visé (entre $5$ et $10$), arrondir à une valeur simple.
\item \textbf{Écrire les classes} avec la convention $[a\,;b[$, en fermant la dernière classe des deux côtés.
\item \textbf{Calculer le centre} de chaque classe : $c_i = \dfrac{a_i+b_i}{2}$.
\item \textbf{Compter} les individus de chaque classe (dépouillement de la liste, par exemple avec des traits de pointage), puis reporter les effectifs $n_i$.
\item \textbf{Vérifier} que la somme des effectifs vaut $N$, puis calculer les fréquences $f_i = \dfrac{n_i}{N}$ et contrôler que leur somme vaut $1$.
\end{enumerate}
\end{methode}

\courssub{Exemple guidé : les 120 notes du proviseur}

\begin{exoresolu}[Regroupement des 120 notes en classes d'amplitude 4]
\textbf{Énoncé.} Le dépouillement des $120$ notes de mathématiques du lycée technique de Nkongsamba donne les comptages suivants : $9$ notes dans $[0\,;4[$, $21$ notes dans $[4\,;8[$, $36$ notes dans $[8\,;12[$, $33$ notes dans $[12\,;16[$ et $21$ notes dans $[16\,;20]$. Construire le tableau complet (classes, centres, effectifs, fréquences en fraction et en pourcentage) et vérifier les totaux.
\tcblower
\textbf{Solution détaillée.}

\textbf{Étape 1 et 2.} L'amplitude est imposée : $4$ points. Les classes sont $[0\,;4[$, $[4\,;8[$, $[8\,;12[$, $[12\,;16[$, $[16\,;20]$.

\textbf{Étape 3.} Les centres :
\begin{align*}
c_1 &= \frac{0+4}{2} = 2, &
c_2 &= \frac{4+8}{2} = 6, \\
c_3 &= \frac{8+12}{2} = 10, &
c_4 &= \frac{12+16}{2} = 14, \\
c_5 &= \frac{16+20}{2} = 18.
\end{align*}

\textbf{Étape 4.} Les effectifs sont donnés par le dépouillement : $n_1 = 9$, $n_2 = 21$, $n_3 = 36$, $n_4 = 33$, $n_5 = 21$.

\textbf{Étape 5.} Vérification de l'effectif total :
\[ N = 9 + 21 + 36 + 33 + 21 = 120. \checkmark \]
Calcul des fréquences avec $f_i = \dfrac{n_i}{N} = \dfrac{n_i}{120}$ :
\begin{align*}
f_1 &= \frac{9}{120} = 0,075 = 7,5\,\%, &
f_2 &= \frac{21}{120} = 0,175 = 17,5\,\%, \\
f_3 &= \frac{36}{120} = 0,30 = 30\,\%, &
f_4 &= \frac{33}{120} = 0,275 = 27,5\,\%, \\
f_5 &= \frac{21}{120} = 0,175 = 17,5\,\%.
\end{align*}
Contrôle : $0,075 + 0,175 + 0,30 + 0,275 + 0,175 = 1$. $\checkmark$

\textbf{Tableau complet :}
\begin{center}
\begin{tabularx}{\linewidth}{|c|c|c|X|X|}
\hline
\rowcolor{popblueL} Classe & Centre $c_i$ & Effectif $n_i$ & Fréquence $f_i$ & Fréquence en \% \\
\hline
$[0\,;4[$ & $2$ & $9$ & $\dfrac{9}{120}=0,075$ & $7,5\,\%$ \\
\hline
$[4\,;8[$ & $6$ & $21$ & $\dfrac{21}{120}=0,175$ & $17,5\,\%$ \\
\hline
$[8\,;12[$ & $10$ & $36$ & $\dfrac{36}{120}=0,30$ & $30\,\%$ \\
\hline
$[12\,;16[$ & $14$ & $33$ & $\dfrac{33}{120}=0,275$ & $27,5\,\%$ \\
\hline
$[16\,;20]$ & $18$ & $21$ & $\dfrac{21}{120}=0,175$ & $17,5\,\%$ \\
\hline
\rowcolor{popcream} \textbf{Total} & -- & $N=120$ & $1$ & $100\,\%$ \\
\hline
\end{tabularx}
\end{center}

\textbf{Lecture.} La classe $[8\,;12[$ est la \textbf{classe modale} : $30\,\%$ des élèves. En additionnant les fréquences des classes $[8\,;12[$, $[12\,;16[$ et $[16\,;20]$, le proviseur lit immédiatement que $30 + 27,5 + 17,5 = 75\,\%$ des élèves ont au moins $8/20$, et que seuls $7,5\,\%$ sont en dessous de $4/20$. La décision pédagogique devient possible : un soutien ciblé pour le petit groupe en difficulté.
\end{exoresolu}

\begin{attention}[Erreurs fréquentes à éviter]
\begin{itemize}
\item \textbf{Oublier la convention $[a\,;b[$} et compter deux fois une valeur sur une borne : le total dépasse alors $N$.
\item \textbf{Confondre amplitude et centre} : l'amplitude est $b-a$, le centre est $\dfrac{a+b}{2}$.
\item \textbf{Calculer les fréquences avec un mauvais total} : toujours vérifier d'abord que $\sum n_i = N$ avant de diviser.
\item \textbf{Arrondir trop tôt} : garde les fractions $\dfrac{n_i}{N}$ ou au moins trois décimales, sinon la somme des pourcentages peut donner $99,9\,\%$ ou $100,1\,\%$.
\item \textbf{Oublier de fermer la dernière classe} : la borne maximale doit être incluse.
\end{itemize}
\end{attention}

\begin{exercice}[À toi de jouer]
Un atelier de mécanique de Douala mesure le diamètre (en mm) de $60$ arbres moteurs. Les valeurs s'étendent de $49{,}6$ mm à $50{,}4$ mm. Le chef d'atelier regroupe les mesures en $4$ classes d'amplitude $0{,}2$ : $[49,6\,;49,8[$, $[49,8\,;50[$, $[50\,;50,2[$, $[50,2\,;50,4]$, avec les effectifs $6$, $18$, $27$, $9$.
\begin{enumerate}
\item Vérifier l'effectif total $N$.
\item Calculer le centre de chaque classe.
\item Calculer la fréquence en pourcentage de chaque classe et vérifier que la somme vaut $100\,\%$.
\item Identifier la classe modale.
\end{enumerate}
\end{exercice}

\begin{corrige}[Exercice 1]
\begin{enumerate}
\item $N = 6 + 18 + 27 + 9 = 60$. $\checkmark$
\item Centres : $\dfrac{49,6+49,8}{2}=49,7$ ; $\dfrac{49,8+50}{2}=49,9$ ; $\dfrac{50+50,2}{2}=50,1$ ; $\dfrac{50,2+50,4}{2}=50,3$.
\item Fréquences : $\dfrac{6}{60}=10\,\%$ ; $\dfrac{18}{60}=30\,\%$ ; $\dfrac{27}{60}=45\,\%$ ; $\dfrac{9}{60}=15\,\%$. Somme : $10+30+45+15=100\,\%$. $\checkmark$
\item La classe modale est $[50\,;50,2[$ avec un effectif de $27$ (fréquence $45\,\%$).
\end{enumerate}
\end{corrige}

\begin{aretenir}[L'essentiel de la section]
\begin{itemize}
\item On regroupe en \cle{classes} quand les valeurs sont trop nombreuses ou continues.
\item Classe $[a\,;b[$ : amplitude $b-a$, centre $\dfrac{a+b}{2}$ ; borne de gauche incluse, borne de droite exclue (dernière classe fermée).
\item Effectif $n_i$, effectif total $N = \sum n_i$, fréquence $f_i = \dfrac{n_i}{N}$, et $\sum f_i = 1$.
\item La \cle{classe modale} est celle de plus grand effectif.
\item En pratique : $5$ à $10$ classes, amplitude constante, et toujours vérifier les totaux.
\end{itemize}
\end{aretenir}

\begin{savaistu}[Transition vers la suite]
Maintenant que le tableau est construit, la suite de la leçon montrera comment cumuler les effectifs et les fréquences (croissants et décroissants) pour répondre encore plus vite aux questions du type « combien d'élèves ont moins de $12/20$ ? », puis comment représenter graphiquement la série par un histogramme et des courbes cumulatives.
\end{savaistu}

\coursec{Effectifs et fréquences cumulés}

Dans la section précédente, nous avons appris à regrouper les données d'une série statistique en classes d'amplitude égale et à calculer les effectifs $n_i$ et les fréquences $f_i$ de chaque classe. Mais le proviseur qui étudie les notes de ses élèves ne se contente pas de savoir combien d'élèves ont une note entre 10 et 12 : il veut surtout répondre à des questions comme « combien d'élèves ont \emph{moins de} 12/20 ? » ou « combien ont \emph{au moins} 10/20 ? ». Ces questions font intervenir des \cle{accumulations} de plusieurs classes à la fois. C'est précisément le rôle des \cle{effectifs cumulés} et des \cle{fréquences cumulées}, que nous étudions maintenant.

\courssub{Effectif cumulé croissant (ECC)}

\textbf{Intuition.} Imaginons les élèves rangés en file indienne, de la plus petite note à la plus grande. L'effectif cumulé croissant d'une classe, c'est le nombre d'élèves que l'on compte en partant du début de la file jusqu'à la fin de cette classe. On \emph{accumule} les effectifs de bas en haut du tableau, comme une citerne qui se remplit classe après classe.

\begin{definition}[Effectif cumulé croissant]
Dans une série statistique regroupée en classes ordonnées, l'\cle{effectif cumulé croissant} (ECC) d'une classe est le nombre d'individus dont la valeur du caractère est \textbf{strictement inférieure à la borne supérieure} de cette classe.

On l'obtient en additionnant les effectifs de toutes les classes situées au-dessus (y compris la classe elle-même) :
\[
\mathrm{ECC}_i = n_1 + n_2 + \dots + n_i .
\]
\end{definition}

\begin{methode}[Remplir une colonne d'effectifs cumulés croissants]
\begin{enumerate}
\item Écrire les classes dans l'ordre croissant, de haut en bas.
\item Sur la première ligne, recopier l'effectif $n_1$ : c'est le premier $\mathrm{ECC}$.
\item Sur chaque ligne suivante, \textbf{additionner} l'effectif de la ligne à l'ECC de la ligne précédente : $\mathrm{ECC}_i = \mathrm{ECC}_{i-1} + n_i$.
\item \textbf{Vérifier} : le dernier $\mathrm{ECC}$ doit être égal à l'effectif total $N$. Si ce n'est pas le cas, une erreur d'addition s'est glissée.
\end{enumerate}
\end{methode}

\courssub{Effectif cumulé décroissant (ECD)}

\textbf{Intuition.} Reprenons notre file d'élèves rangés de la plus petite à la plus grande note. L'effectif cumulé décroissant d'une classe, c'est le nombre d'élèves que l'on compte en partant de la \emph{fin} de la file (les meilleures notes) jusqu'au \emph{début} de cette classe. On accumule cette fois de bas en haut du tableau.

\begin{definition}[Effectif cumulé décroissant]
Dans une série statistique regroupée en classes ordonnées, l'\cle{effectif cumulé décroissant} (ECD) d'une classe est le nombre d'individus dont la valeur du caractère est \textbf{supérieure ou égale à la borne inférieure} de cette classe.

On l'obtient en additionnant les effectifs de toutes les classes situées au-dessous (y compris la classe elle-même) :
\[
\mathrm{ECD}_i = n_i + n_{i+1} + \dots + n_k .
\]
\end{definition}

\begin{methode}[Remplir une colonne d'effectifs cumulés décroissants]
\begin{enumerate}
\item Partir de la \textbf{dernière ligne} du tableau : recopier l'effectif $n_k$ de la dernière classe.
\item \textbf{Remonter} ligne par ligne en additionnant : $\mathrm{ECD}_i = \mathrm{ECD}_{i+1} + n_i$.
\item \textbf{Vérifier} : le premier $\mathrm{ECD}$ (ligne du haut) doit être égal à l'effectif total $N$.
\end{enumerate}
\end{methode}

\begin{propriete}[Valeurs remarquables des cumuls]
Pour une série d'effectif total $N$ :
\begin{itemize}
\item le \textbf{dernier} effectif cumulé croissant vaut $N$ (tous les individus ont une valeur inférieure à la dernière borne supérieure) ;
\item le \textbf{premier} effectif cumulé décroissant vaut $N$ (tous les individus ont une valeur supérieure ou égale à la première borne inférieure).
\end{itemize}
De plus, pour chaque ligne $i$ : $\mathrm{ECC}_i + \mathrm{ECD}_{i+1} = N$ (résultat admis, que nous vérifierons sur l'exemple).
\end{propriete}

\begin{attention}[« Moins de $b$ » ou « au moins $a$ » ?]
L'erreur la plus fréquente au Probatoire est de confondre les deux types de cumuls. Retiens bien :
\begin{itemize}
\item « \textbf{moins de} $b$ » ou « \textbf{inférieur à} $b$ » $\longrightarrow$ on lit un \textbf{ECC} à la classe de borne supérieure $b$ ;
\item « \textbf{au moins} $a$ » ou « \textbf{supérieur ou égal à} $a$ » $\longrightarrow$ on lit un \textbf{ECD} à la classe de borne inférieure $a$.
\end{itemize}
Réfère-toi toujours aux \cle{bornes} des classes, jamais aux centres. Dire « moins de 12 » et lire l'ECD, c'est donner le nombre d'élèves ayant \emph{au moins} 12 : c'est exactement l'inverse de la réponse attendue !
\end{attention}

\courssub{Fréquences cumulées}

Les fréquences cumulées traduisent les mêmes informations que les effectifs cumulés, mais en \cle{proportion} (fraction) ou en \cle{pourcentage} de l'effectif total. Elles sont très pratiques pour comparer deux groupes de tailles différentes (par exemple deux ateliers d'effectifs différents).

\begin{definition}[Fréquences cumulées]
La \cle{fréquence cumulée croissante} (FCC) d'une classe est le quotient de son ECC par l'effectif total :
\[
\mathrm{FCC}_i = \dfrac{\mathrm{ECC}_i}{N}.
\]
La \cle{fréquence cumulée décroissante} (FCD) d'une classe est le quotient de son ECD par l'effectif total :
\[
\mathrm{FCD}_i = \dfrac{\mathrm{ECD}_i}{N}.
\]
On les exprime en fraction décimale (entre 0 et 1) ou en pourcentage (entre 0 \% et 100 \%). On peut aussi les obtenir en additionnant les fréquences $f_i$ de la même façon que les effectifs.
\end{definition}

\begin{aretenir}[À retenir absolument]
\begin{itemize}
\item ECC : on cumule de \textbf{haut en bas} ; le dernier vaut $N$ ; il répond à « moins de (borne supérieure) ».
\item ECD : on cumule de \textbf{bas en haut} ; le premier vaut $N$ ; il répond à « au moins (borne inférieure) ».
\item La dernière FCC vaut 1 (soit 100 \%) ; la première FCD vaut 1 (soit 100 \%).
\item $\mathrm{ECC}_i + \mathrm{ECD}_{i+1} = N$ et $\mathrm{FCC}_i + \mathrm{FCD}_{i+1} = 1$.
\end{itemize}
\end{aretenir}

\courssub{Exemple complet : les notes du proviseur}

Le proviseur d'un lycée technique de Douala a relevé les notes de mathématiques des 50 élèves d'une classe de Première technique. Les notes ont été regroupées en classes d'amplitude 2 (section précédente). Nous allons compléter \textbf{pas à pas} le tableau à sept colonnes.

\textbf{Étape 1 : les classes et les effectifs.} Les effectifs $n_i$ sont connus après le dépouillement.

\textbf{Étape 2 : colonne des ECC (de haut en bas).}
\begin{align*}
\mathrm{ECC}_1 &= 3 \\
\mathrm{ECC}_2 &= 3 + 7 = 10 \\
\mathrm{ECC}_3 &= 10 + 12 = 22 \\
\mathrm{ECC}_4 &= 22 + 15 = 37 \\
\mathrm{ECC}_5 &= 37 + 9 = 46 \\
\mathrm{ECC}_6 &= 46 + 4 = 50 \quad \checkmark \text{ (on retrouve bien } N = 50\text{)}
\end{align*}

\textbf{Étape 3 : colonne des ECD (de bas en haut).}
\begin{align*}
\mathrm{ECD}_6 &= 4 \\
\mathrm{ECD}_5 &= 4 + 9 = 13 \\
\mathrm{ECD}_4 &= 13 + 15 = 28 \\
\mathrm{ECD}_3 &= 28 + 12 = 40 \\
\mathrm{ECD}_2 &= 40 + 7 = 47 \\
\mathrm{ECD}_1 &= 47 + 3 = 50 \quad \checkmark \text{ (on retrouve bien } N = 50\text{)}
\end{align*}

\textbf{Étape 4 : fréquences et fréquences cumulées.} On divise chaque effectif (cumulé ou non) par $N = 50$, puis on convertit en pourcentage. Par exemple $f_3 = \dfrac{12}{50} = 0,24 = 24\,\%$ et $\mathrm{FCC}_3 = \dfrac{22}{50} = 0,44 = 44\,\%$.

\begin{center}
\small
\begin{tabularx}{\linewidth}{|c|X|X|X|X|X|X|}
\hline
\rowcolor{popblueL} \textbf{Classe} & \textbf{$n_i$} & \textbf{ECC} & \textbf{ECD} & \textbf{$f_i$ (\%)} & \textbf{FCC (\%)} & \textbf{FCD (\%)} \\
\hline
$[0\,;2[$ & 3 & 3 & 50 & 6 & 6 & 100 \\
\hline
$[2\,;4[$ & 7 & 10 & 47 & 14 & 20 & 94 \\
\hline
$[4\,;6[$ & 12 & 22 & 40 & 24 & 44 & 80 \\
\hline
$[6\,;8[$ & 15 & 37 & 28 & 30 & 74 & 56 \\
\hline
$[8\,;10[$ & 9 & 46 & 13 & 18 & 92 & 26 \\
\hline
$[10\,;12[$ & 4 & 50 & 4 & 8 & 100 & 8 \\
\hline
\rowcolor{popgoldL} \textbf{Total} & \textbf{50} & & & \textbf{100} & & \\
\hline
\end{tabularx}
\end{center}

\textit{Remarque : pour la lisibilité de l'exemple, les notes ont été regroupées sur l'intervalle $[0\,;12[$. La démarche serait identique avec des classes couvrant $[0\,;20]$.}

\textbf{Vérification de la propriété} $\mathrm{ECC}_i + \mathrm{ECD}_{i+1} = N$ : par exemple pour $i = 3$, $\mathrm{ECC}_3 + \mathrm{ECD}_4 = 22 + 28 = 50 = N$. La propriété est bien vérifiée sur chaque ligne.

\begin{popfigure}
\begin{tikzpicture}[scale=0.9]
% Escalier des ECC
\draw[->,thick,popdark] (0,0) -- (7.2,0) node[below left,font=\small]{classes};
\draw[->,thick,popdark] (0,0) -- (0,5.6) node[above,font=\small]{ECC};
\foreach \y/\lab in {0.3/3, 1/10, 2.2/22, 3.7/37, 4.6/46, 5/50}{
  \draw[popmuted,dashed] (0,\y) -- (7,\y);
  \node[left,font=\small,popdark] at (0,\y) {\lab};
}
\draw[very thick,popblue] (0,0.3) -- (1,0.3) -- (1,1) -- (2,1) -- (2,2.2) -- (3,2.2) -- (3,3.7) -- (4,3.7) -- (4,4.6) -- (5,4.6) -- (5,5) -- (6,5);
\fill[popblue!25] (0,0) rectangle (1,0.3);
\fill[popblue!18] (1,0) rectangle (2,1);
\fill[popblue!12] (2,0) rectangle (3,2.2);
\fill[popblue!08] (3,0) rectangle (4,3.7);
\fill[popblue!05] (4,0) rectangle (5,4.6);
\fill[popblue!03] (5,0) rectangle (6,5);
\draw[very thick,popblue] (0,0.3) -- (1,0.3) -- (1,1) -- (2,1) -- (2,2.2) -- (3,2.2) -- (3,3.7) -- (4,3.7) -- (4,4.6) -- (5,4.6) -- (5,5) -- (6,5);
\foreach \x/\lab in {0.5/{[0;2[}, 1.5/{[2;4[}, 2.5/{[4;6[}, 3.5/{[6;8[}, 4.5/{[8;10[}, 5.5/{[10;12[}}{
  \node[below,font=\scriptsize,popdark] at (\x,0) {\lab};
}
\foreach \p in {(1,1),(2,2.2),(3,3.7),(4,4.6),(5,5),(6,5)}{
  \fill[poporange] \p circle (2.2pt);
}
\node[popblue,font=\small\bfseries] at (3.4,5.3) {Accumulation des effectifs};
\end{tikzpicture}
\legende{Principe d'accumulation : chaque marche de l'escalier ajoute l'effectif $n_i$ d'une classe à la somme des précédentes. La dernière marche atteint $N = 50$.}
\end{popfigure}

\courssub{Interprétation concrète des cumuls}

Tout l'intérêt des colonnes cumulées est de répondre \emph{directement}, sans aucun calcul supplémentaire, aux questions que se posent un chef d'établissement, un chef d'atelier ou un contremaître.

\begin{exoresolu}[Questions du proviseur]
Énoncé. À l'aide du tableau précédent, répondre aux questions suivantes :
\begin{enumerate}
\item Combien d'élèves ont obtenu une note \textbf{inférieure à 8} ?
\item Combien d'élèves ont obtenu \textbf{au moins 6} ?
\item Quel \textbf{pourcentage} d'élèves a une note inférieure à 10 ?
\end{enumerate}
\tcblower
Solution détaillée.
\begin{enumerate}
\item « Inférieure à 8 » fait intervenir la borne supérieure 8 : on lit l'\textbf{ECC} de la classe $[6\,;8[$, soit $\mathrm{ECC}_4 = 37$. Donc \textbf{37 élèves} ont une note inférieure à 8.
\item « Au moins 6 » fait intervenir la borne inférieure 6 : on lit l'\textbf{ECD} de la classe $[6\,;8[$, soit $\mathrm{ECD}_4 = 28$. Donc \textbf{28 élèves} ont au moins 6.
\item « Inférieure à 10 » : on lit la \textbf{FCC} de la classe $[8\,;10[$, soit $\mathrm{FCC}_5 = 92\,\%$. Donc \textbf{92 \%} des élèves ont une note inférieure à 10.
\end{enumerate}
Vérification de cohérence pour la question 2 : $28 + 22 = 50 = N$, car les 22 élèves de $\mathrm{ECC}_3$ (note $< 6$) et les 28 élèves de $\mathrm{ECD}_4$ (note $\geq 6$) recouvrent toute la classe. \checkmark
\end{exoresolu}

\begin{attention}[Pièges de rédaction au Probatoire]
\begin{itemize}
\item Ne jamais additionner les fréquences cumulées entre elles : une FCC se calcule à partir des ECC (ou en cumulant les $f_i$), mais la somme de la colonne FCC n'a \textbf{aucune signification}.
\item Ne pas confondre fréquence et fréquence cumulée : $f_4 = 30\,\%$ signifie « 30 \% des élèves ont une note dans $[6\,;8[$ », tandis que $\mathrm{FCC}_4 = 74\,\%$ signifie « 74 \% ont une note inférieure à 8 ».
\item Toujours conclure par une phrase en français avec l'unité ou le pourcentage : « 28 élèves », « 92 \% des élèves ».
\end{itemize}
\end{attention}

\begin{savaistu}[Les cumuls dans la vie professionnelle]
Au Cameroun, les entreprises de distribution d'eau et d'électricité (ENEO, CAMWATER) utilisent des effectifs cumulés pour établir leurs \emph{tranches de facturation} : « combien d'abonnés consomment moins de 220 kWh par mois ? ». Les services du recensement (BUCREP) publient les pyramides des âges à partir de données cumulées. Maîtriser ECC et ECD, c'est déjà lire les tableaux de bord de la vie économique nationale !
\end{savaistu}

\begin{exercice}[À toi de jouer]
Un atelier de menuiserie de Yaoundé mesure les longueurs (en cm) de 40 planches découpées. Les classes et effectifs sont : $[148\,;150[$ : 5 ; $[150\,;152[$ : 11 ; $[152\,;154[$ : 14 ; $[154\,;156[$ : 7 ; $[156\,;158[$ : 3.
\begin{enumerate}
\item Construire le tableau complet à sept colonnes (classe, $n_i$, ECC, ECD, $f_i$ en \%, FCC en \%, FCD en \%).
\item Combien de planches mesurent moins de 154 cm ?
\item Quel pourcentage de planches mesure au moins 152 cm ?
\end{enumerate}
\end{exercice}

\begin{corrige}[Exercice 1]
1. ECC de haut en bas : 5 ; $5+11=16$ ; $16+14=30$ ; $30+7=37$ ; $37+3=40$ ($=N$ \checkmark). ECD de bas en haut : 3 ; $3+7=10$ ; $10+14=24$ ; $24+11=35$ ; $35+5=40$ ($=N$ \checkmark). Fréquences : $f_i = n_i/40$, soit 12,5 \% ; 27,5 \% ; 35 \% ; 17,5 \% ; 7,5 \%. FCC : 12,5 \% ; 40 \% ; 75 \% ; 92,5 \% ; 100 \%. FCD : 100 \% ; 87,5 \% ; 60 \% ; 25 \% ; 7,5 \%.

2. « Moins de 154 » : ECC de la classe $[152\,;154[$, soit \textbf{30 planches}.

3. « Au moins 152 » : FCD de la classe $[152\,;154[$, soit \textbf{60 \%} des planches.
\end{corrige}

En résumé, les effectifs et fréquences cumulés transforment un tableau brut en un véritable outil de décision : une simple lecture de ligne suffit pour répondre à « combien en dessous de… ? » ou « combien au-dessus de… ? ». Dans la section suivante, nous donnerons une image à ces nombres en construisant l'\emph{histogramme} d'une série groupée en classes.

\coursec{Histogramme d'une série groupée en classes}

Après avoir organisé les données en classes et calculé les effectifs et fréquences cumulés (croissants et décroissants), nous disposons d'un tableau de synthèse clair. Mais l'œil humain saisit souvent mieux une image qu'un tableau de nombres. C'est tout l'objet de l'histogramme : \textbf{représenter graphiquement la répartition des effectifs (ou des fréquences) par classe}. Nous verrons comment le construire rigoureusement, comment le lire pour en extraire la classe modale, et comment le compléter par le polygone des effectifs.

\courssub{Principes de construction}

Rappelons qu'une série groupée en classes partitionne l'intervalle des valeurs observées en \textbf{classes contiguës} (l'extrémité supérieure d'une classe est l'extrémité inférieure de la suivante). Chaque classe $C_i = [a_i ; a_{i+1}[$ (ou $]a_i ; a_{i+1}]$ pour la dernière) possède une \textbf{amplitude} $h_i = a_{i+1} - a_i$ et un \textbf{effectif} $n_i$ (ou une fréquence $f_i = n_i / N$).

\begin{definition}[Histogramme d'une série groupée en classes]
Soit une variable quantitative continue (ou discrète à nombreuses modalités) groupée en $k$ classes $C_1, \dots, C_k$ d'amplitudes $h_1, \dots, h_k$ et d'effectifs $n_1, \dots, n_k$ (ou fréquences $f_1, \dots, f_k$).
L'\textbf{histogramme} est le graphique constitué de $k$ rectangles juxtaposés (sans espace entre eux) tels que :
\begin{itemize}
    \item La base du $i$-ième rectangle s'étend sur l'intervalle de la classe $C_i$ (largeur $h_i$ sur l'axe des abscisses).
    \item La hauteur du $i$-ième rectangle est choisie de sorte que \textbf{l'aire du rectangle soit proportionnelle à l'effectif $n_i$ (ou à la fréquence $f_i$)}.
\end{itemize}
\end{definition}

\courssub{Cas des classes de même amplitude (cas courant au Probatoire)}

Lorsque toutes les classes ont la \textbf{même amplitude} $h$ (ex. : classes de largeur 4, 5, 10...), la condition « aire proportionnelle à l'effectif » se simplifie considérablement :
\[
\text{Aire}_i = h \times \text{Hauteur}_i \propto n_i \quad \Longrightarrow \quad \text{Hauteur}_i \propto n_i
\]
On peut donc prendre directement \textbf{l'effectif $n_i$ (ou la fréquence $f_i$) comme hauteur du rectangle} (à une unité d'échelle près). C'est ce que nous ferons dans la quasi-totalité des exercices de Première Technique.

\begin{attention}[Piège classique : les espaces entre les rectangles]
\emph{Un histogramme n'est PAS un diagramme en bâtons.}
\begin{itemize}
    \item \textbf{Diagramme en bâtons} (variable discrète ou qualitative) : les bâtons sont \textbf{séparés} par des espaces, car les modalités sont distinctes et non ordonnées de manière continue.
    \item \textbf{Histogramme} (variable continue groupée en classes) : les rectangles sont \textbf{juxtaposés (collés les uns aux autres)}, car les classes sont contiguës et couvrent tout l'intervalle des valeurs sans trou.
\end{itemize}
Laisser des espaces entre les rectangles d'un histogramme est une \textbf{erreur grave de rédaction} qui sera pénalisée au Probatoire.
\end{attention}

\courssub{Méthode : Construire un histogramme en 4 étapes}

\begin{methode}[Construction d'un histogramme (classes d'amplitude égale)]
\begin{enumerate}
    \item \textbf{Axes et graduation (abscisses) :} Sur l'axe horizontal, reporter les \textbf{bornes des classes} (ex. : 0, 4, 8, 12, 16, 20). Choisir une échelle adaptée à l'amplitude des classes. L'origine (0) n'est pas obligatoire si la première classe ne commence pas à 0 ; on peut utiliser une \textbf{brise d'axe} (zigzag) pour signaler la coupure.
    \item \textbf{Graduation (ordonnées) :} Sur l'axe vertical, choisir une échelle pour les effectifs (ou fréquences). L'unité doit permettre de tracer le rectangle le plus haut sans dépasser la feuille. Graduation régulière (ex. : 1 cm pour 10 élèves).
    \item \textbf{Tracé des rectangles :} Pour chaque classe $[a_i ; a_{i+1}[$ d'effectif $n_i$, tracer un rectangle :
    \begin{itemize}
        \item Base : segment $[a_i ; a_{i+1}]$ sur l'axe des abscisses.
        \item Hauteur : proportionnelle à $n_i$ (selon l'échelle choisie).
        \item Les rectangles se touchent par leurs côtés verticaux (pas d'espace).
    \end{itemize}
    \item \textbf{Finalisation :} Donner un \textbf{titre} précis (ex. : « Histogramme des notes sur 20 de la classe de 1re T3 »), légender les axes (variable étudiée + unité ; effectif ou fréquence), indiquer l'échelle.
\end{enumerate}
\end{methode}

\courssub{Illustration : Histogramme des notes (classes d'amplitude 4)}

Prenons l'exemple classique des 120 notes sur 20, regroupées en classes d'amplitude 4.

\begin{center}
\begin{tabular}{|c|c|c|c|}\hline
\textbf{Classe (notes)} & \textbf{Effectif $n_i$} & \textbf{Fréquence $f_i$} & \textbf{Milieu $x_i$} \\ \hline
$[0 ; 4[$ & 5 & 0,042 & 2 \\ \hline
$[4 ; 8[$ & 18 & 0,150 & 6 \\ \hline
$[8 ; 12[$ & \textbf{42} & \textbf{0,350} & \textbf{10} \\ \hline
$[12 ; 16[$ & 35 & 0,292 & 14 \\ \hline
$[16 ; 20]$ & 20 & 0,167 & 18 \\ \hline
\textbf{Total} & \textbf{120} & \textbf{1} &  \\ \hline
\end{tabular}
\end{center}

Voici l'histogramme correspondant, tracé avec \texttt{pgfplots} (option \texttt{ybar interval} qui gère automatiquement les bases sur les bornes de classes). Le paramètre \texttt{ybar interval=1} (valeur par défaut) garantit que les rectangles sont juxtaposés sans espace.

\begin{popfigure}
\begin{tikzpicture}
\begin{axis}[
    width=\linewidth,
    height=8cm,
    ybar interval=1, % largeur = amplitude de la classe (rectangles juxtaposés)
    xtick={0,4,8,12,16,20},
    xticklabels={0,4,8,12,16,20},
    xlabel={Notes sur 20},
    ylabel={Effectif},
    ymin=0, ymax=50,
    ytick={0,10,20,30,40,50},
    title={Histogramme des notes de 120 élèves (classes d'amplitude 4)},
    title style={yshift=0.5cm, font=\bfseries},
    axis lines*=left,
    enlarge x limits=false, % important pour ybar interval : pas de marge avant 0 et après 20
    bar cycle list={popblue!70, poporange!70, popgreen!70, poppurple!70, poppink!70},
    grid=major,
    grid style={popline},
]
\addplot+ [draw=popdark, fill=popblue!60] coordinates {
    (0,5)   % [0;4[
    (4,18)  % [4;8[
    (8,42)  % [8;12[  <- Modale
    (12,35) % [12;16[
    (16,20) % [16;20]
    (20,0)  % point final fictif pour fermer la dernière barre dans ybar interval
};
\end{axis}
\end{tikzpicture}
\legende{Histogramme des 120 notes. La classe modale $[8;12[$ apparaît clairement comme le rectangle le plus haut.}
\end{popfigure}

\courssub{Lecture et interprétation : la classe modale}

L'histogramme permet une lecture visuelle immédiate de la répartition.

\begin{definition}[Classe modale d'une série groupée en classes]
La \textbf{classe modale} est la classe qui possède \textbf{l'effectif le plus grand} (ou la fréquence la plus grande). Sur l'histogramme, c'est la classe dont le rectangle a \textbf{la plus grande hauteur} (puisque les amplitudes sont égales).
\end{definition}

Dans notre exemple, l'effectif maximal est $42$ pour la classe $[8 ; 12[$. C'est donc la \textbf{classe modale}. On dit que les notes sont « concentrées autour de 10 ».

\begin{aretenir}
\begin{itemize}
    \item La classe modale donne une \textbf{estimation grossière du mode} (valeur la plus fréquente) pour une variable continue.
    \item Si l'histogramme a \textbf{deux pics} de hauteur comparable (deux classes modales distinctes), on parle de distribution \textbf{bimodale} (ex. : mélange de deux populations : garçons/filles, débutants/experts).
    \item La forme de l'histogramme (symétrique, étirée à droite « asymétrie positive », étirée à gauche « asymétrie négative ») renseigne sur la dispersion et la concentration des données.
\end{itemize}
\end{aretenir}

\courssub{Le polygone des effectifs (ou polygone des fréquences)}

Souvent superposé à l'histogramme, le polygone permet de visualiser la tendance générale de la répartition et de comparer plusieurs séries sur le même graphique.

\begin{definition}[Polygone des effectifs d'une série groupée en classes]
Le \textbf{polygone des effectifs} est la ligne brisée obtenue en reliant entre eux les \textbf{milieux des sommets supérieurs} des rectangles de l'histogramme.
Pour la première et la dernière classe, on prolonge la ligne brisée jusqu'à l'axe des abscisses en ajoutant une classe fictive d'effectif nul de part et d'autre (milieu $x_0$ avant la première classe, $x_{k+1}$ après la dernière).
\end{definition}

\begin{methode}[Tracer le polygone des effectifs]
\begin{enumerate}
    \item Calculer le \textbf{milieu} $x_i = \frac{a_i + a_{i+1}}{2}$ de chaque classe.
    \item Reporter les points $M_i(x_i ; n_i)$ (ou $M_i(x_i ; f_i)$) dans le repère de l'histogramme.
    \item Ajouter $M_0(x_0 ; 0)$ et $M_{k+1}(x_{k+1} ; 0)$ où $x_0 = x_1 - h$ et $x_{k+1} = x_k + h$ ($h$ amplitude commune).
    \item Relier les points $M_0, M_1, \dots, M_k, M_{k+1}$ par des segments de droite (ligne brisée).
\end{enumerate}
\end{methode}

Voici le même histogramme avec le polygone des effectifs superposé (en rouge).

\begin{popfigure}
\begin{tikzpicture}
\begin{axis}[
    width=\linewidth,
    height=8cm,
    ybar interval=1,
    xtick={0,4,8,12,16,20},
    xticklabels={0,4,8,12,16,20},
    xlabel={Notes sur 20},
    ylabel={Effectif},
    ymin=0, ymax=50,
    ytick={0,10,20,30,40,50},
    title={Histogramme et polygone des effectifs (120 notes)},
    title style={yshift=0.5cm, font=\bfseries},
    axis lines*=left,
    enlarge x limits=false,
    bar cycle list={popblue!50},
    grid=major,
    grid style={popline},
    legend pos=north east,
    legend cell align=left,
]
% Histogramme
\addplot+ [draw=popdark, fill=popblue!50, fill opacity=0.6] coordinates {
    (0,5) (4,18) (8,42) (12,35) (16,20) (20,0)
};
\addlegendentry{Histogramme (effectifs)}

% Polygone des effectifs : points (milieu, effectif) + extrémités nulles
% Milieux : 2, 6, 10, 14, 18. Extrémités : -2 (0), 22 (0).
\addplot+ [popred, thick, mark=*, mark size=2.5pt, mark options={fill=popred}] coordinates {
    (-2,0)  % classe fictive avant [0;4[
    (2,5)   % milieu [0;4[
    (6,18)  % milieu [4;8[
    (10,42) % milieu [8;12[  <- sommet modale
    (14,35) % milieu [12;16[
    (18,20) % milieu [16;20]
    (22,0)  % classe fictive après [16;20]
};
\addlegendentry{Polygone des effectifs}
\end{axis}
\end{tikzpicture}
\legende{Superposition de l'histogramme et du polygone des effectifs. Le polygone passe par les milieux des sommets des rectangles et rejoint l'axe des abscisses aux extrémités.}
\end{popfigure}

\courssub{Exemple résolu : Analyse de l'histogramme du Proviseur}

\begin{exoresolu}[Histogramme des notes du Proviseur]
\textbf{Énoncé :} Le proviseur d'un lycée technique a relevé les notes sur 20 obtenues par 120 élèves de Première à l'épreuve de mathématiques du premier trimestre. Les notes sont regroupées en classes d'amplitude 4 comme suit :

\begin{center}
\begin{tabular}{|c|c|c|}\hline
\textbf{Classe} & \textbf{Effectif} & \textbf{Fréquence} \\ \hline
$[0 ; 4[$ & 5 & 0,042 \\ \hline
$[4 ; 8[$ & 18 & 0,150 \\ \hline
$[8 ; 12[$ & 42 & 0,350 \\ \hline
$[12 ; 16[$ & 35 & 0,292 \\ \hline
$[16 ; 20]$ & 20 & 0,167 \\ \hline
\textbf{Total} & \textbf{120} & \textbf{1} \\ \hline
\end{tabular}
\end{center}

\textbf{Questions :}
\begin{enumerate}
    \item Construire l'histogramme de cette série statistique (échelle : 1 cm pour 4 unités en abscisses, 1 cm pour 5 élèves en ordonnées).
    \item Déterminer la classe modale et l'interpréter.
    \item Tracer le polygone des effectifs sur le même graphique.
    \item Quelle proportion d'élèves a obtenu une note \textbf{strictement inférieure à 12} ? \textbf{Supérieure ou égale à 12} ?
\end{enumerate}

\tcblower

\textbf{Solution détaillée :}

\noindent \textbf{1. Construction de l'histogramme}
\begin{itemize}
    \item \textbf{Axe des abscisses (notes) :} Échelle 1 cm $\rightarrow$ 4 unités. Les bornes 0, 4, 8, 12, 16, 20 sont espacées de 1 cm. On trace les 5 rectangles de base 1 cm juxtaposés.
    \item \textbf{Axe des ordonnées (effectifs) :} Échelle 1 cm $\rightarrow$ 5 élèves. Effectif max = 42 $\rightarrow$ hauteur max = 8.4cm (tenable sur une feuille A4).
    \item \textbf{Hauteurs des rectangles :}
    \begin{align*}
        [0;4[ &: 5 \to 1,0 \text{ cm} \\
        [4;8[ &: 18 \to 3,6 \text{ cm} \\
        [8;12[ &: 42 \to 8,4 \text{ cm} \quad \text{(le plus haut)} \\
        [12;16[ &: 35 \to 7,0 \text{ cm} \\
        [16;20] &: 20 \to 4,0 \text{ cm}
    \end{align*}
    \item \textbf{Titre :} « Histogramme des notes de Mathématiques (120 élèves de 1re Tech) ».
\end{itemize}

\noindent \textbf{2. Classe modale}
L'effectif maximal est $n_3 = 42$ pour la classe $[8 ; 12[$.
$\rightarrow$ \textbf{La classe modale est $[8 ; 12[$.}
\textit{Interprétation :} C'est dans cet intervalle que se trouve le plus grand nombre d'élèves (42 sur 120, soit 35 \%). La performance « typique » du groupe se situe autour de 10/20.

\noindent \textbf{3. Polygone des effectifs}
Milieux des classes : $x_1=2,\; x_2=6,\; x_3=10,\; x_4=14,\; x_5=18$.
Classes fictives d'effectif nul : $x_0 = -2$ (avant 0), $x_6 = 22$ (après 20).
Points à relier : $M_0(-2;0),\; M_1(2;5),\; M_2(6;18),\; M_3(10;42),\; M_4(14;35),\; M_5(18;20),\; M_6(22;0)$.
On trace la ligne brisée $M_0 M_1 M_2 M_3 M_4 M_5 M_6$ en pointillés rouges (par exemple) par-dessus les rectangles.

\noindent \textbf{4. Proportions (lecture sur effectifs cumulés croissants)}
\begin{itemize}
    \item Notes $< 12$ : classes $[0;4[$, $[4;8[$, $[8;12[$. Effectif cumulé = $5 + 18 + 42 = 65$.
    Proportion = $\frac{65}{120} \approx 0,542 = 54,2\%$.
    \item Notes $\ge 12$ : classes $[12;16[$, $[16;20]$. Effectif = $35 + 20 = 55$.
    Proportion = $\frac{55}{120} \approx 0,458 = 45,8\%$.
    \item \textit{Vérification :} $65 + 55 = 120$. $54,2\% + 45,8\% = 100\%$.
\end{itemize}
\end{exoresolu}

\courssub{Complément (horizon Terminale) : Classes d'amplitude inégale}

\begin{savaistu}[Pourquoi l'aire et pas la hauteur ?]
Au Probatoire, vous n'aurez quasiment que des classes de \textbf{même amplitude}. Mais en Terminale (et dans la vie professionnelle : histogrammes d'âges en démographie, de salaires en économie, de précipitations en météorologie), les classes ont souvent des amplitudes variables (ex. : $[0;20[$, $[20;40[$, $[40;50[$, $[50;60[$, $[60;80[$).

Si on gardait la règle « hauteur = effectif », une classe large (amplitude 20) avec 10 individus ferait un rectangle énorme visuellement comparé à une classe étroite (amplitude 5) avec 10 individus aussi. L'œil serait trompé : on croirait qu'il y a plus de monde dans la classe large.

\emph{La règle universelle} est donc : \textbf{L'aire du rectangle est proportionnelle à l'effectif.}
\[
\text{Hauteur}_i = \frac{\text{Effectif}_i}{\text{Amplitude}_i} \times \text{constante}
\]
Cette hauteur s'appelle la \textbf{densité d'effectif} (ou densité de fréquence). L'histogramme devient alors une estimation graphique de la \textbf{densité de probabilité} de la variable aléatoire continue sous-jacente.
\end{savaistu}

\begin{attention}[Résumé des points de vigilance pour le Probatoire]
\begin{itemize}
    \item \textbf{Juxtaposition :} Rectangles collés (pas d'espace). \cle{Zéro tolérance} sur ce point.
    \item \textbf{Échelles :} Toujours indiquées sur les axes (ex. : 1 cm $\leftrightarrow$ 5 élèves). L'origine des ordonnées doit être à 0 (ou brisure signalée).
    \item \textbf{Unités :} Axe horizontal = unité de la variable (notes, mm, kg, €). Axe vertical = effectif (sans unité) ou fréquence (en \%, ou nombre décimal).
    \item \textbf{Classe modale :} Se lit sur la \textbf{hauteur} (si amplitudes égales) ou sur l'\textbf{aire} (si amplitudes inégales). Ne pas confondre avec le milieu de la classe modale.
    \item \textbf{Polygone :} Passe par les \textbf{milieux} des sommets, rejoint l'axe des abscisses aux extrémités.
\end{itemize}
\end{attention}

\courssub{Exercice d'application immédiate}

\begin{exercice}[Durée de vie de piles (en heures)]
On a testé la durée de vie de 200 piles alcalines. Résultats groupés en classes :

\begin{center}
\begin{tabular}{|c|c|}\hline
\textbf{Durée (heures)} & \textbf{Effectif} \\ \hline
$[0 ; 50[$ & 10 \\ \hline
$[50 ; 100[$ & 25 \\ \hline
$[100 ; 150[$ & 55 \\ \hline
$[150 ; 200[$ & 60 \\ \hline
$[200 ; 300[$ & 50 \\ \hline
\textbf{Total} & \textbf{200} \\ \hline
\end{tabular}
\end{center}

\textbf{Remarquez :} La dernière classe a une amplitude de 100, les autres 50. (C'est un exemple « hors programme strict » pour vous préparer à la Terminale, mais ici on vous demande juste l'histogramme « naïf » hauteur = effectif, puis la correction).

\begin{enumerate}
    \item Tracer l'histogramme « naïf » (hauteur = effectif) sur papier millimétré. Quelle classe semble modale visuellement ?
    \item Calculer les \textbf{densités d'effectif} (effectif / amplitude) pour chaque classe.
    \item Tracer l'histogramme \textbf{correct} (hauteur = densité). Quelle est la \textbf{vraie} classe modale (celle de plus grande aire = plus grand effectif) ?
    \item Comparer les deux graphiques : que cache l'histogramme naïf ?
\end{enumerate}
\end{exercice}

\begin{corrige}[Exercice Durée de vie de piles]
\begin{enumerate}
    \item \textbf{Histogramme naïf :} Hauteurs : 10, 25, 55, 60, 50. Visuellement, la classe $[150;200[$ (hauteur 60) est la plus haute, mais la classe $[200;300[$ (hauteur 50) est deux fois plus large : son aire visuelle (5000 unités\texteuro{}aire) dépasse celle de la classe modale (3000 unités\texteuro{}aire). L'œil est trompé.
    \item \textbf{Densités $d_i = n_i / h_i$ :}
    \begin{align*}
        [0;50[ &: 10 / 50 = \mathbf{0,2} \\
        [50;100[ &: 25 / 50 = \mathbf{0,5} \\
        [100;150[ &: 55 / 50 = \mathbf{1,1} \\
        [150;200[ &: 60 / 50 = \mathbf{1,2} \quad \text{(densité max)} \\
        [200;300[ &: 50 / 100 = \mathbf{0,5}
    \end{align*}
    \item \textbf{Histogramme correct (aire = effectif) :} Hauteurs = densités ci-dessus.
    La classe modale (plus grand effectif = plus grande aire) reste $[150;200[$ ($n=60$).
    Mais visuellement, le rectangle $[200;300[$ est maintenant \textbf{large (100) et bas (0,5)}, son aire = $100 \times 0,5 = 50$ (effectif 50). Il ne domine plus le graphique.
    \item \textbf{Conclusion :} L'histogramme naïf \textbf{sur-représente visuellement les classes larges} et sous-représente les classes étroites. Seul l'histogramme à aires proportionnelles (densités) donne une image fidèle de la répartition de la variable continue.
\end{enumerate}
\end{corrige}

\courssub{Synthèse de la section}

\begin{aretenir}[Ce qu'il faut retenir pour l'examen]
\begin{itemize}
    \item L'\textbf{histogramme} représente une série groupée en classes par des \textbf{rectangles juxtaposés} (bases = bornes des classes).
    \item \textbf{Classes même amplitude} : Hauteur $\propto$ Effectif (ou Fréquence). C'est le cas standard au Probatoire.
    \item \textbf{Classe modale} = classe d'effectif maximal = rectangle le plus haut (si même amplitude).
    \item \textbf{Polygone des effectifs} : ligne brisée reliant les \textbf{milieux des sommets} des rectangles, prolongée jusqu'à l'axe des abscisses par des classes fictives d'effectif nul.
    \item \textbf{Erreur interdite :} Espaces entre les rectangles (confusion avec diagramme en bâtons).
    \item \textbf{Terminale / Pro :} Si amplitudes inégales $\rightarrow$ Hauteur = Densité = Effectif / Amplitude. Aire = Effectif.
\end{itemize}
\end{aretenir}

La section suivante (« Courbes cumulatives ») exploitera les effectifs cumulés calculés précédemment pour construire une représentation graphique permettant de lire directement les quartiles et la médiane.

\coursec{Courbes cumulatives}

Après avoir construit l'histogramme d'une série groupée en classes, nous disposons d'une vision globale de la répartition des valeurs. Mais l'histogramme ne permet pas de lire directement « combien de valeurs sont inférieures à un seuil donné » ou « quelle valeur dépasse 75 \% des observations ». Pour répondre à ces questions, nous utilisons les \textbf{courbes cumulatives} (aussi appelées \emph{ogives}). Elles transforment les effectifs cumulés en une représentation graphique continue, immédiatement exploitable.

\courssub{Courbe cumulative croissante}

\emph{Intuition.} Imaginez que vous notiez, pour chaque note possible, le nombre d'élèves ayant obtenu \textbf{au plus} cette note. En reliant ces points, vous obtenez une courbe qui monte progressivement de 0 jusqu'à l'effectif total $N$. C'est exactement la courbe cumulative croissante.

\begin{definition}[Courbe cumulative croissante]
Soit une série statistique regroupée en $k$ classes d'intervalles $[a_1; a_2[$, $[a_2; a_3[$, …, $[a_k; a_{k+1}[$, d'effectifs respectifs $n_1, n_2, \dots, n_k$ et d'effectifs cumulés croissants $ECC_1, ECC_2, \dots, ECC_k$ (avec $ECC_i = n_1 + \dots + n_i$). La \textbf{courbe cumulative croissante} est la courbe obtenue en traçant les points de coordonnées $(a_{i+1}\,;\, ECC_i)$ pour $i = 1 \dots k$, en ajoutant le point de départ $(a_1\,;\, 0)$, et en reliant ces points par des segments de droite.
\end{definition}

\begin{methode}[Tracer une courbe cumulative croissante]
\begin{enumerate}
\item Construire le tableau des effectifs cumulés croissants ($ECC$).
\item Choisir une échelle adaptée sur l'axe des abscisses (bornes des classes) et sur l'axe des ordonnées (effectifs de 0 à $N$).
\item Placer le point de départ $(\text{borne inférieure de la 1re classe}\,;\, 0)$.
\item Pour chaque classe, placer le point $(\text{borne supérieure de la classe}\,;\, ECC \text{ de cette classe})$.
\item Relier tous les points par des segments de droite, dans l'ordre croissant des abscisses.
\end{enumerate}
\end{methode}

\begin{attention}[Erreur fréquente : abscisses des points]
Ne placez \textbf{jamais} les effectifs cumulés croissants en face des \emph{centres} des classes (valeurs moyennes). Les $ECC$ correspondent au nombre de valeurs \textbf{inférieures ou égales à la borne supérieure} de la classe. Les points se placent donc aux \textbf{bornes supérieures} des classes. Le point de départ se place à la borne inférieure de la première classe, d'ordonnée 0.
\end{attention}

\begin{popfigure}
\begin{tikzpicture}
\begin{axis}[
    width=\linewidth,
    height=0.55\linewidth,
    xlabel={Notes},
    ylabel={Effectifs cumulés croissants},
    xmin=0, xmax=20,
    ymin=0, ymax=32,
    xtick={0,2,4,6,8,10,12,14,16,18,20},
    ytick={0,5,10,15,20,25,30},
    grid=major,
    grid style={popline},
    axis lines=middle,
    enlargelimits={abs=0.5},
    clip=false,
]
% Data: classes [0;2[ [2;4[ [4;6[ [6;8[ [8;10[ [10;12[ [12;14[ [14;16[ [16;18[ [18;20]
% Effectifs: 1, 2, 3, 5, 7, 6, 4, 2, 1, 0  (total 31)
% ECC: 1, 3, 6, 11, 18, 24, 28, 30, 31, 31
% Points: (0,0), (2,1), (4,3), (6,6), (8,11), (10,18), (12,24), (14,28), (16,30), (18,31), (20,31)
\addplot[popcurve, very thick, mark=*, mark size=2.5, mark options={fill=popblue, draw=popdark}] coordinates {
    (0,0) (2,1) (4,3) (6,6) (8,11) (10,18) (12,24) (14,28) (16,30) (18,31) (20,31)
};
% Vertical dashed line at x=10 for reading
\addplot[dashed, poporange, thick] coordinates {(10,0) (10,18)};
\addplot[dashed, poporange, thick] coordinates {(0,18) (10,18)};
\node[poporange, anchor=south west] at (10,18) {\small $(10;\,18)$};
\node[poporange, anchor=north] at (10,0) {\small $10$};
\node[poporange, anchor=east] at (0,18) {\small $18$};
\end{axis}
\end{tikzpicture}
\legende{Courbe cumulative croissante des notes d'une classe de 31 élèves. Lecture graphique : 18 élèves ont une note $\le 10$.}
\end{popfigure}

\courssub{Courbe cumulative décroissante}

\emph{Intuition symétrique.} Au lieu de compter « au plus », comptons « au moins ». Pour chaque note, notons le nombre d'élèves ayant obtenu \textbf{au moins} cette note. La courbe part de $N$ à la borne inférieure de la première classe et descend jusqu'à 0 à la borne supérieure de la dernière classe.

\begin{definition}[Courbe cumulative décroissante]
Avec les mêmes notations, les effectifs cumulés décroissants sont $ECD_i = n_i + n_{i+1} + \dots + n_k$ (avec $ECD_1 = N$). La \textbf{courbe cumulative décroissante} est la courbe obtenue en traçant les points de coordonnées $(a_i\,;\, ECD_i)$ pour $i = 1 \dots k$, en ajoutant le point final $(a_{k+1}\,;\, 0)$, et en reliant ces points par des segments.
\end{definition}

\begin{methode}[Tracer une courbe cumulative décroissante]
\begin{enumerate}
\item Construire le tableau des effectifs cumulés décroissants ($ECD$). On a $ECD_1 = N$, puis $ECD_{i+1} = ECD_i - n_i$.
\item Placer le point de départ $(\text{borne inférieure de la 1re classe}\,;\, N)$.
\item Pour chaque classe, placer le point $(\text{borne inférieure de la classe}\,;\, ECD \text{ de cette classe})$.
\item Ajouter le point final $(\text{borne supérieure de la dernière classe}\,;\, 0)$.
\item Relier tous les points par des segments.
\end{enumerate}
\end{methode}

\begin{attention}[Bornes inférieures pour la décroissante]
Pour la courbe décroissante, les points se placent aux \textbf{bornes inférieures} des classes (sauf le dernier point à la borne supérieure finale). C'est l'erreur symétrique de celle de la courbe croissante.
\end{attention}

\courssub{Les deux courbes sur un même repère}

Tracer les deux courbes sur le même graphique permet de visualiser leur position relative et d'estimer la médiane.

\begin{propriete}[Intersection des deux courbes cumulatives (admise)]
Les courbes cumulatives croissante et décroissante se coupent en un unique point. L'ordonnée de ce point est voisine de $\dfrac{N}{2}$. Son abscisse donne une \textbf{estimation graphique de la médiane} de la série statistique. Cette propriété est admise en Première technique ; la médiane sera étudiée rigoureusement en Leçon 16.
\end{propriete}

\begin{popfigure}
\begin{tikzpicture}
\begin{axis}[
    width=\linewidth,
    height=0.6\linewidth,
    xlabel={Notes},
    ylabel={Effectifs cumulés},
    xmin=0, xmax=20,
    ymin=0, ymax=33,
    xtick={0,2,4,6,8,10,12,14,16,18,20},
    ytick={0,5,10,15,20,25,30},
    grid=major,
    grid style={popline},
    axis lines=middle,
    enlargelimits={abs=0.5},
    clip=false,
    legend style={at={(0.98,0.02)}, anchor=south east, fill=white, draw=popline},
]
% Croissante
\addplot[popblue, very thick, mark=*, mark size=2, mark options={fill=popblue, draw=popdark}] coordinates {
    (0,0) (2,1) (4,3) (6,6) (8,11) (10,18) (12,24) (14,28) (16,30) (18,31) (20,31)
};
\addlegendentry{Croissante (ECC)}
% Décroissante
\addplot[poporange, very thick, mark=square*, mark size=2, mark options={fill=poporange, draw=popdark}] coordinates {
    (0,31) (2,30) (4,28) (6,25) (8,20) (10,13) (12,7) (14,3) (16,1) (18,0) (20,0)
};
\addlegendentry{Décroissante (ECD)}
% Intersection approx at x=9.5, y=15.5
\addplot[only marks, mark=*, mark size=4, popgreen] coordinates {(9.5,15.5)};
\node[popgreen, anchor=south west] at (9.5,15.5) {\textbf{I}};
% Median vertical line
\addplot[dashed, popgreen, thick] coordinates {(9.5,0) (9.5,15.5)};
\node[popgreen, anchor=north] at (9.5,0) {\small $\tilde{x} \approx 9,5$};
% N/2 horizontal line
\addplot[dashed, popmuted, thin] coordinates {(0,15.5) (20,15.5)};
\node[popmuted, anchor=east] at (0,15.5) {\small $N/2 = 15,5$};
\end{axis}
\end{tikzpicture}
\legende{Courbes cumulatives croissante (bleue) et décroissante (orange) sur le même repère. Leur intersection I donne une estimation graphique de la médiane $\tilde{x} \approx 9,5$.}
\end{popfigure}

\courssub{Lecture graphique et interpolation linéaire}

Les courbes cumulatives permettent deux types de lecture :

\begin{itemize}
\item \textbf{Abscisse $\to$ Ordonnée :} Donnée une valeur $x$ (borne de classe ou valeur intermédiaire), on lit l'effectif cumulé correspondant. Si $x$ n'est pas une borne de classe, on utilise l'\textbf{interpolation linéaire} sur le segment concerné : on suppose que les valeurs sont uniformément réparties à l'intérieur de la classe.
\item \textbf{Ordonnée $\to$ Abscisse :} Donné un effectif cumulé $E$ (ou une fréquence cumulée), on lit la valeur correspondante. C'est l'opération inverse, elle aussi par interpolation linéaire.
\end{itemize}

\begin{exemplebox}[Lecture graphique sur la courbe cumulative croissante]
Reprenons la série des notes (31 élèves). On veut connaître le nombre d'élèves ayant obtenu \textbf{strictement moins de 10,4}.
\begin{itemize}
\item Sur l'axe des abscisses, repérer $10,4$ (dans la classe $[10;12[$).
\item Monter verticalement jusqu'à la courbe croissante.
\item Lire l'ordonnée par interpolation : le segment va de $(10;18)$ à $(12;24)$. La pente est $\frac{24-18}{12-10} = 3$ élèves par point de note.
\item Pour $x = 10,4$, l'ordonnée est $18 + 3 \times (10,4 - 10) = 18 + 1,2 = 19,2$.
\item \textbf{Interprétation :} Environ 19 élèves ont une note $< 10,4$ (on arrondit à l'entier le plus proche, soit 19).
\end{itemize}
\end{exemplebox}

\begin{exoresolu}[Lecture inverse : trouver le seuil pour un pourcentage donné]
\textbf{Énoncé.} Dans la même classe de 31 élèves, on souhaite déterminer la note en dessous de laquelle se situent les $25\%$ des élèves les plus faibles (premier quartile $Q_1$).
\textbf{Solution.}
\begin{enumerate}
\item $25\%$ de $31$ élèves correspondent à $0,25 \times 31 = 7,75$ élèves. On cherche l'abscisse pour laquelle $ECC \approx 7,75$ (ou $8$).
\item Sur la courbe croissante, repérer l'ordonnée $7,75$ (entre $6$ et $11$).
\item Le segment concerné relie $(6;6)$ à $(8;11)$. Pente $= \frac{11-6}{8-6} = 2,5$ élèves par point.
\item Partant de $(6;6)$, il faut monter de $7,75 - 6 = 1,75$ élèves. L'abscisse correspondante est $6 + \frac{1,75}{2,5} = 6 + 0,7 = 6,7$.
\item \textbf{Conclusion :} Environ $25\%$ des élèves ont une note inférieure à $6,7/20$. C'est une estimation graphique du premier quartile.
\end{enumerate}
\tcblower
\textbf{Solution détaillée :} L'interpolation linéaire suppose une répartition uniforme dans la classe $[6;8[$ (effectif 5). Le calcul exact par la formule du quartile (Leçon 16) donnera une valeur très proche.
\end{exoresolu}

\courssub{Version en fréquences cumulées (pourcentages)}

Il est souvent plus commode de travailler en \textbf{fréquences cumulées} (en $\%$) plutôt qu'en effectifs. Les courbes ont alors la même forme, mais l'axe des ordonnées va de $0\%$ à $100\%$.

\begin{definition}[Courbes de fréquences cumulées]
On divise chaque effectif cumulé par l'effectif total $N$ et on multiplie par $100$. La courbe cumulative croissante en fréquences passe par $(\text{borne inf. 1re classe}; 0\%)$ et $(\text{borne sup. dernière classe}; 100\%)$. La courbe décroissante passe par $(\text{borne inf. 1re classe}; 100\%)$ et $(\text{borne sup. dernière classe}; 0\%)$.
\end{definition}

\begin{aretenir}[Avantages de la version en \%]
\begin{itemize}
\item Lecture directe des seuils usuels : $25\%$ (premier quartile), $50\%$ (médiane), $75\%$ (troisième quartile), $90\%$, $95\%$, etc.
\item Comparaison immédiate entre deux populations de tailles différentes (ex. : notes de deux classes de 30 et 35 élèves).
\item Utilisation standard dans les domaines techniques : normes de qualité, courbes de croissance, analyses de fiabilité.
\end{itemize}
\end{aretenir}

\begin{exemplebox}[Courbes en pourcentages — seuils directs]
Pour la classe de 31 élèves, $N=31$. Les fréquences cumulées croissantes (en $\%$) aux bornes supérieures sont :
\[
\begin{array}{c|cccccccccc}
\text{Borne sup.} & 2 & 4 & 6 & 8 & 10 & 12 & 14 & 16 & 18 & 20 \\ \hline
FCC\ (\%) & 3,2 & 9,7 & 19,4 & 35,5 & 58,1 & 77,4 & 90,3 & 96,8 & 100 & 100
\end{array}
\]
Sur le graphique en $\%$, on lit directement :
\begin{itemize}
\item Médiane ($50\%$) $\approx 9,5$ (même abscisse qu'en effectifs).
\item Premier quartile ($25\%$) $\approx 6,7$.
\item Troisième quartile ($75\%$) : on cherche $75\%$ sur l'axe vertical $\to$ abscisse $\approx 11,7$.
\end{itemize}
\end{exemplebox}

\courssub{Exemple complet : notes d'un contrôle technique}

\begin{exoresolu}[Étude complète avec courbes cumulatives]
\textbf{Énoncé.} Un formateur note 40 apprentis sur 20 lors d'un examen pratique. La série groupée est :

\begin{center}
\begin{tabular}{|c|c|c|c|c|c|c|}\hline
Classes (notes) & $[0;4[$ & $[4;8[$ & $[8;12[$ & $[12;16[$ & $[16;20]$ \\ \hline
Effectifs $n_i$ & 3 & 7 & 14 & 12 & 4 \\ \hline
\end{tabular}
\end{center}

\begin{enumerate}
\item Construire le tableau des effectifs et fréquences cumulés (croissants et décroissants).
\item Tracer les deux courbes cumulatives en effectifs sur un même repère.
\item Estimer graphiquement la médiane.
\item Déterminer graphiquement le nombre d'apprentis ayant obtenu au moins 14/20.
\item Exprimer en pourcentage la proportion d'apprentis ayant moins de 10/20.
\end{enumerate}

\textbf{Solution.}

\noindent \textbf{1. Tableaux des effectifs cumulés.}

\begin{center}
\begin{tabular}{|c|c|c|c|c|c|}\hline
Classes & $[0;4[$ & $[4;8[$ & $[8;12[$ & $[12;16[$ & $[16;20]$ \\ \hline
$n_i$ & 3 & 7 & 14 & 12 & 4 \\ \hline
$ECC_i$ & 3 & 10 & 24 & 36 & 40 \\ \hline
$FCC_i\ (\%)$ & 7,5 & 25 & 60 & 90 & 100 \\ \hline
$ECD_i$ & 40 & 37 & 30 & 16 & 4 \\ \hline
$FCD_i\ (\%)$ & 100 & 92,5 & 75 & 40 & 10 \\ \hline
\end{tabular}
\end{center}

\noindent \textbf{2. Tracé des courbes.}
\begin{itemize}
\item \textbf{Croissante :} points $(0;0)$, $(4;3)$, $(8;10)$, $(12;24)$, $(16;36)$, $(20;40)$.
\item \textbf{Décroissante :} points $(0;40)$, $(4;37)$, $(8;30)$, $(12;16)$, $(16;4)$, $(20;0)$.
\end{itemize}

\begin{popfigure}
\begin{tikzpicture}
\begin{axis}[
    width=\linewidth,
    height=0.55\linewidth,
    xlabel={Notes},
    ylabel={Effectifs cumulés},
    xmin=0, xmax=20,
    ymin=0, ymax=42,
    xtick={0,4,8,12,16,20},
    ytick={0,5,10,15,20,25,30,35,40},
    grid=major,
    grid style={popline},
    axis lines=middle,
    enlargelimits={abs=0.5},
    clip=false,
    legend style={at={(0.02,0.98)}, anchor=north west, fill=white, draw=popline},
]
% Croissante
\addplot[popblue, very thick, mark=*, mark size=2.5, mark options={fill=popblue, draw=popdark}] coordinates {
    (0,0) (4,3) (8,10) (12,24) (16,36) (20,40)
};
\addlegendentry{Croissante (ECC)}
% Décroissante
\addplot[poporange, very thick, mark=square*, mark size=2.5, mark options={fill=poporange, draw=popdark}] coordinates {
    (0,40) (4,37) (8,30) (12,16) (16,4) (20,0)
};
\addlegendentry{Décroissante (ECD)}
% Intersection
\addplot[only marks, mark=*, mark size=4, popgreen] coordinates {(10.7,20)};
\node[popgreen, anchor=south west] at (10.7,20) {\textbf{I}};
\addplot[dashed, popgreen, thick] coordinates {(10.7,0) (10.7,20)};
\node[popgreen, anchor=north] at (10.7,0) {\small $\tilde{x} \approx 10,7$};
\addplot[dashed, popmuted, thin] coordinates {(0,20) (20,20)};
\node[popmuted, anchor=east] at (0,20) {\small $N/2 = 20$};
% Reading for x=14 (at least 14)
\addplot[dashed, poppurple, thick] coordinates {(14,0) (14,30)};
\addplot[dashed, poppurple, thick] coordinates {(0,30) (14,30)};
\node[poppurple, anchor=south west] at (14,30) {\small $(14;\,30)$};
\node[poppurple, anchor=north] at (14,0) {\small $14$};
\node[poppurple, anchor=east] at (0,30) {\small $30$};
% Reading for x=10 (less than 10)
\addplot[dashed, poppink, thick] coordinates {(10,0) (10,17)};
\addplot[dashed, poppink, thick] coordinates {(0,17) (10,17)};
\node[poppink, anchor=south west] at (10,17) {\small $(10;\,17)$};
\node[poppink, anchor=north] at (10,0) {\small $10$};
\node[poppink, anchor=east] at (0,17) {\small $17$};
\end{axis}
\end{tikzpicture}
\legende{Courbes cumulatives pour les 40 apprentis. Intersection I $\to$ médiane $\approx 10,7$. Lecture à $x=14$ : ECC $=30 \Rightarrow$ 10 apprentis $\ge 14$. Lecture à $x=10$ : ECC $\approx 17 \Rightarrow$ 17 apprentis $< 10$ ($42,5\%$).}
\end{popfigure}

\noindent \textbf{3. Médiane graphique.} Les deux courbes se coupent au point d'ordonnée $20 = N/2$. L'abscisse du point d'intersection est $\approx 10,7$. \textbf{Estimation de la médiane : $10,7/20$.}

\noindent \textbf{4. Apprentis ayant au moins 14/20.} On lit sur la courbe \textbf{décroissante} (ou on calcule $N - ECC(14)$). À $x=14$ (dans la classe $[12;16[$), interpolation sur la décroissante entre $(12;16)$ et $(16;4)$ : pente $= \frac{4-16}{16-12} = -3$. Pour $x=14$, $ECD = 16 + (-3)\times(14-12) = 10$. \textbf{10 apprentis ont $\ge 14/20$.}

\noindent \textbf{5. Proportion ayant moins de 10/20.} Sur la courbe croissante, à $x=10$ (dans $[8;12[$), interpolation entre $(8;10)$ et $(12;24)$ : pente $= \frac{24-10}{12-8} = 3,5$. $ECC(10) = 10 + 3,5 \times (10-8) = 17$. Fréquence $= \frac{17}{40} \times 100 = 42,5\%$. \textbf{42,5\% des apprentis ont moins de 10/20.}
\tcblower
\textbf{Remarques.} La médiane exacte (calculée en Leçon 16) sera dans la classe $[8;12[$ car $ECC(8)=10 < 20 < ECC(12)=24$. L'estimation graphique $10,7$ est très proche de la valeur exacte. Pour la question 4, on aurait pu lire sur la croissante : $ECC(14) \approx 30 \Rightarrow N - 30 = 10$.
\end{exoresolu}

\begin{savaistu}[Les courbes de croissance des enfants au Cameroun]
Les carnets de santé utilisés dans les centres de santé camerounais (selon les standards OMS) contiennent des \textbf{courbes de fréquences cumulées} appelées \emph{percentiles}. Pour un âge donné, la courbe donne le poids (ou la taille) en fonction du percentile : le percentile 50 est la médiane, le percentile 3 correspond à « 3\% des enfants ont un poids inférieur ». Un enfant dont le point (âge, poids) se situe sous la courbe du percentile 3 est considéré en \emph{insuffisance pondérale} et nécessite une prise en charge nutritionnelle. C'est une application directe, vitale, des courbes cumulatives en pourcentage !
\end{savaistu}

\courssub{Synthèse des points clés}

\begin{aretenir}[Ce qu'il faut retenir sur les courbes cumulatives]
\begin{itemize}
\item \textbf{Croissante :} points aux \textbf{bornes supérieures} $(x_{sup}; ECC)$, départ à $(\text{borne inf. 1re classe}; 0)$. Monte de 0 à $N$.
\item \textbf{Décroissante :} points aux \textbf{bornes inférieures} $(x_{inf}; ECD)$, départ à $(\text{borne inf. 1re classe}; N)$, fin à $(\text{borne sup. dernière classe}; 0)$. Descend de $N$ à 0.
\item \textbf{Lecture graphique :} interpolation linéaire sur les segments pour toute valeur intermédiaire.
\item \textbf{Médiane graphique :} abscisse de l'intersection des deux courbes (ordonnée $\approx N/2$). Propriété admise en 1re Tech.
\item \textbf{Version \% :} même forme, ordonnées de 0\% à 100\%. Lecture directe des quartiles (25\%, 50\%, 75\%).
\item \textbf{Applications techniques :} contrôle qualité (seuils de conformité), maintenance (durée de vie), gestion (délais), santé (courbes de croissance OMS).
\end{itemize}
\end{aretenir}

\begin{exercice}[Entraînement : courbes cumulatives en autonomie]
Une entreprise de BTP a relevé les durées (en jours) de 50 chantiers terminés :
\begin{center}
\begin{tabular}{|c|c|c|c|c|c|}\hline
Durées (jours) & $[0;10[$ & $[10;20[$ & $[20;30[$ & $[30;40[$ & $[40;50]$ \\ \hline
Nombre de chantiers & 5 & 12 & 18 & 10 & 5 \\ \hline
\end{tabular}
\end{center}
\begin{enumerate}
\item Compléter le tableau des effectifs et fréquences cumulés croissants et décroissants.
\item Tracer les deux courbes cumulatives en effectifs sur un même repère.
\item Estimer graphiquement la médiane des durées.
\item Déterminer graphiquement le nombre de chantiers ayant duré \textbf{plus de 35 jours}.
\item Quelle proportion de chantiers a duré \textbf{au maximum 25 jours} ? (Réponse en $\%$).
\end{enumerate}
\end{exercice}

\begin{corrige}[Exercice : courbes cumulatives]
\begin{enumerate}
\item Tableau :
\begin{center}
\begin{tabular}{|c|c|c|c|c|c|}\hline
Classes & $[0;10[$ & $[10;20[$ & $[20;30[$ & $[30;40[$ & $[40;50]$ \\ \hline
$n_i$ & 5 & 12 & 18 & 10 & 5 \\ \hline
$ECC_i$ & 5 & 17 & 35 & 45 & 50 \\ \hline
$FCC_i\ (\%)$ & 10 & 34 & 70 & 90 & 100 \\ \hline
$ECD_i$ & 50 & 45 & 33 & 15 & 5 \\ \hline
$FCD_i\ (\%)$ & 100 & 90 & 66 & 30 & 10 \\ \hline
\end{tabular}
\end{center}
\item Courbes : Croissante $(0;0)$, $(10;5)$, $(20;17)$, $(30;35)$, $(40;45)$, $(50;50)$. Décroissante $(0;50)$, $(10;45)$, $(20;33)$, $(30;15)$, $(40;5)$, $(50;0)$.
\item Intersection vers l'ordonnée $25$ ($N/2$). Abscisse $\approx 24$ jours. Médiane $\approx 24$ jours.
\item Plus de 35 jours $\Rightarrow$ courbe décroissante à $x=35$ (classe $[30;40[$). Interpolation : $(30;15)$ à $(40;5)$, pente $-1$. $ECD(35) = 15 - 1\times5 = 10$. \textbf{10 chantiers}.
\item Au max 25 jours $\Rightarrow$ courbe croissante à $x=25$ (classe $[20;30[$). Interpolation : $(20;17)$ à $(30;35)$, pente $1,8$. $ECC(25) = 17 + 1,8\times5 = 26$. Fréquence $= 26/50 = 52\%$. \textbf{52\% des chantiers}.
\end{enumerate}
\end{corrige}

\coursec{Étude complète d'une situation concrète}

Après avoir maîtrisé la construction des histogrammes et des courbes cumulatives, nous sommes désormais en mesure de mener une étude statistique complète, de la collecte des données brutes jusqu'à la rédaction d'une conclusion argumentée. C'est l'aboutissement de cette leçon : transformer des nombres en informations exploitables pour la prise de décision technique ou gestionnaire.

\courssub{La démarche en 5 étapes : de la donnée brute à la décision}

Toute étude statistique sérieuse suit un protocole rigoureux. En classe de Première technique, nous l'appliquons systématiquement en cinq temps forts.

\begin{methode}[Fiche de démarche complète — Étude d'une série groupée en classes]
\textbf{Étape 1 : Regrouper les données brutes en classes.}
\begin{itemize}
    \item Déterminer l'étendue $E = x_{\max} - x_{\min}$.
    \item Choisir un nombre de classes $k$ (généralement entre 5 et 15, règle de Sturges : $k \approx 1 + 3,3 \log_{10}(N)$).
    \item Fixer l'amplitude $a = \frac{E}{k}$ (arrondie à une valeur « commode » : 1, 2, 5, 10, 20, 50, 100...).
    \item Définir les bornes des classes (souvent de la forme $[a ; b[$ sauf la dernière $[a ; b]$) sans chevauchement ni trou.
\end{itemize}

\textbf{Étape 2 : Compléter le tableau statistique.}
\begin{itemize}
    \item Compter les effectifs $n_i$ pour chaque classe.
    \item Calculer les effectifs cumulés croissants $N_i$ et décroissants $N'_i$.
    \item Calculer les fréquences $f_i = \frac{n_i}{N}$ (en \%, $f_i \times 100$) et les fréquences cumulées $F_i = \frac{N_i}{N}$.
    \item Vérifier : $\sum n_i = N$ (effectif total) et $\sum f_i = 1$ (ou 100\%).
\end{itemize}

\textbf{Étape 3 : Représenter graphiquement.}
\begin{itemize}
    \item \textbf{Histogramme :} Sur l'axe des abscisses, les bornes des classes (échelle réelle). Sur l'axe des ordonnées, la \textbf{densité} $d_i = \frac{f_i}{a_i}$ (ou l'effectif $n_i$ si amplitudes égales). Surface du rectangle = fréquence de la classe.
    \item \textbf{Courbes cumulatives :} Polygone des fréquences cumulées croissantes (points $(b_i ; F_i)$ où $b_i$ est la borne supérieure de la classe) et/ou décroissantes.
\end{itemize}

\textbf{Étape 4 : Lire et extraire les indicateurs clés.}
\begin{itemize}
    \item \textbf{Classe modale :} Classe de densité maximale (rectangle le plus haut sur l'histogramme à amplitudes égales, ou de plus grande surface si amplitudes inégales).
    \item \textbf{Seuils (quartiles, déciles, médiane) :} Lecture graphique sur la courbe cumulative croissante (abscisse pour une ordonnée donnée) ou calcul par interpolation linéaire dans la classe contenant le seuil.
    \item \textbf{Pourcentages :} Lecture directe sur les fréquences cumulées (ex : « Quel \% de salaires $< 150\,000$ FCFA ? »).
\end{itemize}

\textbf{Étape 5 : Conclure par une phrase argumentée.}
\begin{itemize}
    \item Reformuler la question initiale.
    \item Apporter la réponse chiffrée (indicateurs lus ou calculés).
    \item Interpréter dans le contexte réel (atelier, gestion, production).
    \item \textbf{Exemple :} « La classe modale des salaires est $[100\,000 ; 150\,000[$ ; 50\% des employés gagnent moins de 142\,000 FCFA (médiane) ; l'entreprise doit revoir sa grille pour réduire l'écart interquartile. »
\end{itemize}
\end{methode}

\begin{popfigure}
\begin{tikzpicture}[
    node distance=1.2cm and 2.5cm,
    every node/.style={font=\small},
    box/.style={rectangle, draw=popdark, thick, rounded corners, fill=popcream, minimum width=3.2cm, minimum height=1.1cm, align=center},
    arrow/.style={-Stealth, thick, popblue},
    decision/.style={diamond, draw=poporange, thick, aspect=2, fill=poporangeL, minimum width=2.5cm, minimum height=1.2cm, align=center}
]
\node[box, align=center] (start){Données brutes\\ (liste de $N$ valeurs)};
\node[box, below of=start, align=center] (step1){Étape 1 : Regroupement\\ Choix classes, amplitude, bornes};
\node[box, below of=step1, align=center] (step2){Étape 2 : Tableau complet\\ $n_i, N_i, f_i, F_i, \text{centres } x_i$};
\node[decision, below of=step2, align=center] (check){Vérifications :\\ $\sum n_i = N$ ? $\sum f_i = 1$ ?};
\node[box, below left=1.2cm and 1.5cm of check, align=center] (step3a){Étape 3a : Histogramme\\ Densités $d_i = f_i/a_i$};
\node[box, below right=1.2cm and 1.5cm of check, align=center] (step3b){Étape 3b : Courbes cumulées\\ Polygones $F_i$ croiss./décroiss.};
\node[box, below=1.5cm of step3a.south west, anchor=north east, xshift=-1.5cm, align=center] (step4){Étape 4 : Lecture indicateurs\\ Classe modale, médiane, quartiles, \%};
\node[box, below=1.5cm of step3b.south east, anchor=north west, xshift=1.5cm, align=center] (step5){Étape 5 : Conclusion rédigée\\ Contexte + Chiffres + Décision};
\draw[arrow] (start) -- (step1);
\draw[arrow] (step1) -- (step2);
\draw[arrow] (step2) -- (check);
\draw[arrow] (check) -| (step3a);
\draw[arrow] (check) -| (step3b);
\draw[arrow] (step3a) |- (step4);
\draw[arrow] (step3b) |- (step4);
\draw[arrow] (step4) -- (step5);
\end{tikzpicture}
\legende{Organigramme de la démarche complète en 5 étapes pour l'étude d'une série statistique groupée en classes.}
\end{popfigure}

\courssub{Rédiger une conclusion argumentée : le triptyque Contexte — Chiffres — Interprétation}

Au Probatoire, une conclusion sans chiffres est vide ; des chiffres sans phrase sont muets. La conclusion doit être une \textbf{phrase unique, complète, contextualisée}.

\begin{exoresolu}[Synthèse rédigée — Étude des notes du Proviseur]
\textbf{Contexte :} Le proviseur d'un lycée technique a relevé les notes sur 20 de 120 élèves de Première à l'épreuve de mathématiques du 1er trimestre. La série a été groupée en classes d'amplitude 2.

\textbf{Tableau récapitulatif (extraits) :}
\begin{center}
\begin{tabularx}{\linewidth}{|c|X|c|c|c|}\hline
Classe & Centre $x_i$ & Effectif $n_i$ & Fréquence $f_i$ & Fréquence cumulée $F_i$ \\ \hline
$[0 ; 2[$ & 1 & 2 & 1,7\% & 1,7\% \\
$[2 ; 4[$ & 3 & 5 & 4,2\% & 5,8\% \\
$[4 ; 6[$ & 5 & 12 & 10,0\% & 15,8\% \\
$[6 ; 8[$ & 7 & 18 & 15,0\% & 30,8\% \\
$[8 ; 10[$ & 9 & \textbf{28} & \textbf{23,3\%} & 54,2\% \\
$[10 ; 12[$ & 11 & 22 & 18,3\% & 72,5\% \\
$[12 ; 14[$ & 13 & 16 & 13,3\% & 85,8\% \\
$[14 ; 16[$ & 15 & 10 & 8,3\% & 94,2\% \\
$[16 ; 18[$ & 17 & 5 & 4,2\% & 98,3\% \\
$[18 ; 20]$ & 19 & 2 & 1,7\% & 100\% \\ \hline
\textbf{Total} & & \textbf{120} & \textbf{100\%} & \\ \hline
\end{tabularx}
\end{center}

\textbf{Indicateurs lus (graphiquement ou par interpolation linéaire) :}
\begin{itemize}
    \item \textbf{Classe modale :} $[8 ; 10[$ (effectif max 28, fréquence 23,3\%).
    \item \textbf{Médiane (50\% = 60 élèves) :} Dans la classe $[8 ; 10[$. $N_{\text{prec}}=37$, $n_{\text{med}}=28$, $h=2$.
    $M_e \approx 8 + \frac{60 - 37}{28} \times 2 \approx 9,6$.
    \item \textbf{1er quartile $Q_1$ (25\% = 30 élèves) :} Dans $[6 ; 8[$. $N_{\text{prec}}=19$, $n=18$, $h=2$.
    $Q_1 \approx 6 + \frac{30 - 19}{18} \times 2 \approx 7,2$.
    \item \textbf{3ème quartile $Q_3$ (75\% = 90 élèves) :} Dans $[12 ; 14[$. $N_{\text{prec}}=87$, $n=16$, $h=2$.
    $Q_3 \approx 12 + \frac{90 - 87}{16} \times 2 \approx 12,4$.
    \item \textbf{Pourcentage de notes $< 10$ :} $F_{[8;10[} = 54,2\%$.
    \item \textbf{Pourcentage de notes $\ge 12$ (seuil admission) :} $1 - F_{[10;12[} = 1 - 0,725 = 27,5\%$.
\end{itemize}

\begin{aretenir}[Conclusion type attendue à l'examen]
« L'analyse des notes de mathématiques des 120 élèves de Première révèle une \textbf{classe modale située entre 8 et 10 sur 20} (23,3\% des élèves), traduisant une concentration des résultats autour de la moyenne. La \textbf{médiane est proche de 9,6}, indiquant que la moitié des élèves a obtenu moins de 9,6. L'écart interquartile ($Q_3 - Q_1 \approx 5,2$ points) montre une dispersion modérée. Cependant, \textbf{seuls 27,5\% des élèves atteignent la note de 12}, seuil souvent requis pour une admission sereine en Terminale ; une action pédagogique de remise à niveau s'impose pour le groupe des 54,2\% d'élèves en dessous de 10. »
\end{aretenir}
\end{exoresolu}

\courssub{Critiquer un regroupement : les pièges à éviter}

Un mauvais regroupement fausse toute l'analyse. Voici les trois erreurs cardinales à détecter et corriger.

\begin{attention}[Erreurs fréquentes de regroupement — Check-list de contrôle]
\begin{enumerate}
    \item \textbf{Amplitude mal choisie (trop grande ou trop petite).}
    \begin{itemize}
        \item \emph{Trop grande (ex: 2 classes pour 200 données) :} Perte totale d'information, forme de la distribution masquée, classe modale non significative.
        \item \emph{Trop petite (ex: amplitude 0,5 pour des notes entières sur 20) :} Beaucoup de classes vides ($n_i=0$), histogramme « en peigne » illisible, instabilité de la classe modale.
        \item \emph{Règle pratique :} Viser 8 à 12 classes non vides. Amplitude « ronde » (1, 2, 5, 10, 20, 50, 100...).
    \end{itemize}
    \item \textbf{Classes qui se chevauchent ou laissent des trous.}
    \begin{itemize}
        \item \emph{Chevauchement :} $[0 ; 10[$, $[5 ; 15[$, $[10 ; 20[$. Une valeur 7 appartient à deux classes $\rightarrow$ effectif total $> N$.
        \item \emph{Trous :} $[0 ; 10[$, $[11 ; 20[$. La valeur 10,5 n'appartient à aucune classe $\rightarrow$ effectif total $< N$.
        \item \emph{Correction :} Notation standard $[a ; b[$ pour toutes les classes sauf la dernière $[a ; b]$. Borne supérieure de la classe $i$ = borne inférieure de la classe $i+1$.
    \end{itemize}
    \item \textbf{Effectif total incorrect.}
    \begin{itemize}
        \item Oubli de lignes (ex: classe « 0 » ou classe « valeurs extrêmes » omises).
        \item Erreur de saisie dans le tableur (doublons, suppression accidentelle).
        \item \emph{Test systématique :} $\sum n_i \stackrel{?}{=} N_{\text{initial}}$. Si $\neq$, tout le tableau est faux.
    \end{itemize}
\end{enumerate}
\end{attention}

\courssub{Comparer deux séries groupées : la règle des fréquences}

En milieu professionnel (comparaison de deux machines, deux classes, deux sites), les effectifs totaux $N_1$ et $N_2$ sont presque toujours différents. \textbf{Comparer les effectifs bruts $n_i$ est une erreur grave.} On doit comparer les \textbf{fréquences $f_i$ (ou pourcentages)}.

\begin{exemplebox}[Comparaison de deux classes de Première — Notes de Mathématiques]
\textbf{Situation :} Deux classes passent le même contrôle.
\begin{itemize}
    \item \textbf{Classe A (1re TMA) :} $N_A = 35$ élèves.
    \item \textbf{Classe B (1re TEB) :} $N_B = 28$ élèves.
\end{itemize}
Les notes sont groupées en classes d'amplitude 2 : $[0;4[$, $[4;8[$, $[8;12[$, $[12;16[$, $[16;20]$.

\begin{center}
\begin{tabular}{|c|c|c|c|c|c|}\hline
Classe & \multicolumn{2}{c|}{Effectifs bruts} & \multicolumn{2}{c|}{Fréquences (\%)} & Analyse \\ \hline
& $n_{i,A}$ & $n_{i,B}$ & $f_{i,A}$ & $f_{i,B}$ & \\ \hline
$[0 ; 4[$ & 3 & 2 & 8,6\% & 7,1\% & Faibles résultats similaires \\ \hline
$[4 ; 8[$ & 8 & 6 & 22,9\% & 21,4\% & Niveau moyen comparable \\ \hline
$[8 ; 12[$ & \textbf{14} & \textbf{10} & \textbf{40,0\%} & \textbf{35,7\%} & \textbf{Classe modale pour les deux} \\ \hline
$[12 ; 16[$ & 7 & 6 & 20,0\% & 21,4\% & Bons élèves \\ \hline
$[16 ; 20]$ & 3 & 4 & 8,6\% & 14,3\% & \textbf{Classe B a plus d'excellents} \\ \hline
\textbf{Total} & \textbf{35} & \textbf{28} & \textbf{100\%} & \textbf{100\%} & \\ \hline
\end{tabular}
\end{center}

\textbf{Erreur de débutant :} « La classe A a 14 élèves dans la modale contre 10 pour la classe B, donc la classe A est meilleure. »
\textbf{Analyse correcte (fréquences) :} La classe modale est $[8;12[$ pour les deux (40,0\% vs 35,7\%). La différence majeure est dans la dernière classe : \textbf{14,3\% d'excellents en Classe B contre 8,6\% en Classe A}. La classe B a un profil plus « bipolaire » (plus de très bons, mais effectif total plus faible).
\end{exemplebox}

\begin{popfigure}
\begin{tikzpicture}
\begin{groupplot}[
    group style={group size=2 by 1, horizontal sep=2.5cm, vertical sep=0pt},
    width=0.45\linewidth,
    height=6cm,
    ymin=0, ymax=45,
    xtick={1,2,3,4,5},
    xticklabels={$[0;4[$, $[4;8[$, $[8;12[$, $[12;16[$, $[16;20]$},
    x tick label style={rotate=45, anchor=east, font=\footnotesize},
    ylabel={Effectif},
    ylabel style={font=\small},
    title style={font=\normalsize\bfseries, yshift=0.3cm},
    enlarge x limits=0.15,
    grid=major, grid style={popline},
]
\nextgroupplot[ybar=0pt, bar width=12pt, title={Classe A (TMA) — $N=35$}]
\addplot[popblue, fill=popblue!60] coordinates {(1,3) (2,8) (3,14) (4,7) (5,3)};
\nextgroupplot[ybar=0pt, bar width=12pt, title={Classe B (TEB) — $N=28$}]
\addplot[poporange, fill=poporange!60] coordinates {(1,2) (2,6) (3,10) (4,6) (5,4)};
\end{groupplot}
\end{tikzpicture}
\legende{Comparaison de deux histogrammes en effectifs bruts. Visuellement, la classe A semble « plus grosse » partout à cause de $N_A > N_B$. Seule la comparaison des fréquences (pourcentages) est pertinente.}
\end{popfigure}

\courssub{Lien avec la Leçon 16 : Centre de classe, moyenne et médiane approchées}

Dans la leçon suivante (Caractéristiques de position et de dispersion), nous aurons besoin de calculer une \textbf{moyenne} et une \textbf{médiane} pour une série groupée. Comme les valeurs exactes à l'intérieur d'une classe sont inconnues, nous utilisons une \textbf{hypothèse de travail} :

\begin{propriete}[Hypothèse du centre de classe — Base du calcul approché]
Pour une série groupée en classes, on \textbf{remplace toutes les valeurs d'une classe par le centre (milieu) de cette classe} $x_i = \frac{\text{borne inf} + \text{borne sup}}{2}$.
\begin{itemize}
    \item \textbf{Moyenne approchée :} $\bar{x} \approx \frac{\sum n_i x_i}{\sum n_i} = \sum f_i x_i$.
    \item \textbf{Médiane approchée (par interpolation linéaire) :} Dans la classe médiane $[a ; b[$ d'amplitude $h$, d'effectif cumulé précédent $N_{\text{prec}}$ et d'effectif $n_{\text{med}}$ :
    \[ M_e \approx a + \frac{\frac{N}{2} - N_{\text{prec}}}{n_{\text{med}}} \times h \]
    (Équivalent en fréquences : $M_e \approx a + \frac{0,5 - F_{\text{prec}}}{f_{\text{med}}} \times h$).
\end{itemize}
\textbf{Validité :} Cette approximation est d'autant meilleure que l'amplitude des classes est petite et que la distribution est homogène à l'intérieur de la classe modale.
\end{propriete}

\begin{savaistu}[Les statistiques de l'INS Cameroun — Des classes pour piloter le pays]
L'Institut National de la Statistique (INS) du Cameroun publie régulièrement l'\textbf{Enquête Camerounaise Auprès des Ménages (ECAM)}. Les revenus, les dépenses de consommation, les tailles des exploitations agricoles (cacao, café, coton) y sont \textbf{systématiquement présentés en classes} (ex: tranches de revenus mensuels : $[0 ; 50\,000[$, $[50\,000 ; 100\,000[$, ...).
\begin{itemize}
    \item \textbf{Pourquoi ?} Confidentialité (anonymat), lisibilité (pas de liste de 10 000 ménages), comparabilité dans le temps.
    \item \textbf{Usage concret :} La \textbf{classe modale des revenus} identifie le « Camerounais moyen » ; la \textbf{médiane} (souvent bien inférieure à la moyenne à cause des très hauts revenus) guide le SMIG ; les \textbf{quartiles/quintiles} définissent les seuils de pauvreté (premier quintile) et de classe moyenne.
    \item \textbf{Lien technique :} Les ingénieurs de l'ENEO ou de la CDC (Cameroon Development Corporation) utilisent ces mêmes méthodes (classes, histogrammes, courbes de Lorenz) pour dimensionner les transformateurs électriques ou prévoir les besoins en engrais par hectare de plants de cacao.
\end{itemize}
\end{savaistu}

\courssub{Applications concrètes : Trois situations professionnelles types}

\begin{enumerate}
    \item \textbf{Répartition des salaires (Gestion / Comptabilité).}
    Données : 250 employés, salaires de 45\,000 à 1\,200\,000 FCFA.
    \textbf{Regroupement :} Amplitude 100\,000 FCFA (classes $[40k;140k[$, $[140k;240k[$, ...).
    \textbf{Enjeu :} Identifier la classe modale (salaire le plus fréquent), calculer la médiane (négociation syndicale), repérer les écarts (équité). L'histogramme révèle souvent une distribution asymétrique à droite (quelques cadres très bien payés tirent la moyenne vers le haut).
    
    \item \textbf{Tailles de plants de cacaoyers (Agronomie / CDC-SODECAO).}
    Données : Hauteurs (cm) de 500 plants de 6 mois en pépinière.
    \textbf{Regroupement :} Amplitude 5 cm (classes $[20;25[$, $[25;30[$, ...).
    \textbf{Enjeu :} La classe modale donne le stade dominant. Si la médiane $< 30$ cm, le lot n'est pas prêt pour la mise en terre (repiquage). L'écart interquartile mesure l'homogénéité du lot (qualité du matériel végétal).
    
    \item \textbf{Consommation d'électricité ENEO (Maintenance / Exploitation).}
    Données : Consommation mensuelle (kWh) de 1\,200 abonnés basse tension dans un quartier de Douala.
    \textbf{Regroupement :} Amplitude 50 kWh (classes $[0;50[$, $[50;100[$, $[100;150[$, ...).
    \textbf{Enjeu :} La courbe cumulative permet de fixer les tranches de facturation (tarification sociale $< 110$ kWh, tarif normal, tarif majoré $> 300$ kWh). Le 9ème décile ($D_9$) identifie les « gros consommateurs » pour proposer des offres « Heures Creuses » ou détecter des fraudes (branchements sauvages).
\end{enumerate}

\begin{exercice}[Entraînement — Situation ENEO]
Un agent ENEO relève la consommation mensuelle (en kWh) de 40 abonnés :
\[ 85, 120, 45, 210, 330, 95, 150, 60, 180, 240, 75, 110, 290, 55, 135, 195, 310, 40, 165, 220, 90, 125, 260, 70, 140, 200, 300, 50, 100, 175, 230, 80, 115, 275, 65, 105, 185, 250, 35 \]

\begin{enumerate}
    \item Regrouper en classes d'amplitude 50 kWh en partant de 0. Compléter le tableau (effectifs, cumulés, fréquences, centres).
    \item Tracer l'histogramme (densités) et la courbe cumulative croissante.
    \item Déterminer la classe modale, la médiane (graphique et calcul), le 1er et 3ème quartile.
    \item Quel pourcentage d'abonnés bénéficie du tarif social ($< 110$ kWh) ?
    \item Rédiger la conclusion technique pour le chef de service « Facturation ».
\end{enumerate}
\end{exercice}

\begin{corrige}[Exercice — Situation ENEO]
\textbf{1. Tableau statistique (classes $[0;50[$, $[50;100[$, $[100;150[$, $[150;200[$, $[200;250[$, $[250;300[$, $[300;350]$).}
\begin{center}
\begin{tabular}{|c|c|c|c|c|c|}\hline
Classe & Centre $x_i$ & $n_i$ & $N_i$ (cumulé) & $f_i$ (\%) & $F_i$ (\%) \\ \hline
$[0 ; 50[$ & 25 & 3 & 3 & 7,5 & 7,5 \\
$[50 ; 100[$ & 75 & 10 & 13 & 25,0 & 32,5 \\
$[100 ; 150[$ & 125 & \textbf{8} & 21 & \textbf{20,0} & 52,5 \\
$[150 ; 200[$ & 175 & 6 & 27 & 15,0 & 67,5 \\
$[200 ; 250[$ & 225 & 5 & 32 & 12,5 & 80,0 \\
$[250 ; 300[$ & 275 & 4 & 36 & 10,0 & 90,0 \\
$[300 ; 350]$ & 325 & 4 & 40 & 10,0 & 100,0 \\ \hline
\textbf{Total} & & \textbf{40} & & \textbf{100,0} & \\ \hline
\end{tabular}
\end{center}
Vérification : $\sum n_i = 40 = N$. $\sum f_i = 100\%$.

\textbf{2. Graphiques :} 
\begin{itemize}
    \item \textbf{Histogramme :} Amplitude constante $a=50$. Densités $d_i = f_i / 50$ (avec $f_i$ en décimal). Exemple classe modale $[100;150[$ : $f=0,20$, $d = 0,20 / 50 = 0,004$. Les hauteurs des rectangles sont proportionnelles aux fréquences.
    \item \textbf{Courbe cumulative croissante :} Points $(50;0,075)$, $(100;0,325)$, $(150;0,525)$, $(200;0,675)$, $(250;0,80)$, $(300;0,90)$, $(350;1)$. Relier par segments.
\end{itemize}

\textbf{3. Indicateurs (calculs par interpolation linéaire) :}
\begin{itemize}
    \item \textbf{Classe modale :} $[100 ; 150[$ (fréquence max 20,0\%).
    \item \textbf{Médiane (50\% = 20ème valeur) :} Dans $[100 ; 150[$ ($N_{\text{prec}}=13$, $n_{\text{med}}=8$, $h=50$).
    $M_e = 100 + \frac{20 - 13}{8} \times 50 = 100 + 43,75 = \textbf{143,75 kWh}$.
    \item \textbf{$Q_1$ (25\% = 10ème valeur) :} Dans $[50 ; 100[$ ($N_{\text{prec}}=3$, $n=10$, $h=50$).
    $Q_1 = 50 + \frac{10 - 3}{10} \times 50 = 50 + 35 = \textbf{85 kWh}$.
    \item \textbf{$Q_3$ (75\% = 30ème valeur) :} Dans $[200 ; 250[$ ($N_{\text{prec}}=27$, $n=5$, $h=50$).
    $Q_3 = 200 + \frac{30 - 27}{5} \times 50 = 200 + 30 = \textbf{230 kWh}$.
\end{itemize}

\textbf{4. Tarif social ($< 110$ kWh) :} Lecture sur courbe cumulative à $x=110$ (dans $[100;150[$).
$F(100) = 0,325$. Fréquence classe $[100;150[$ : $f = 0,20$.
Interpolation : $F(110) \approx 0,325 + \frac{110-100}{50} \times 0,20 = 0,325 + 0,04 = 0,365$.
\textbf{36,5\% des abonnés} sont en tarif social.

\textbf{5. Conclusion :}
« L'étude de la consommation de 40 abonnés ENEO montre une \textbf{classe modale entre 100 et 150 kWh/mois} (20,0\% des clients). La \textbf{médiane est de 143,75 kWh} : la moitié des abonnés consomme moins de cette valeur. L'écart interquartile ($Q_3 - Q_1 = 145$ kWh) traduit une \textbf{forte dispersion} des comportements. \textbf{36,5\% des abonnés entrent dans la tranche tarif sociale ($< 110$ kWh)}. Le service facturation doit prévoir une communication ciblée vers les 20\% de gros consommateurs ($> 250$ kWh) pour l'option « Heures Creuses » et surveiller les 10\% au-delà de 300 kWh (risque fraude). »
\end{corrige}

\newpage

\coursec{Activité d'intégration}

\coursec{Leçon 11 : Séries statistiques regroupées en classes}

\courssub{Objectifs de l'activité}

Cette activité d'intégration mobilise \emph{l'ensemble} des savoir-faire de la leçon sur des données réelles collectées par les élèves eux-mêmes. À l'issue du travail, chaque élève doit être capable de :

\begin{itemize}
\item regrouper une série statistique brute en \cle{classes} d'amplitude donnée ;
\item construire un tableau complet : \cle{effectifs}, \cle{effectifs cumulés croissants} (ECC), \cle{effectifs cumulés décroissants} (ECD), \cle{fréquences} et \cle{fréquences cumulées croissantes} (FCC) ;
\item tracer un \cle{histogramme} et une \cle{courbe cumulative} sur papier millimétré ;
\item lire graphiquement la \cle{classe modale}, un effectif cumulé et une valeur médiane approchée ;
\item \textbf{rédiger et justifier} une démarche et une conclusion, compétence exigée au Probatoire technique.
\end{itemize}

\courssub{Situation}

\textbf{Enquête statistique dans la classe : les dépenses hebdomadaires de transport.}

Au lycée technique de Nkouloulou (Douala), la classe de Première technique compte 30 élèves. La plupart viennent au lycée en moto-taxi (\emph{bendskin}) ou en taxi collectif. Le conseiller d'orientation souhaite connaître le budget transport des élèves afin de proposer à l'administration l'ouverture d'une ligne de transport scolaire.

Chaque élève a relevé sa dépense hebdomadaire de transport (aller-retour, 5 jours), en francs CFA. Voici les 30 données brutes recueillies :

\begin{center}
\begin{tabular}{|c|c|c|c|c|c|}
\hline
\rowcolor{popblueL}
1200 & 2500 & 800 & 3000 & 1500 & 2200 \\
\hline
1750 & 900 & 2800 & 1300 & 2000 & 1600 \\
\hline
3500 & 1100 & 2400 & 750 & 1900 & 2700 \\
\hline
1400 & 2100 & 1000 & 3200 & 1800 & 1250 \\
\hline
2300 & 1500 & 950 & 2600 & 1700 & 1450 \\
\hline
\end{tabular}
\end{center}

Le proviseur demande : \emph{« Quelle est la dépense de transport la plus fréquente ? Combien d'élèves dépensent moins de 2000 FCFA ? À partir de quelle dépense se situe la moitié la plus aisée des élèves ? »} Pour répondre, la classe décide de regrouper les données en classes d'amplitude 500 FCFA, en commençant à 500 FCFA : $[500\,;\,1000[$, $[1000\,;\,1500[$, \dots

\courssub{Consignes}

Travail par groupes de 4 élèves. Matériel : papier millimétré, règle, calculatrice, crayons de couleur.

\begin{enumerate}
\item \textbf{Regroupement.} Ranger les 30 données dans les classes d'amplitude 500 FCFA. Justifier le nombre de classes retenu. Vérifier que la somme des effectifs vaut 30.
\item \textbf{Tableau complet.} Construire le tableau statistique comportant : les classes, les centres de classes, les effectifs $n_i$, les ECC, les ECD, les fréquences $f_i$ (en \%) et les FCC (en \%).
\item \textbf{Histogramme.} Tracer l'histogramme des effectifs (classes d'amplitudes égales). Identifier graphiquement la ou les \cle{classes modales}.
\item \textbf{Courbe cumulative.} Tracer la courbe des effectifs cumulés croissants (polygone des ECC). Expliquer pourquoi on place les points aux extrémités droites des classes.
\item \textbf{Exploitation.} Répondre aux trois questions du proviseur :
\begin{enumerate}
\item Quelle(s) est/sont la/les classe(s) modale(s) ?
\item Combien d'élèves dépensent moins de 2000 FCFA ? Quel pourcentage cela représente-t-il ?
\item À partir de quelle dépense se situe la moitié la plus aisée des élèves ? (Lire la médiane sur la courbe cumulative.)
\end{enumerate}
\item \textbf{Rédaction.} Chaque groupe rédige une conclusion de trois phrases au proviseur, en justifiant chaque réponse par le tableau ou un graphique.
\item \textbf{Comparaison.} En correction collective, comparer les résultats des groupes et discuter les éventuelles différences de lecture graphique.
\end{enumerate}

\courssub{Ressources à mobiliser}

\begin{aretenir}[Boîte à outils de la leçon]
\begin{itemize}
\item \textbf{Classe(s) modale(s)} : classe(s) ayant le plus grand effectif (amplitudes égales).
\item \textbf{Effectif cumulé croissant} en une borne $b$ : nombre de données \emph{strictement inférieures} à $b$ ; on l'obtient en additionnant les effectifs des classes situées avant $b$.
\item \textbf{Effectif cumulé décroissant} : nombre de données \emph{supérieures ou égales} à la borne inférieure de la classe.
\item \textbf{Fréquence} d'une classe : $f_i = \dfrac{n_i}{N}$, avec $N$ l'effectif total ; en pourcentage : $f_i \times 100$.
\item \textbf{Courbe cumulative croissante} : on place les points aux \emph{extrémités droites} des classes, d'ordonnées les ECC, en partant du point $(\text{première borne}\,;\,0)$.
\item \textbf{Médiane} (lecture graphique) : abscisse du point de la courbe cumulative d'ordonnée $\dfrac{N}{2}$.
\end{itemize}
\end{aretenir}

\courssub{Démarche et corrigé commenté}

\begin{exoresolu}[Corrigé complet de l'enquête]
Répondre aux consignes 1 à 6 avec les données de la classe.
\tcblower
\textbf{Étape 1 : regroupement en classes.}

La plus petite donnée est 750 FCFA, la plus grande 3500 FCFA. Avec des classes d'amplitude 500 FCFA à partir de 500, on obtient 6 classes : de $[500\,;\,1000[$ à $[3000\,;\,3500]$. On dépouille les données une à une (méthode des pointages) :

\begin{center}
\begin{tabular}{|c|c|c|c|c|c|c|}
\hline
\rowcolor{popblueL}
Classes (FCFA) & Centre $x_i$ & $n_i$ & ECC & ECD & $f_i$ (\%) & FCC (\%) \\
\hline
$[500\,;\,1000[$ & 750 & 4 & 4 & 30 & 13,3 & 13,3 \\
\hline
$[1000\,;\,1500[$ & 1250 & 7 & 11 & 26 & 23,3 & 36,7 \\
\hline
$[1500\,;\,2000[$ & 1750 & 7 & 18 & 19 & 23,3 & 60,0 \\
\hline
$[2000\,;\,2500[$ & 2250 & 5 & 23 & 12 & 16,7 & 76,7 \\
\hline
$[2500\,;\,3000[$ & 2750 & 4 & 27 & 7 & 13,3 & 90,0 \\
\hline
$[3000\,;\,3500]$ & 3250 & 3 & 30 & 3 & 10,0 & 100 \\
\hline
\rowcolor{popcream}
Total & & 30 & & & 100 & \\
\hline
\end{tabular}
\end{center}

\emph{Justification} : $4+7+7+5+4+3 = 30 = N$ : le dépouillement est correct. Les fréquences sont calculées par $f_i = \dfrac{n_i}{30}\times 100$ ; par exemple $f_2 = \dfrac{7}{30}\times 100 \approx 23,3\,\%$. Les ECC s'obtiennent en additionnant de haut en bas ($4$ ; $4+7=11$ ; $11+7=18$ ; \dots), les ECD de bas en haut ($3$ ; $3+4=7$ ; $7+5=12$ ; \dots).

\textbf{Étape 2 : histogramme.}

Les classes ont toutes la même amplitude (500 FCFA) : les hauteurs des rectangles sont donc proportionnelles aux effectifs. Échelle conseillée : 1 cm pour 500 FCFA en abscisse, 1 cm pour 1 élève en ordonnée (ou 1 cm pour 2 élèves).

\begin{popfigure}
\begin{tikzpicture}[x=0.018cm,y=0.55cm]
\draw[->] (400,0) -- (3700,0) node[below left]{FCFA};
\draw[->] (500,0) -- (500,8.5) node[below left]{effectif};
\foreach \x in {500,1000,1500,2000,2500,3000,3500} \draw (\x,0) node[below,rotate=45,anchor=east]{\scriptsize \x};
\foreach \y in {1,2,3,4,5,6,7} \draw (500,\y) node[left]{\scriptsize \y};
\fill[popblue!60,draw=popblue] (500,0) rectangle (1000,4);
\fill[poporange!60,draw=poporange] (1000,0) rectangle (1500,7);
\fill[poporange!60,draw=poporange] (1500,0) rectangle (2000,7);
\fill[popblue!60,draw=popblue] (2000,0) rectangle (2500,5);
\fill[popblue!60,draw=popblue] (2500,0) rectangle (3000,4);
\fill[popblue!60,draw=popblue] (3000,0) rectangle (3500,3);
\end{tikzpicture}
\legende{Histogramme des dépenses hebdomadaires de transport (en orange, les deux classes modales).}
\end{popfigure}

\textbf{Étape 3 : courbe cumulative croissante.}

On place les points $(500\,;\,0)$, $(1000\,;\,4)$, $(1500\,;\,11)$, $(2000\,;\,18)$, $(2500\,;\,23)$, $(3000\,;\,27)$, $(3500\,;\,30)$, puis on les relie par des segments. On place les points aux extrémités \emph{droites} car l'ECC en une borne compte toutes les données inférieures à cette borne.

\begin{popfigure}
\begin{tikzpicture}[x=0.018cm,y=0.24cm]
\draw[->] (400,0) -- (3700,0) node[below left]{FCFA};
\draw[->] (500,0) -- (500,33) node[below left]{ECC};
\foreach \x in {500,1000,1500,2000,2500,3000,3500} \draw (\x,0) node[below,rotate=45,anchor=east]{\scriptsize \x};
\foreach \y in {5,10,15,20,25,30} \draw (500,\y) node[left]{\scriptsize \y};
\draw[popcurve] (500,0) -- (1000,4) -- (1500,11) -- (2000,18) -- (2500,23) -- (3000,27) -- (3500,30);
\foreach \p in {(1000,4),(1500,11),(2000,18),(2500,23),(3000,27),(3500,30)} \fill[poporange] \p circle (0.35);
\draw[dashed,popgreen] (500,15) -- (1786,15) -- (1786,0) node[below]{\scriptsize $Me \approx 1786$};
\end{tikzpicture}
\legende{Courbe des effectifs cumulés croissants et lecture de la médiane.}
\end{popfigure}

\textbf{Étape 4 : réponses aux questions du proviseur.}

\begin{enumerate}
\item \textbf{Classes modales} : les classes $[1000\,;\,1500[$ et $[1500\,;\,2000[$ ont le plus grand effectif ($n_2 = n_3 = 7$) et les rectangles les plus hauts sur l'histogramme. La dépense la plus fréquente se situe donc entre 1000 et 2000 FCFA par semaine (deux classes modales).
\item \textbf{Moins de 2000 FCFA} : on lit l'ECC à la borne 2000 : $18$ élèves dépensent moins de 2000 FCFA, soit $\dfrac{18}{30}\times 100 = 60\,\%$ de la classe (on retrouve la FCC de la classe $[1500\,;\,2000[$).
\item \textbf{Moitié la plus aisée} : $\dfrac{N}{2} = 15$. Sur la courbe cumulative, l'horizontale d'ordonnée 15 coupe le polygone dans la classe $[1500\,;\,2000[$, à une abscisse d'environ 1786 FCFA (interpolation linéaire : $1500 + 500 \times \frac{15-11}{7} \approx 1786$). La moitié la plus aisée des élèves dépense donc \emph{au moins} environ 1790 FCFA par semaine (médiane $Me \approx 1790$ FCFA).
\end{enumerate}

\textbf{Étape 5 : conclusion rédigée (modèle).}

\emph{« Les dépenses hebdomadaires de transport les plus fréquentes se situent entre 1000 et 2000 FCFA, car les classes $[1000\,;\,1500[$ et $[1500\,;\,2000[$ ont le plus grand effectif (7 élèves chacune). 60\,\% des élèves (18 sur 30) dépensent moins de 2000 FCFA, comme le montre l'effectif cumulé croissant lu à la borne 2000. Enfin, la courbe cumulative indique une médiane d'environ 1790 FCFA : la moitié la plus aisée des élèves dépense au moins cette somme, ce qui justifie l'étude d'une ligne de transport scolaire à tarif inférieur. »}
\end{exoresolu}

\begin{attention}[Erreurs fréquentes à signaler en correction]
\begin{itemize}
\item Oublier de vérifier que la somme des effectifs vaut $N = 30$ : une donnée mal pointée fausse tout le tableau.
\item Placer les points de la courbe cumulative aux \emph{centres} des classes au lieu des extrémités droites.
\item Confondre ECC et ECD : l'ECC à 2000 compte ceux qui dépensent \emph{moins} de 2000 FCFA.
\item Conclure « 15 élèves » pour la médiane : 15 est l'\emph{effectif} $\frac{N}{2}$, la médiane est une \emph{dépense} lue sur l'axe des abscisses.
\item Écrire des fréquences dont la somme ne fait pas 100\,\% (arrondis à contrôler ; ici $13,3+23,3+23,3+16,7+13,3+10,0 = 99,9 \approx 100$).
\item Ne pas signaler qu'il y a deux classes modales quand les effectifs maximaux sont égaux.
\end{itemize}
\end{attention}

\courssub{Bilan de l'activité}

Cette enquête a permis de mettre en évidence que :

\begin{itemize}
\item les données brutes, trop nombreuses, ne « parlent » pas : le \cle{regroupement en classes} rend la série lisible et exploitable ;
\item le tableau des \cle{effectifs et fréquences cumulés} répond immédiatement à des questions du type « combien dépensent moins de \dots » ou « plus de \dots » ;
\item l'\cle{histogramme} visualise la ou les classes modales, tandis que la \cle{courbe cumulative} permet d'estimer la médiane par simple lecture graphique ;
\item une conclusion statistique rigoureuse s'appuie toujours sur un \emph{élément précis} du tableau ou du graphique : c'est cette justification écrite qui est évaluée au Probatoire technique.
\end{itemize}

\begin{savaistu}[Le transport scolaire, un vrai enjeu statistique]
Dans les grandes villes du Cameroun comme Douala et Yaoundé, les dépenses de transport représentent une part importante du budget des familles. Les services de l'urbanisme et les sociétés de transport utilisent exactement ces méthodes (regroupement en classes, courbes cumulatives) pour fixer les tarifs et décider de l'ouverture de nouvelles lignes : les statistiques que tu viens de pratiquer servent à prendre de vraies décisions.
\end{savaistu}

\newpage

\coursec{Méthodes et exercices résolus}

\courssub{Regroupement en classes : choix des bornes et amplitude}

\begin{methode}[Construire un tableau de regroupement en classes]
    \textbf{Objectif :} Transformer une série brute en série groupée en classes d'amplitude constante (ou variable).
    \begin{enumerate}
        \item \textbf{Étendue :} Calculer $E = x_{\max} - x_{\min}$.
        \item \textbf{Nombre de classes $k$ :} Choisir $k$ entre 5 et 15 (souvent $k \approx \sqrt{n}$ ou règle de Sturges $k = 1 + 3,3 \log_{10} n$). Arrondir à l'entier supérieur.
        \item \textbf{Amplitude $a$ :} $a = \frac{E}{k}$. Arrondir $a$ à une valeur « commode » (multiple de 1, 2, 5, 10, 0,1, 0,2, 0,5\ldots) pour faciliter la lecture.
        \item \textbf{Bornes des classes :} Choisir la borne inférieure de la première classe $b_0 \le x_{\min}$ (souvent un multiple de $a$). Les classes sont $[b_0; b_0+a[$, $[b_0+a; b_0+2a[$, \ldots, la dernière classe fermante $[b_{k-1}; b_k]$.
        \item \textbf{Effectifs :} Compter les valeurs dans chaque classe. Vérifier $\sum n_i = n$.
    \end{enumerate}
    \textbf{Repères :} Pour des mesures techniques (dimensions, tensions, temps), privilégier des bornes « rondes » (ex: 10,0 ; 10,5 ; 11,0 mm) plutôt que les valeurs exactes min/max.
\end{methode}

\begin{exoresolu}[Regroupement des diamètres de tiges tournées]
    \textbf{Énoncé :} Un tourneur contrôle le diamètre (en mm) de 40 tiges. Résultats :
    \[10,12; 10,14; 10,16; 10,18; 10,21; 10,23; 10,25; 10,28; 10,31; 10,34;\]
    \[10,36; 10,39; 10,42; 10,44; 10,46; 10,49; 10,51; 10,53; 10,55; 10,58;\]
    \[10,62; 10,64; 10,66; 10,69; 10,71; 10,73; 10,75; 10,78; 10,81; 10,83;\]
    \[10,85; 10,88; 10,91; 10,93; 10,95; 10,98; 11,01; 11,03; 11,05; 11,08.\]
    Regrouper ces données en classes d'amplitude 0,10 mm en partant de 10,10 mm. Compléter le tableau : classes, effectifs, fréquences (en \%), fréquences cumulées croissantes.
    \tcblower
    \textbf{Solution détaillée :}
    \begin{enumerate}
        \item \textbf{Vérification des paramètres :} $x_{\min}=10,12$, $x_{\max}=11,08$. $E = 0,96$. Amplitude imposée $a=0,10$. Borne de départ $b_0=10,10$. Nombre de classes $k=10$ (de 10,10 à 11,10).
        \item \textbf{Construction des classes :}
        \begin{align*}
            C_1 &: [10,10; 10,20[ \quad C_2 : [10,20; 10,30[ \quad C_3 : [10,30; 10,40[ \quad C_4 : [10,40; 10,50[ \\
            C_5 &: [10,50; 10,60[ \quad C_6 : [10,60; 10,70[ \quad C_7 : [10,70; 10,80[ \quad C_8 : [10,80; 10,90[ \\
            C_9 &: [10,90; 11,00[ \quad C_{10} : [11,00; 11,10] \text{ (dernière classe fermante)}
        \end{align*}
        \item \textbf{Comptage des effectifs $n_i$ :} On parcourt la liste brute. Chaque classe contient exactement 4 valeurs.
        \item \textbf{Fréquences $f_i = \frac{n_i}{40}$ :} En \%, $f_i(\%) = \frac{n_i}{40}\times 100 = 10,0\%$.
        \item \textbf{Fréquences cumulées croissantes $F_i$ :} $F_1=f_1$, $F_2=F_1+f_2$, \ldots, $F_{10}=100\%$.
    \end{enumerate}
    \begin{center}
    \begin{tabularx}{\linewidth}{|c|>{\centering\arraybackslash}X|>{\centering\arraybackslash}X|>{\centering\arraybackslash}X|>{\centering\arraybackslash}X|}\hline
    \rowcolor{popblueL}
    \textbf{Classes (mm)} & \textbf{Effectifs $n_i$} & \textbf{Fréquences $f_i$ (\%)} & \textbf{Fréq. cumulées $F_i$ (\%)} & \textbf{Milieu $x_i$} \\ \hline
    $[10,10; 10,20[$ & 4 & 10,0 & 10,0 & 10,15 \\ \hline
    $[10,20; 10,30[$ & 4 & 10,0 & 20,0 & 10,25 \\ \hline
    $[10,30; 10,40[$ & 4 & 10,0 & 30,0 & 10,35 \\ \hline
    $[10,40; 10,50[$ & 4 & 10,0 & 40,0 & 10,45 \\ \hline
    $[10,50; 10,60[$ & 4 & 10,0 & 50,0 & 10,55 \\ \hline
    $[10,60; 10,70[$ & 4 & 10,0 & 60,0 & 10,65 \\ \hline
    $[10,70; 10,80[$ & 4 & 10,0 & 70,0 & 10,75 \\ \hline
    $[10,80; 10,90[$ & 4 & 10,0 & 80,0 & 10,85 \\ \hline
    $[10,90; 11,00[$ & 4 & 10,0 & 90,0 & 10,95 \\ \hline
    $[11,00; 11,10]$ & 4 & 10,0 & 100,0 & 11,05 \\ \hline
    \textbf{Total} & \textbf{40} & \textbf{100,0} & & \\ \hline
    \end{tabularx}
    \end{center}
    \textbf{Conclusion :} La répartition est uniforme (distribution rectangulaire), ce qui est inhabituel pour un processus réel ; cela suggère une simulation ou un arrondi des mesures. Le tableau est prêt pour le tracé de l'histogramme.
\end{exoresolu}

\courssub{Effectifs et fréquences cumulés (croissants et décroissants)}

\begin{methode}[Calculer et interpréter les effectifs et fréquences cumulés]
    \textbf{Données :} Tableau de classes $[a_i; b_i[$ avec effectifs $n_i$, total $N$.
    \begin{enumerate}
        \item \textbf{Croissants ( « inférieurs à » ) :}
        \begin{itemize}
            \item Effectif cumulé croissant $N_i = \sum_{k=1}^{i} n_k$. $N_k = N$.
            \item Fréquence cumulée croissante $F_i = \frac{N_i}{N}$. $F_k = 1$ (ou 100\%).
            \item \emph{Interprétation :} $N_i$ = nombre de valeurs $<$ borne sup de la classe $i$. $F_i$ = proportion de valeurs $<$ cette borne.
        \end{itemize}
        \item \textbf{Décroissants ( « supérieurs ou égaux à » ) :}
        \begin{itemize}
            \item Effectif cumulé décroissant $N'_i = \sum_{k=i}^{k} n_k$. $N'_1 = N$.
            \item Fréquence cumulée décroissante $F'_i = \frac{N'_i}{N}$. $F'_1 = 1$.
            \item \emph{Interprétation :} $N'_i$ = nombre de valeurs $\ge$ borne inf de la classe $i$.
        \end{itemize}
        \item \textbf{Formule de passage :} $N'_i = N - N_{i-1}$ (avec $N_0=0$). $F'_i = 1 - F_{i-1}$.
    \end{enumerate}
    \textbf{Attention :} Bien distinguer « effectif » (nombre, entier) et « fréquence » (proportion, décimal ou \%). Au Probatoire, préciser l'unité (\% ou nombre).
\end{methode}

\begin{exoresolu}[Cumulés croissants et décroissants : temps de cycle]
    \textbf{Énoncé :} Une étude de poste donne les temps de cycle (en secondes) pour 50 opérations :
    \begin{center}
    \begin{tabular}{|c|c|c|c|c|c|c|c|}\hline
    \rowcolor{popblueL} Classes (s) & $[10;12[$ & $[12;14[$ & $[14;16[$ & $[16;18[$ & $[18;20[$ & $[20;22[$ & $[22;24[$ \\ \hline
    Effectifs & 5 & 8 & 12 & 15 & 6 & 3 & 1 \\ \hline
    \end{tabular}
    \end{center}
    1. Compléter le tableau avec les fréquences (\%), les effectifs et fréquences cumulés croissants et décroissants.
    2. Combien d'opérations durent \emph{strictement moins de 18 s} ? \emph{Au moins 16 s} ?
    \tcblower
    \textbf{Solution détaillée :}
    \begin{enumerate}
        \item \textbf{Total $N=50$. Fréquences $f_i = n_i/50$ (en \% : $\times 2$).}
        \item \textbf{Croissants :} $N_1=5, N_2=13, N_3=25, N_4=40, N_5=46, N_6=49, N_7=50$.
        $F_i = N_i/50$ (en \% : $10\%, 26\%, 50\%, 80\%, 92\%, 98\%, 100\%$).
        \item \textbf{Décroissants :} $N'_7=1, N'_6=4, N'_5=10, N'_4=25, N'_3=37, N'_2=45, N'_1=50$.
        $F'_i = N'_i/50$ (en \% : $2\%, 8\%, 20\%, 50\%, 74\%, 90\%, 100\%$).
        \item \textbf{Vérification :} $N'_4 = N - N_3 = 50 - 25 = 25$. $F'_4 = 1 - F_3 = 1 - 0,5 = 0,5$. OK.
    \end{enumerate}
    \begin{center}
    \begin{tabularx}{\linewidth}{|c|>{\centering\arraybackslash}X|>{\centering\arraybackslash}X|>{\centering\arraybackslash}X|>{\centering\arraybackslash}X|>{\centering\arraybackslash}X|>{\centering\arraybackslash}X|}\hline
    \rowcolor{popblueL}
    \textbf{Classes} & \textbf{$n_i$} & \textbf{$f_i$(\%)} & \textbf{$N_i$ (croit)} & \textbf{$F_i$(\%)} & \textbf{$N'_i$ (décroit)} & \textbf{$F'_i$(\%)} \\ \hline
    $[10;12[$ & 5 & 10 & 5 & 10 & 50 & 100 \\ \hline
    $[12;14[$ & 8 & 16 & 13 & 26 & 45 & 90 \\ \hline
    $[14;16[$ & 12 & 24 & 25 & 50 & 37 & 74 \\ \hline
    $[16;18[$ & 15 & 30 & 40 & 80 & 25 & 50 \\ \hline
    $[18;20[$ & 6 & 12 & 46 & 92 & 10 & 20 \\ \hline
    $[20;22[$ & 3 & 6 & 49 & 98 & 4 & 8 \\ \hline
    $[22;24[$ & 1 & 2 & 50 & 100 & 1 & 2 \\ \hline
    \end{tabularx}
    \end{center}
    \textbf{Réponses aux questions :}
    \begin{itemize}
        \item « Strictement moins de 18 s » $\rightarrow$ valeurs dans $[10;18[$. Effectif cumulé croissant de la classe $[16;18[$ : $N_4 = 40$ opérations.
        \item « Au moins 16 s » $\rightarrow$ valeurs $\ge 16$. Effectif cumulé décroissant de la classe $[16;18[$ : $N'_4 = 25$ opérations.
    \end{itemize}
    \textbf{Conclusion :} 40 opérations sont $< 18$s ; 25 opérations sont $\ge 16$s. La médiane (classe où $F_i$ dépasse 50\%) est dans $[16;18[$.
\end{exoresolu}

\courssub{Histogramme d'une série groupée en classes}

\begin{methode}[Tracer l'histogramme (densité / aire proportionnelle à l'effectif)]
    \textbf{Règle fondamentale :} \cle{L'aire du rectangle = effectif (ou fréquence) de la classe}. Hauteur = \cle{Densité} $d_i = \frac{n_i}{a_i}$ (ou $\frac{f_i}{a_i}$).
    \begin{enumerate}
        \item \textbf{Axe des abscisses :} Échelle adaptée aux bornes des classes. Représenter les bornes exactes. Laisser un espace pour l'origine si $x_{\min} > 0$ (trait de rupture $\backslash\backslash$).
        \item \textbf{Axe des ordonnées :} Échelle pour les densités $d_i = n_i / a_i$. Unité : « effectifs par unité de $X$ » (ex: pièces/mm, opérations/seconde).
        \item \textbf{Rectangles :} Base = amplitude $a_i$ (largeur de la classe). Hauteur = densité $d_i$.
        \item \textbf{Classes d'amplitude inégale :} Indispensable d'utiliser la densité. Si amplitudes égales, hauteur $\propto$ effectif (on peut tracer directement les effectifs, mais l'axe doit être étiqueté « Densité » ou « Effectifs » avec précision).
        \item \textbf{Classe ouverte (ex: $[a; +\infty[$) :} Ne pas tracer (ou tracer une flèche) car aire infinie. Au Probatoire, on suppose généralement une amplitude finie pour la dernière classe ou on ne la trace pas.
    \end{enumerate}
    \textbf{Formule à retenir :} $\text{Aire}_i = a_i \times d_i = n_i$. $\sum \text{Aire}_i = N$ (ou 1 pour les fréquences).
\end{methode}

\begin{exoresolu}[Histogramme des consommations électriques]
    \textbf{Énoncé :} Consommation mensuelle (en kWh) de 80 abonnés :
    \begin{center}
    \begin{tabular}{|c|c|c|c|c|c|}\hline
    \rowcolor{popblueL} Classes (kWh) & $[0;50[$ & $[50;100[$ & $[100;200[$ & $[200;300[$ & $[300;500[$ \\ \hline
    Effectifs & 10 & 20 & 30 & 15 & 5 \\ \hline
    \end{tabular}
    \end{center}
    1. Calculer les densités. 2. Tracer l'histogramme (schéma coté). 3. Quelle proportion d'abonnés consomme entre 100 et 200 kWh ?
    \tcblower
    \textbf{Solution détaillée :}
    \begin{enumerate}
        \item \textbf{Amplitudes $a_i$ et Densités $d_i = n_i / a_i$ :}
        \begin{align*}
            [0;50[: & a=50, n=10 \Rightarrow d = 10/50 = 0,20 \\
            [50;100[: & a=50, n=20 \Rightarrow d = 20/50 = 0,40 \\
            [100;200[: & a=100, n=30 \Rightarrow d = 30/100 = 0,30 \\
            [200;300[: & a=100, n=15 \Rightarrow d = 15/100 = 0,15 \\
            [300;500[: & a=200, n=5 \Rightarrow d = 5/200 = 0,025
        \end{align*}
        \item \textbf{Vérification des aires :} $50\times0,20=10$, $50\times0,40=20$, $100\times0,30=30$, $100\times0,15=15$, $200\times0,025=5$. Somme $=80=N$. OK.
        \item \textbf{Schéma de l'histogramme (côté) :}
        \begin{popfigure}
        \begin{tikzpicture}[scale=0.7, xscale=0.02, yscale=2]
        % Axes
        \draw[->, thick] (-10,0) -- (510,0) node[right] {$x$ (kWh)};
        \draw[->, thick] (0,-0.05) -- (0,0.45) node[above] {$d$ (abonnés/kWh)};
        % Rupture origine
        \draw[thick] (-5,-0.02) -- (-5,0.02);
        \draw[thick] (5,-0.02) -- (5,0.02);
        % Rectangles
        \fill[popblue!50, draw=popblue, thick] (0,0) rectangle (50,0.20);
        \fill[popblue!50, draw=popblue, thick] (50,0) rectangle (100,0.40);
        \fill[popblue!50, draw=popblue, thick] (100,0) rectangle (200,0.30);
        \fill[popblue!50, draw=popblue, thick] (200,0) rectangle (300,0.15);
        \fill[popblue!50, draw=popblue, thick] (300,0) rectangle (500,0.025);
        % Graduations
        \foreach \x/\label in {0/0, 50/50, 100/100, 200/200, 300/300, 500/500}
            \draw (\x, -0.01) -- (\x, 0.01) node[below, font=\small] {\label};
        \foreach \y/\label in {0.1/0,1, 0.2/0,2, 0.3/0,3, 0.4/0,4}
            \draw (-5, \y) -- (5, \y) node[left, font=\small] {\label};
        % Densités écrites sur barres
        \node at (25,0.22) {\small 0,20};
        \node at (75,0.42) {\small 0,40};
        \node at (150,0.32) {\small 0,30};
        \node at (250,0.17) {\small 0,15};
        \node at (400,0.045) {\small 0,025};
        \end{tikzpicture}
        \legende{Histogramme des consommations (densités en abonnés/kWh).}
        \end{popfigure}
        \item \textbf{Proportion $[100;200[$ :} Effectif $=30$. Fréquence $=30/80 = 0,375 = 37,5\%$.
    \end{enumerate}
    \textbf{Conclusion :} L'histogramme montre une concentration autour de 100-200 kWh. La classe $[300;500[$ est large mais peu dense (queue de distribution). 37,5\% des abonnés consomment entre 100 et 200 kWh.
\end{exoresolu}

\courssub{Courbes cumulées (polygones) et lecture graphique des quartiles}

\begin{methode}[Tracer et exploiter les courbes cumulées croissantes]
    \textbf{Principe :} Représenter les points $(x_i; F_i)$ où $x_i$ = borne \emph{supérieure} de la classe $i$ (pour croissante) et $F_i$ = fréquence cumulée croissante (en \%). Relier par segments (polygone) ou courbe lissée.
    \begin{enumerate}
        \item \textbf{Points à tracer (croissante) :} $(b_1; F_1), (b_2; F_2), \ldots, (b_k; 100)$. Ajouter $(b_0; 0)$ pour partir de l'axe.
        \item \textbf{Axes :} Abscisses = variable $X$ (bornes de classes). Ordonnées = fréquences cumulées en \% (0 à 100) ou proportions (0 à 1).
        \item \textbf{Lecture graphique (quartiles, médiane, centiles) :}
        \begin{itemize}
            \item Médiane $M_e$ (50\%) : Abscisse du point d'ordonnée 50 sur la courbe.
            \item Premier quartile $Q_1$ (25\%) : Abscisse du point d'ordonnée 25.
            \item Troisième quartile $Q_3$ (75\%) : Abscisse du point d'ordonnée 75.
            \item Écart interquartile $EIQ = Q_3 - Q_1$.
        \end{itemize}
        \item \textbf{Courbe décroissante :} Points $(a_i; F'_i)$ où $a_i$ = borne \emph{inférieure}. $F'_1=100$ à $x_{\min}$, $F'_k=0$ à $x_{\max}$ (ou dernière borne). Utile pour lire « nombre de valeurs $\ge$ seuil ».
    \end{enumerate}
    \textbf{Précision Probatoire :} La lecture graphique donne une \emph{estimation}. Pour le calcul exact (formule), voir méthode suivante. Au Probatoire, on demande souvent : « Déterminer graphiquement... » $\rightarrow$ tracer des pointillés, lire sur l'axe, encadrer la valeur.
\end{methode}

\begin{exoresolu}[Courbe cumulée et quartiles : résistance à la traction]
    \textbf{Énoncé :} On teste la résistance $R$ (en MPa) de 100 éprouvettes.
    \begin{center}
    \begin{tabular}{|c|c|c|c|c|c|c|}\hline
    \rowcolor{popblueL} Classes (MPa) & $[300;320[$ & $[320;340[$ & $[340;360[$ & $[360;380[$ & $[380;400[$ & $[400;420[$ \\ \hline
    Effectifs & 5 & 15 & 30 & 30 & 15 & 5 \\ \hline
    \end{tabular}
    \end{center}
    1. Construire le tableau des fréquences cumulées croissantes. 2. Tracer la courbe cumulée croissante (schéma). 3. Déterminer graphiquement la médiane, $Q_1$, $Q_3$ et l'écart interquartile.
    \tcblower
    \textbf{Solution détaillée :}
    \begin{enumerate}
        \item \textbf{Tableau cumulé ($N=100 \Rightarrow$ effectifs = fréquences en \%) :}
        \begin{center}
        \begin{tabular}{|c|c|c|c|}\hline
        \rowcolor{popblueL} Classes & $n_i$ & $N_i$ & $F_i$(\%) \\ \hline
        $[300;320[$ & 5 & 5 & 5 \\ \hline
        $[320;340[$ & 15 & 20 & 20 \\ \hline
        $[340;360[$ & 30 & 50 & 50 \\ \hline
        $[360;380[$ & 30 & 80 & 80 \\ \hline
        $[380;400[$ & 15 & 95 & 95 \\ \hline
        $[400;420[$ & 5 & 100 & 100 \\ \hline
        \end{tabular}
        \end{center}
        \item \textbf{Points à tracer (bornes sup, $F_i$) :} $(300;0), (320;5), (340;20), (360;50), (380;80), (400;95), (420;100)$.
        \item \textbf{Schéma de la courbe cumulée :}
        \begin{popfigure}
        \begin{tikzpicture}[scale=0.8, xscale=0.02, yscale=0.06]
        % Axes
        \draw[->, thick] (290,0) -- (430,0) node[right] {$R$ (MPa)};
        \draw[->, thick] (290,0) -- (290,105) node[above] {$F$ (\%)};
        % Grille légère
        \foreach \x in {300,320,...,420} \draw[popline] (\x,0) -- (\x,100);
        \foreach \y in {0,25,50,75,100} \draw[popline] (290,\y) -- (420,\y);
        % Courbe
        \draw[poporange, very thick] (300,0) -- (320,5) -- (340,20) -- (360,50) -- (380,80) -- (400,95) -- (420,100);
        % Points
        \foreach \p in {(300,0),(320,5),(340,20),(360,50),(380,80),(400,95),(420,100)}
            \fill[poporange] \p circle (2pt);
        % Lectures graphiques (pointillés)
        \draw[dashed, popgreen, thick] (290,25) -- (335,25) -- (335,0) node[below, font=\small, popgreen] {$Q_1 \approx 335$};
        \draw[dashed, popgreen, thick] (290,50) -- (360,50) -- (360,0) node[below, font=\small, popgreen] {$M_e \approx 360$};
        \draw[dashed, popgreen, thick] (290,75) -- (385,75) -- (385,0) node[below, font=\small, popgreen] {$Q_3 \approx 385$};
        \end{tikzpicture}
        \legende{Courbe cumulée croissante et lecture des quartiles.}
        \end{popfigure}
        \item \textbf{Lectures (interpolation linéaire sur les segments pour valeurs exactes) :}
        \begin{itemize}
            \item $Q_1 (25\%)$ : Sur segment $(320;20) \to (340;50)$. Pente $=30/20=1,5$ \%/MPa. $25-20=5\% \Rightarrow \Delta x = 5/1,5 \approx 3,33$. $Q_1 \approx 323,3$ MPa.
            \item $M_e (50\%)$ : Point $(360;50)$ exact. $M_e = 360$ MPa.
            \item $Q_3 (75\%)$ : Sur segment $(360;50) \to (380;80)$. Pente $=30/20=1,5$. $75-50=25\% \Rightarrow \Delta x = 25/1,5 \approx 16,67$. $Q_3 \approx 376,7$ MPa.
        \end{itemize}
        \item \textbf{Écart interquartile :} $EIQ = Q_3 - Q_1 \approx 376,7 - 323,3 = 53,4$ MPa.
    \end{enumerate}
    \textbf{Conclusion :} La médiane est 360 MPa. 50\% des éprouvettes ont une résistance entre 323 et 377 MPa environ. La courbe cumulée permet une estimation rapide ; le calcul par interpolation (formule) est plus précis.
\end{exoresolu}

\courssub{Étude complète : position, dispersion, interprétation technique}

\begin{methode}[Étude complète d'une série groupée en classes (Moyenne, Médiane, Quartiles, Écart-type)]
    \textbf{Données :} Classes $[x_{inf}; x_{sup}[$, effectifs $n_i$, total $N$. Milieux $x_i = \frac{x_{inf}+x_{sup}}{2}$. Amplitudes $h_i = x_{sup} - x_{inf}$.
    \begin{enumerate}
        \item \textbf{Moyenne $\bar{x}$ :} $\bar{x} = \frac{1}{N}\sum n_i x_i$. (Unité = unité de $X$).
        \item \textbf{Variance $s^2$ et Écart-type $s$ :}
        $s^2 = \frac{1}{N}\sum n_i (x_i - \bar{x})^2 = \frac{\sum n_i x_i^2}{N} - \bar{x}^2$ (formule de König-Huygens).
        $s = \sqrt{s^2}$. (Unité = unité de $X$).
        \item \textbf{Médiane $M_e$ (calcul exact par interpolation linéaire) :}
        Trouver la classe médiane $j$ tq $N_{j-1} < N/2 \le N_j$.
        $M_e = x_{inf,j} + \frac{\frac{N}{2} - N_{j-1}}{n_j} \times h_j$.
        \item \textbf{Quartiles $Q_1, Q_3$ (calcul exact) :}
        Classe de $Q_1$ : $N_{j-1} < N/4 \le N_j$. $Q_1 = x_{inf,j} + \frac{N/4 - N_{j-1}}{n_j} \times h_j$.
        Classe de $Q_3$ : $N_{j-1} < 3N/4 \le N_j$. $Q_3 = x_{inf,j} + \frac{3N/4 - N_{j-1}}{n_j} \times h_j$.
        \item \textbf{Coefficient de variation $CV = \frac{s}{\bar{x}} \times 100$ (en \%) :} Permet de comparer la dispersion relative. $CV < 15\%$ : série homogène / peu dispersée.
        \item \textbf{Interprétation technique :} Comparer $\bar{x}$ à la cible (nominale). $s$ ou $EIQ$ $\rightarrow$ régularité du processus. $CV$ $\rightarrow$ capacité du processus.
    \end{enumerate}
    \textbf{Formules à connaître par cœur (Probatoire) :} $\bar{x}$, $s^2$ (König-Huygens), $M_e$, $Q_1, Q_3$ (interpolation linéaire).
\end{methode}

\begin{exoresolu}[Étude complète : diamètre d'axes moteurs]
    \textbf{Énoncé :} Contrôle du diamètre $D$ (mm) de 200 axes. Spécification : $25,00 \pm 0,10$ mm.
    \begin{center}
    \begin{tabular}{|c|c|c|c|c|c|c|c|}\hline
    \rowcolor{popblueL} Classes (mm) & $[24,80;24,85[$ & $[24,85;24,90[$ & $[24,90;24,95[$ & $[24,95;25,00[$ & $[25,00;25,05[$ & $[25,05;25,10[$ & $[25,10;25,15[$ \\ \hline
    Effectifs & 5 & 15 & 30 & 50 & 60 & 30 & 10 \\ \hline
    \end{tabular}
    \end{center}
    1. Calculer la moyenne $\bar{D}$ et l'écart-type $s$.
    2. Calculer la médiane $M_e$ et les quartiles $Q_1, Q_3$.
    3. Calculer le coefficient de variation $CV$. La série est-elle homogène ?
    4. Quel pourcentage d'axes est hors tolérance ($<24,90$ ou $>25,10$) ? (Utiliser les effectifs cumulés).
    \tcblower
    \textbf{Solution détaillée :}
    \begin{enumerate}
        \item \textbf{Tableau de calcul (milieux $x_i$, $n_i x_i$, $n_i x_i^2$) :} Amplitude constante $h=0,05$.
        \begin{center}
        \begin{tabular}{|c|c|c|c|c|c|}\hline
        \rowcolor{popblueL} Classes & $x_i$ & $n_i$ & $n_i x_i$ & $n_i x_i^2$ & $N_i$ \\ \hline
        $[24,80;24,85[$ & 24,825 & 5 & 124,125 & 3081,38 & 5 \\ \hline
        $[24,85;24,90[$ & 24,875 & 15 & 373,125 & 9281,48 & 20 \\ \hline
        $[24,90;24,95[$ & 24,925 & 30 & 747,750 & 18637,67 & 50 \\ \hline
        $[24,95;25,00[$ & 24,975 & 50 & 1248,750 & 31187,53 & 100 \\ \hline
        $[25,00;25,05[$ & 25,025 & 60 & 1501,500 & 37575,04 & 160 \\ \hline
        $[25,05;25,10[$ & 25,075 & 30 & 752,250 & 18862,67 & 190 \\ \hline
        $[25,10;25,15[$ & 25,125 & 10 & 251,250 & 6312,66 & 200 \\ \hline \hline
        \textbf{Total} & & \textbf{200} & \textbf{4998,75} & \textbf{124938,43} & \\ \hline
        \end{tabular}
        \end{center}
        \textbf{Moyenne :} $\bar{D} = \frac{4998,75}{200} = 24,99375 \approx 24,994$ mm.
        \textbf{Variance (König-Huygens) :} 
        $\frac{\sum n_i x_i^2}{N} = \frac{124938,43}{200} = 624,69215$.
        $\bar{x}^2 = (24,99375)^2 = 624,68754$.
        $s^2 = 624,69215 - 624,68754 = 0,00461$.
        $s = \sqrt{0,00461} \approx 0,0679$ mm.
        \item \textbf{Médiane ($N/2=100$) :} Classe médiane = $[24,95;25,00[$ ($N_3=50 < 100 \le N_4=100$).
        $M_e = 24,95 + \frac{100 - 50}{50} \times 0,05 = 24,95 + 0,05 = 25,000$ mm.
        \textbf{$Q_1$ ($N/4=50$) :} Classe $[24,90;24,95[$ ($N_2=20 < 50 \le N_3=50$). Borne inf $x_{inf}=24,90$.
        $Q_1 = 24,90 + \frac{50 - 20}{30} \times 0,05 = 24,90 + 0,05 = 24,950$ mm.
        \textbf{$Q_3$ ($3N/4=150$) :} Classe $[25,00;25,05[$ ($N_4=100 < 150 \le N_5=160$). Borne inf $x_{inf}=25,00$.
        $Q_3 = 25,00 + \frac{150 - 100}{60} \times 0,05 = 25,00 + 0,04167 = 25,0417$ mm.
        \item \textbf{Coefficient de variation :} $CV = \frac{0,0679}{24,994} \times 100 \approx 0,27\%$.
        $CV < 15\% \Rightarrow$ Série \textbf{très homogène} (processus très stable).
        \item \textbf{Hors tolérance $[24,90; 25,10]$ :}
        \begin{itemize}
            \item $< 24,90$ : Classes $[24,80;24,85[$ et $[24,85;24,90[$. Effectifs $= 5+15=20$. Fréquence $= 10\%$.
            \item $> 25,10$ : Classe $[25,10;25,15[$. Effectif $= 10$. Fréquence $= 5\%$.
            \item Total hors tolérance $= 20+10=30$ axes. Pourcentage $= 15\%$.
        \end{itemize}
    \end{enumerate}
    \textbf{Conclusion :} Le processus est centré sur la cible (Moyenne $\approx 24,994 \approx 25,00$, Médiane $=25,00$) et très précis ($s \approx 0,068$ mm, $CV=0,27\%$). Cependant, 15\% de la production est hors tolérance (principalement côté bas), ce qui indique un décalage léger ou une dispersion non négligeable par rapport à l'intervalle de tolérance $\pm 0,10$ mm ($6s \approx 0,41$ mm $> 0,20$ mm). Le processus n'est pas capable ($C_p = \frac{IT}{6s} \approx \frac{0,20}{0,41} < 1$).
\end{exoresolu}

\courssub{Méthode : Lecture et exploitation d'un histogramme ou courbe cumulée fournie}

\begin{methode}[Analyser un document graphique (histogramme ou courbe) sans tableau de données]
    \textbf{Situation fréquente au Probatoire :} Le tableau n'est pas donné, seul le graphique l'est.
    \begin{enumerate}
        \item \textbf{Histogramme :} Lire les bornes des classes sur l'axe des abscisses. Lire les hauteurs (densités) sur l'axe des ordonnées. \textbf{Calculer les effectifs :} $n_i = \text{Aire}_i = \text{Largeur}_i \times \text{Hauteur}_i$. Reconstituer le tableau.
        \item \textbf{Courbe cumulée :} Lire les points caractéristiques (bornes de classes, valeurs aux pourcentages 25, 50, 75). Relever les coordonnées $(x; F)$.
        \item \textbf{Estimation de la moyenne :} Si histogramme, utiliser les milieux des classes et effectifs reconstitués. Si seule la courbe : estimation grossière par le centre de gravité visuel ou par la médiane si symétrique.
        \item \textbf{Comparaison de deux séries :} Superposition de deux histogrammes (transparence) ou deux courbes cumulées. Comparer positions (médianes/moyennes) et dispersions (largeur à 50\%, écart interquartile).
    \end{enumerate}
    \textbf{Réflexe :} Toujours vérifier l'échelle des ordonnées (effectifs ? fréquences ? densités ? \%). Une hauteur de 10 peut signifier 10 pièces, 10\% ou 10 pièces/mm selon l'axe.
\end{methode}

\begin{exoresolu}[Reconstitution depuis un histogramme : dureté Brinell]
    \textbf{Énoncé :} L'histogramme ci-contre représente la dureté (HB) de 50 pièces traitées. L'axe des ordonnées est gradué en \textbf{densité} (effectifs par unité de dureté). Les classes ont une amplitude de 10 HB.
    \begin{popfigure}
    \begin{tikzpicture}[scale=0.7, xscale=0.05, yscale=1.5]
    \draw[->, thick] (140,0) -- (260,0) node[right] {Dureté (HB)};
    \draw[->, thick] (140,0) -- (140,2.2) node[above] {Densité};
    \foreach \x in {150,160,...,250} \draw (\x, -0.05) -- (\x, 0.05) node[below, font=\small] {\x};
    \foreach \y in {0.5,1.0,1.5,2.0} \draw (-5, \y) -- (5, \y) node[left, font=\small] {\y};
    % Barres : effectifs 5, 10, 15, 12, 8 -> densités 0.5, 1.0, 1.5, 1.2, 0.8
    \fill[popblue!50, draw=popblue, thick] (150,0) rectangle (160,0.5);
    \fill[popblue!50, draw=popblue, thick] (160,0) rectangle (170,1.0);
    \fill[popblue!50, draw=popblue, thick] (170,0) rectangle (180,1.5);
    \fill[popblue!50, draw=popblue, thick] (180,0) rectangle (190,1.2);
    \fill[popblue!50, draw=popblue, thick] (190,0) rectangle (200,0.8);
    % Légendes densités
    \node at (155,0.6) {\small 0,5}; \node at (165,1.1) {\small 1,0}; \node at (175,1.6) {\small 1,5}; \node at (185,1.3) {\small 1,2}; \node at (195,0.9) {\small 0,8};
    \end{tikzpicture}
    \legende{Histogramme des duretés (densités).}
    \end{popfigure}
    1. Reconstituer le tableau (classes, effectifs, fréquences \%). 2. Calculer la moyenne et l'écart-type. 3. Estimer la médiane par interpolation.
    \tcblower
    \textbf{Solution détaillée :}
    \begin{enumerate}
        \item \textbf{Reconstitution :} Amplitude $h=10$. $n_i = \text{Densité}_i \times 10$.
        \begin{center}
        \begin{tabular}{|c|c|c|c|c|}\hline
        \rowcolor{popblueL} Classes (HB) & Densité & $n_i$ & $f_i$(\%) & $x_i$ \\ \hline
        $[150;160[$ & 0,5 & 5 & 10 & 155 \\ \hline
        $[160;170[$ & 1,0 & 10 & 20 & 165 \\ \hline
        $[170;180[$ & 1,5 & 15 & 30 & 175 \\ \hline
        $[180;190[$ & 1,2 & 12 & 24 & 185 \\ \hline
        $[190;200[$ & 0,8 & 8 & 16 & 195 \\ \hline \hline
        \textbf{Total} & & \textbf{50} & \textbf{100} & \\ \hline
        \end{tabular}
        \end{center}
        \item \textbf{Moyenne :} $\sum n_i x_i = 5(155)+10(165)+15(175)+12(185)+8(195) = 775+1650+2625+2220+1560 = 8830$.
        $\bar{x} = 8830 / 50 = 176,6$ HB.
        \textbf{Variance :} $\sum n_i x_i^2 = 5(24025)+10(27225)+15(30625)+12(34225)+8(38025) = 120125+272250+459375+410700+304200 = 1566650$.
        $\frac{\sum n_i x_i^2}{N} = 31333$. $\bar{x}^2 = 31187,56$.
        $s^2 = 31333 - 31187,56 = 145,44$. $s \approx 12,06$ HB.
        \item \textbf{Médiane ($N/2=25$) :} $N_1=5, N_2=15, N_3=30$. Classe médiane = $[170;180[$ ($N_2=15 < 25 \le N_3=30$).
        $M_e = 170 + \frac{25 - 15}{15} \times 10 = 170 + \frac{10}{15}\times 10 = 170 + 6,67 = 176,67$ HB.
    \end{enumerate}
    \textbf{Conclusion :} Moyenne $\approx$ Médiane (176,6 vs 176,7) $\rightarrow$ distribution quasi symétrique. Dispersion modérée ($s \approx 12$ HB, $CV \approx 6,8\%$).
\end{exoresolu}

\begin{attention}[Les 5 erreurs qui coûtent des points au Probatoire]
    \begin{enumerate}
        \item \textbf{Confondre effectif et fréquence / densité :} Sur l'histogramme, l'axe des ordonnées est une \textbf{densité} (effectif / amplitude), pas l'effectif brut (sauf si amplitude = 1). Lire « 10 » sur l'axe ne veut pas dire 10 pièces si l'amplitude est 0,5 ! \textbf{Réflexe :} Vérifier l'unité de l'axe (effectifs, \%, densité). Calculer l'aire pour retrouver l'effectif.
        \item \textbf{Oublier l'interpolation linéaire pour médiane/quartiles (calcul exact) :} Écrire « La médiane est dans la classe $[a;b[$ » ne rapporte pas les points. Il faut appliquer la formule : $M_e = x_{inf} + \frac{N/2 - N_{ant}}{n_{cl}} \times h$. Même chose pour $Q_1, Q_3$. \textbf{Attention :} $N_{ant}$ = effectif cumulé \emph{avant} la classe, pas de la classe.
        \item \textbf{Erreur sur les bornes pour la courbe cumulée :} Courbe croissante $\rightarrow$ points aux \textbf{bornes supérieures} ($b_i$). Courbe décroissante $\rightarrow$ points aux \textbf{bornes inférieures} ($a_i$). Inverser les bornes décale toute la lecture graphique.
        \item \textbf{Mauvaise gestion des classes d'amplitude inégale :} Tracer des rectangles de hauteur = effectif (au lieu de densité) $\rightarrow$ histogramme faux visuellement (aires non proportionnelles). Au Probatoire, si amplitudes inégales, \textbf{calculer les densités} impérativement.
        \item \textbf{Arrondis prématurés et perte de précision :} Arrondir la moyenne $\bar{x}$ à 2 décimales \emph{avant} de calculer $s^2 = \frac{\sum n_i x_i^2}{N} - \bar{x}^2$ crée une erreur énorme sur la variance (différence de deux grands nombres proches). \textbf{Règle :} Garder toutes les décimales en mémoire (calculatrice) ou utiliser la formule $\frac{1}{N}\sum n_i (x_i - \bar{x})^2$ avec $\bar{x}$ exact (fraction) ou très précis. Ne donner l'arrondi qu'au résultat final.
    \end{enumerate}
\end{attention}

\newpage

\coursec{Exercices}

\courssub{Je vérifie mes connaissances}

\begin{exercice}[QCM : vocabulaire des classes]
Pour chaque question, une seule réponse est exacte. Recopier la lettre de la bonne réponse.
\begin{enumerate}
\item Une série statistique est regroupée en classes $[0 ; 10[$, $[10 ; 20[$, $[20 ; 30[$. L'amplitude de chaque classe est :
\begin{enumerate}
\item[A.] 5 \quad B. 10 \quad C. 20 \quad D. 30
\end{enumerate}
\item Le centre de la classe $[10 ; 20[$ est :
\begin{enumerate}
\item[A.] 10 \quad B. 12 \quad C. 15 \quad D. 20
\end{enumerate}
\item L'effectif cumulé croissant de la classe $[10 ; 20[$ est :
\begin{enumerate}
\item[A.] l'effectif de la classe $[10 ; 20[$ seule ;
\item[B.] la somme des effectifs des classes $[0 ; 10[$ et $[10 ; 20[$ ;
\item[C.] la somme des effectifs de toutes les classes ;
\item[D.] l'effectif de la dernière classe.
\end{enumerate}
\item Dans un histogramme d'une série groupée en classes de même amplitude, la hauteur de chaque rectangle est proportionnelle :
\begin{enumerate}
\item[A.] au centre de la classe ;
\item[B.] à l'amplitude de la classe ;
\item[C.] à l'effectif de la classe ;
\item[D.] à l'effectif total.
\end{enumerate}
\end{enumerate}
\end{exercice}

\begin{exercice}[Vrai ou faux, en justifiant]
Dire si chaque affirmation est vraie ou fausse, puis justifier la réponse par une phrase ou un calcul.
\begin{enumerate}
\item Le dernier effectif cumulé croissant est toujours égal à l'effectif total.
\item Le premier effectif cumulé décroissant est toujours égal à l'effectif de la première classe.
\item La somme des fréquences de toutes les classes est égale à $1$ (ou à $100\,\%$).
\item La fréquence cumulée croissante peut dépasser $100\,\%$.
\end{enumerate}
\end{exercice}

\begin{exercice}[Définitions]
\begin{enumerate}
\item Définir : classe modale, amplitude d'une classe, centre d'une classe.
\item Expliquer la différence entre un effectif cumulé croissant et un effectif cumulé décroissant.
\item Qu'appelle-t-on courbe cumulative croissante ? Quelles valeurs place-t-on en abscisse ?
\end{enumerate}
\end{exercice}

\begin{exercice}[Calculs directs]
Une enquête porte sur le temps de trajet (en minutes) de $50$ élèves d'un lycée technique de Yaoundé. On obtient :
\begin{center}
\begin{tabular}{|l|c|c|c|c|}
\hline
Temps (min) & $[0 ; 10[$ & $[10 ; 20[$ & $[20 ; 30[$ & $[30 ; 40[$ \\
\hline
Effectif & 8 & 17 & 15 & 10 \\
\hline
\end{tabular}
\end{center}
\begin{enumerate}
\item Vérifier que l'effectif total est bien $50$.
\item Calculer la fréquence (en $\%$) de chaque classe.
\item Calculer les effectifs cumulés croissants.
\item Calculer les effectifs cumulés décroissants.
\item Quelle est la classe modale ?
\end{enumerate}
\end{exercice}

\courssub{Je m'entraîne}

\begin{exercice}[Compléter un tableau]
Un atelier de menuiserie de Bafoussam mesure la longueur (en cm) de $60$ planches :
\begin{center}
\begin{tabular}{|l|c|c|c|c|}
\hline
Longueur (cm) & $[140 ; 150[$ & $[150 ; 160[$ & $[160 ; 170[$ & $[170 ; 180[$ \\
\hline
Effectif & 9 & 21 & 18 & 12 \\
\hline
Fréquence (\%) & & & & \\
\hline
ECC & & & & \\
\hline
ECD & & & & \\
\hline
\end{tabular}
\end{center}
\begin{enumerate}
\item Recopier et compléter le tableau (fréquences en $\%$, arrondies à $0,1\,\%$, ECC et ECD).
\item Déterminer la classe modale.
\item Combien de planches mesurent moins de \SI{160}{cm} ?
\item Combien de planches mesurent au moins \SI{160}{cm} ?
\end{enumerate}
\end{exercice}

\begin{exercice}[Retrouver les effectifs]
Le tableau ci-dessous donne les effectifs cumulés croissants des masses (en kg) de sacs de ciment contrôlés sur un chantier de Douala :
\begin{center}
\begin{tabular}{|l|c|c|c|c|c|}
\hline
Masse (kg) & $[45 ; 47[$ & $[47 ; 49[$ & $[49 ; 51[$ & $[51 ; 53[$ & $[53 ; 55[$ \\
\hline
ECC & 6 & 20 & 52 & 68 & 75 \\
\hline
\end{tabular}
\end{center}
\begin{enumerate}
\item Quel est l'effectif total de la série ?
\item Retrouver l'effectif de chaque classe.
\item Calculer les fréquences (en $\%$) de chaque classe.
\item Calculer les effectifs cumulés décroissants.
\end{enumerate}
\end{exercice}

\begin{exercice}[Construire un histogramme]
Les notes sur $20$ obtenues par $40$ élèves de Première technique à un devoir de mathématiques sont regroupées ainsi :
\begin{center}
\begin{tabular}{|l|c|c|c|c|c|}
\hline
Note & $[0 ; 4[$ & $[4 ; 8[$ & $[8 ; 12[$ & $[12 ; 16[$ & $[16 ; 20]$ \\
\hline
Effectif & 3 & 9 & 16 & 9 & 3 \\
\hline
\end{tabular}
\end{center}
\begin{enumerate}
\item Construire l'histogramme de cette série (on prendra \SI{1}{cm} pour $4$ points en abscisse et \SI{1}{cm} pour $2$ élèves en ordonnée).
\item Déterminer la classe modale.
\item Calculer les fréquences cumulées croissantes (en $\%$).
\item Quel pourcentage d'élèves a obtenu une note inférieure à $12$ ?
\end{enumerate}
\end{exercice}

\begin{exercice}[Courbe cumulative croissante]
Une coopérative agricole de l'Ouest a pesé $100$ sacs de maïs. Les masses (en kg) sont regroupées en classes :
\begin{center}
\begin{tabular}{|l|c|c|c|c|c|}
\hline
Masse (kg) & $[40 ; 45[$ & $[45 ; 50[$ & $[50 ; 55[$ & $[55 ; 60[$ & $[60 ; 65[$ \\
\hline
Effectif & 10 & 22 & 35 & 23 & 10 \\
\hline
\end{tabular}
\end{center}
\begin{enumerate}
\item Dresser le tableau des effectifs cumulés croissants.
\item Tracer la courbe cumulative croissante (polygone des ECC) : on placera en abscisse les bornes supérieures des classes.
\item Estimer graphiquement la masse médiane de la série.
\item Estimer graphiquement le nombre de sacs dont la masse est inférieure à \SI{52}{kg}.
\end{enumerate}
\end{exercice}

\courssub{Je me prépare au Probatoire}

\begin{exercice}[Les plants de cacao — type Probatoire]
Un centre de formation agricole de la région du Centre mesure la taille (en cm) de $80$ plants de cacao six mois après la mise en pépinière. Les résultats sont regroupés dans le tableau incomplet ci-dessous :
\begin{center}
\begin{tabular}{|l|c|c|c|c|c|}
\hline
Taille (cm) & $[30 ; 40[$ & $[40 ; 50[$ & $[50 ; 60[$ & $[60 ; 70[$ & $[70 ; 80[$ \\
\hline
Effectif $n_i$ & 8 & 16 & 28 & & 8 \\
\hline
Fréquence $f_i$ (\%) & & & & & \\
\hline
ECC & & & & & \\
\hline
ECD & & & & & \\
\hline
\end{tabular}
\end{center}
\begin{enumerate}
\item Sachant que l'effectif total est $80$, calculer l'effectif de la classe $[60 ; 70[$.
\item Recopier et compléter le tableau : fréquences en $\%$ (arrondies à l'unité), effectifs cumulés croissants (ECC) et effectifs cumulés décroissants (ECD).
\item Déterminer la classe modale de cette série. Interpréter le résultat par une phrase.
\item Construire l'histogramme de la série (échelle : \SI{1}{cm} pour \SI{10}{cm} de taille en abscisse ; \SI{1}{cm} pour $4$ plants en ordonnée).
\item \begin{enumerate}
\item Combien de plants mesurent moins de \SI{60}{cm} ?
\item Combien de plants mesurent au moins \SI{60}{cm} ?
\item Le technicien déclare : « plus de la moitié des plants mesurent moins de \SI{60}{cm} ». A-t-il raison ? Justifier.
\end{enumerate}
\item Tracer la courbe cumulative croissante et estimer graphiquement la taille médiane des plants.
\end{enumerate}
\end{exercice}

\begin{exercice}[Les salaires d'une PME de Douala — type Probatoire]
Le tableau ci-dessous donne la répartition des salaires mensuels (en milliers de FCFA) des $60$ employés d'une PME de Douala :
\begin{center}
\begin{tabular}{|l|c|c|c|c|c|}
\hline
Salaire (milliers FCFA) & $[50 ; 75[$ & $[75 ; 100[$ & $[100 ; 125[$ & $[125 ; 150[$ & $[150 ; 175[$ \\
\hline
Effectif & 12 & 18 & 15 & 9 & 6 \\
\hline
\end{tabular}
\end{center}
\begin{enumerate}
\item Calculer les fréquences (en $\%$) de chaque classe, arrondies à $0,1\,\%$.
\item Dresser le tableau des effectifs cumulés croissants et des effectifs cumulés décroissants.
\item Construire l'histogramme de cette série.
\item Sur un même graphique, tracer la courbe cumulative croissante et la courbe cumulative décroissante (on placera en abscisse les bornes des classes).
\item \begin{enumerate}
\item Estimer graphiquement le salaire médian de l'entreprise (abscisse du point d'intersection des deux courbes).
\item Estimer graphiquement le pourcentage d'employés gagnant moins de \num{100000} FCFA. Vérifier le résultat par le calcul à l'aide du tableau.
\end{enumerate}
\item Le directeur affirme : « au moins $75\,\%$ des employés gagnent moins de \num{150000} FCFA ». Cette affirmation est-elle exacte ? Justifier la réponse.
\end{enumerate}
\end{exercice}

\begin{exercice}[Comparaison de deux écoles et rapport — type Probatoire]
Deux lycées techniques A et B ont présenté leurs élèves à un même concours d'entrée dans une école de formation technique. Les notes sur $20$ sont regroupées dans les tableaux suivants.

Lycée A ($100$ candidats) :
\begin{center}
\begin{tabular}{|l|c|c|c|c|c|}
\hline
Note & $[0 ; 5[$ & $[5 ; 10[$ & $[10 ; 15[$ & $[15 ; 20]$ \\
\hline
Effectif & 10 & 30 & 45 & 15 \\
\hline
\end{tabular}
\end{center}

Lycée B ($80$ candidats) :
\begin{center}
\begin{tabular}{|l|c|c|c|c|c|}
\hline
Note & $[0 ; 5[$ & $[5 ; 10[$ & $[10 ; 15[$ & $[15 ; 20]$ \\
\hline
Effectif & 4 & 20 & 40 & 16 \\
\hline
\end{tabular}
\end{center}
\begin{enumerate}
\item Pour chaque lycée, calculer les fréquences (en $\%$) de chaque classe.
\item Pour chaque lycée, calculer les fréquences cumulées croissantes (en $\%$).
\item Pourquoi est-il nécessaire de comparer les fréquences plutôt que les effectifs ? Justifier.
\item \begin{enumerate}
\item Quel pourcentage de candidats a obtenu au moins $10/20$ dans le lycée A ? dans le lycée B ?
\item Quel pourcentage de candidats a obtenu au moins $15/20$ dans chaque lycée ?
\end{enumerate}
\item Construire sur deux graphiques distincts les histogrammes des fréquences des deux lycées.
\item En s'appuyant sur les résultats précédents, rédiger un court rapport (cinq à huit lignes) à l'intention du jury du concours, indiquant quel lycée a globalement obtenu les meilleures performances, avec des arguments chiffrés.
\end{enumerate}
\end{exercice}

\newpage

\coursec{Corrigés des exercices}

\coursec{Séries statistiques regroupées en classes — Exercices corrigés}

\courssub{Exercices d'application directe — Vocabulaire et calculs}

\begin{corrige}[Exercice 1 — QCM : vocabulaire des classes]
\begin{enumerate}
\item \textbf{Réponse B.} L'amplitude d'une classe $[a ; b[$ est $b-a$. Ici : $10-0=20-10=30-20=10$.
\item \textbf{Réponse C.} Le centre de $[10 ; 20[$ est $\dfrac{10+20}{2}=15$.
\item \textbf{Réponse B.} L'effectif cumulé croissant d'une classe est la somme des effectifs de cette classe et de toutes les classes précédentes.
\item \textbf{Réponse C.} Pour des classes de même amplitude, la hauteur de chaque rectangle est proportionnelle à l'effectif (ou à la fréquence) de la classe.
\end{enumerate}
\end{corrige}

\begin{corrige}[Exercice 2 — Vrai ou faux, en justifiant]
\begin{enumerate}
\item \textbf{Vrai.} Le dernier effectif cumulé croissant est la somme des effectifs de toutes les classes, c'est-à-dire l'effectif total $N$.
\item \textbf{Faux.} Le premier effectif cumulé \emph{décroissant} est la somme des effectifs de toutes les classes (on cumule « depuis la fin »), donc il est égal à l'effectif total $N$, et non à l'effectif de la première classe. C'est le premier effectif cumulé \emph{croissant} qui est égal à l'effectif de la première classe.
\item \textbf{Vrai.} La somme des fréquences vaut $\dfrac{n_1}{N}+\dfrac{n_2}{N}+\dots=\dfrac{n_1+n_2+\dots}{N}=\dfrac{N}{N}=1$, soit $100\,\%$.
\item \textbf{Faux.} La fréquence cumulée croissante est croissante et sa dernière valeur vaut exactement $100\,\%$ ; elle ne peut donc jamais dépasser $100\,\%$.
\end{enumerate}
\end{corrige}

\begin{corrige}[Exercice 3 — Définitions à connaître]
\begin{enumerate}
\item \begin{itemize}
\item \textbf{Classe modale} : classe dont l'effectif (ou la fréquence) est le plus grand.
\item \textbf{Amplitude d'une classe} $[a ; b[$ : la différence $b-a$.
\item \textbf{Centre d'une classe} $[a ; b[$ : le nombre $\dfrac{a+b}{2}$.
\end{itemize}
\item L'\textbf{effectif cumulé croissant} (ECC) d'une classe est la somme des effectifs de cette classe et de toutes les classes \emph{précédentes} : il répond à la question « combien de valeurs sont \emph{strictement inférieures} à la borne supérieure de la classe ? ». L'\textbf{effectif cumulé décroissant} (ECD) est la somme des effectifs de cette classe et de toutes les classes \emph{suivantes} : il répond à la question « combien de valeurs sont \emph{supérieures ou égales} à la borne inférieure de la classe ? ».
\item La \textbf{courbe cumulative croissante} (polygone des ECC) est la ligne brisée joignant les points dont l'ordonnée est l'effectif (ou la fréquence) cumulé croissant. En \textbf{abscisse}, on place les \textbf{bornes supérieures} des classes (avec, au départ, le point d'abscisse la borne inférieure de la première classe et d'ordonnée $0$).
\end{enumerate}
\end{corrige}

\begin{corrige}[Exercice 4 — Calculs directs sur un tableau complet]
\begin{enumerate}
\item $8+17+15+10=50$. L'effectif total est bien $N=50$.
\item Fréquences : $f_i=\dfrac{n_i}{N}\times 100$.
\[ \frac{8}{50}=16\,\% \quad ; \quad \frac{17}{50}=34\,\% \quad ; \quad \frac{15}{50}=30\,\% \quad ; \quad \frac{10}{50}=20\,\%. \]
Vérification : $16+34+30+20=100\,\%$.
\item ECC : $8$ ; $8+17=25$ ; $25+15=40$ ; $40+10=50$.
\item ECD (en partant de la dernière classe) : $10$ ; $10+15=25$ ; $25+17=42$ ; $42+8=50$.
\item Tableau récapitulatif :
\begin{center}
\begin{tabular}{|l|c|c|c|c|}
\hline
Temps (min) & $[0 ; 10[$ & $[10 ; 20[$ & $[20 ; 30[$ & $[30 ; 40[$ \\
\hline
Effectif & 8 & 17 & 15 & 10 \\
\hline
Fréquence (\%) & 16 & 34 & 30 & 20 \\
\hline
ECC & 8 & 25 & 40 & 50 \\
\hline
ECD & 50 & 42 & 25 & 10 \\
\hline
\end{tabular}
\end{center}
\item La classe modale est $[10 ; 20[$ (effectif maximal $17$).
\end{enumerate}
\end{corrige}

\begin{corrige}[Exercice 5 — Compléter un tableau à partir des effectifs]
\begin{enumerate}
\item Effectif total : $9+21+18+12=60$. Fréquences : $\dfrac{9}{60}=15\,\%$ ; $\dfrac{21}{60}=35\,\%$ ; $\dfrac{18}{60}=30\,\%$ ; $\dfrac{12}{60}=20\,\%$.
\begin{center}
\begin{tabular}{|l|c|c|c|c|}
\hline
Longueur (cm) & $[140 ; 150[$ & $[150 ; 160[$ & $[160 ; 170[$ & $[170 ; 180[$ \\
\hline
Effectif & 9 & 21 & 18 & 12 \\
\hline
Fréquence (\%) & 15 & 35 & 30 & 20 \\
\hline
ECC & 9 & 30 & 48 & 60 \\
\hline
ECD & 60 & 51 & 30 & 12 \\
\hline
\end{tabular}
\end{center}
\item La classe modale est $[150 ; 160[$ (effectif $21$, le plus grand).
\item « Moins de \SI{160}{cm} » correspond à l'ECC de la classe $[150 ; 160[$ : \textbf{30 planches}.
\item « Au moins \SI{160}{cm} » correspond à l'ECD de la classe $[160 ; 170[$ : \textbf{30 planches}. Vérification : $18+12=30$ et $60-30=30$.
\end{enumerate}
\end{corrige}

\begin{corrige}[Exercice 6 — Retrouver les effectifs à partir des ECC]
\begin{enumerate}
\item L'effectif total est le dernier ECC : $N=75$ sacs.
\item L'effectif de la première classe est le premier ECC : $n_1=6$. Pour les autres classes, on fait la différence de deux ECC consécutifs :
\begin{align*}
n_2 &= 20-6 = 14 ; & n_3 &= 52-20 = 32 ;\\
n_4 &= 68-52 = 16 ; & n_5 &= 75-68 = 7.
\end{align*}
Vérification : $6+14+32+16+7=75$.
\item Fréquences : $\dfrac{6}{75}=8\,\%$ ; $\dfrac{14}{75}\approx 18{,}7\,\%$ ; $\dfrac{32}{75}\approx 42{,}7\,\%$ ; $\dfrac{16}{75}\approx 21{,}3\,\%$ ; $\dfrac{7}{75}\approx 9{,}3\,\%$. (Somme : $100\,\%$.)
\item ECD : on cumule depuis la dernière classe : $7$ ; $7+16=23$ ; $23+32=55$ ; $55+14=69$ ; $69+6=75$.
\begin{center}
\begin{tabular}{|l|c|c|c|c|c|}
\hline
Masse (kg) & $[45 ; 47[$ & $[47 ; 49[$ & $[49 ; 51[$ & $[51 ; 53[$ & $[53 ; 55[$ \\
\hline
Effectif & 6 & 14 & 32 & 16 & 7 \\
\hline
Fréquence (\%) & 8 & 18{,}7 & 42{,}7 & 21{,}3 & 9{,}3 \\
\hline
ECC & 6 & 20 & 52 & 68 & 75 \\
\hline
ECD & 75 & 69 & 55 & 23 & 7 \\
\hline
\end{tabular}
\end{center}
\end{enumerate}
\end{corrige}

\courssub{Représentations graphiques — Histogrammes et courbes cumulatives}

\begin{corrige}[Exercice 7 — Construire un histogramme (classes de même amplitude)]
\begin{enumerate}
\item Les classes ont toutes la même amplitude ($4$ points), donc la hauteur de chaque rectangle est proportionnelle à l'effectif. Avec \SI{1}{cm} pour $4$ points en abscisse et \SI{1}{cm} pour $2$ élèves en ordonnée, chaque rectangle a une base de \SI{1}{cm} et des hauteurs respectives de $1{,}5$ ; $4{,}5$ ; $8$ ; $4{,}5$ et $1{,}5$ cm.
\begin{popfigure}
\begin{center}
\begin{tikzpicture}[x=0.25cm,y=0.5cm]
\draw[->,thick] (0,0) -- (21,0) node[below left]{\footnotesize Note};
\draw[->,thick] (0,0) -- (0,18) node[above left]{\footnotesize Effectif};
\foreach \x in {0,4,8,12,16,20} \draw (\x,0.2) -- (\x,-0.2) node[below]{\footnotesize \x};
\foreach \y in {2,4,...,16} \draw (0.2,\y) -- (-0.2,\y) node[left]{\footnotesize \y};
\draw[fill=popblueL,draw=popblue] (0,0) rectangle (4,3);
\draw[fill=popblueL,draw=popblue] (4,0) rectangle (8,9);
\draw[fill=poporangeL,draw=poporange] (8,0) rectangle (12,16);
\draw[fill=popblueL,draw=popblue] (12,0) rectangle (16,9);
\draw[fill=popblueL,draw=popblue] (16,0) rectangle (20,3);
\node at (2,1.5){\footnotesize 3}; \node at (6,4.5){\footnotesize 9}; \node at (10,8){\footnotesize 16};
\node at (14,4.5){\footnotesize 9}; \node at (18,1.5){\footnotesize 3};
\end{tikzpicture}
\end{center}
\legende{Histogramme des notes du devoir de mathématiques (échelle respectée).}
\end{popfigure}
\item La classe modale est $[8 ; 12[$ (effectif $16$).
\item Fréquences : $\dfrac{3}{40}=7{,}5\,\%$ ; $\dfrac{9}{40}=22{,}5\,\%$ ; $\dfrac{16}{40}=40\,\%$ ; $\dfrac{9}{40}=22{,}5\,\%$ ; $\dfrac{3}{40}=7{,}5\,\%$.\\
Fréquences cumulées croissantes : $7{,}5\,\%$ ; $30\,\%$ ; $70\,\%$ ; $92{,}5\,\%$ ; $100\,\%$.
\item « Note inférieure à $12$ » correspond à la fréquence cumulée croissante de la classe $[8 ; 12[$ : \textbf{$70\,\%$} des élèves, soit $28$ élèves ($3+9+16=28$).
\end{enumerate}
\end{corrige}

\begin{corrige}[Exercice 8 — Courbe cumulative croissante et lecture graphique]
\begin{enumerate}
\item Tableau des ECC :
\begin{center}
\begin{tabular}{|l|c|c|c|c|c|}
\hline
Masse (kg) & $[40 ; 45[$ & $[45 ; 50[$ & $[50 ; 55[$ & $[55 ; 60[$ & $[60 ; 65[$ \\
\hline
Effectif & 10 & 22 & 35 & 23 & 10 \\
\hline
ECC & 10 & 32 & 67 & 90 & 100 \\
\hline
\end{tabular}
\end{center}
\item On place en abscisse les bornes supérieures $45, 50, 55, 60, 65$ et en ordonnée les ECC, en partant du point $(40 ; 0)$. On joint les points $(40;0)$, $(45;10)$, $(50;32)$, $(55;67)$, $(60;90)$, $(65;100)$.
\begin{popfigure}
\begin{center}
\begin{tikzpicture}[x=0.32cm,y=0.05cm]
\draw[->,thick] (38,0) -- (67,0) node[below left]{\footnotesize Masse (kg)};
\draw[->,thick] (38,0) -- (38,108) node[above left]{\footnotesize ECC};
\foreach \x in {40,45,...,65} \draw (\x,1.5) -- (\x,-1.5) node[below]{\footnotesize \x};
\foreach \y in {10,20,...,100} \draw (38.4,\y) -- (37.6,\y) node[left]{\footnotesize \y};
\draw[very thick,popblue] (40,0) -- (45,10) -- (50,32) -- (55,67) -- (60,90) -- (65,100);
\foreach \p in {(40,0),(45,10),(50,32),(55,67),(60,90),(65,100)} \fill[popblue] \p circle (0.18);
\draw[dashed,poporange] (38,50) -- (52.3,50) -- (52.3,0);
\node[below,poporange] at (52.3,0){\footnotesize $\approx 52{,}3$};
\draw[dashed,popgreen] (52,0) -- (52,44) -- (38,44);
\node[left,popgreen] at (38,44){\footnotesize $\approx 44$};
\end{tikzpicture}
\end{center}
\legende{Courbe cumulative croissante des masses des sacs de maïs.}
\end{popfigure}
\item La médiane est l'abscisse du point de la courbe d'ordonnée $\dfrac{N}{2}=50$. Graphiquement, on lit $Me \approx \SI{52,3}{kg}$ (environ $\SI{52}{kg}$).
\item Pour une masse inférieure à \SI{52}{kg}, on lit l'ordonnée du point d'abscisse $52$ : environ \textbf{44 sacs}.
\end{enumerate}
\end{corrige}

\courssub{Sujets type Probatoire — Situations concrètes complètes}

\begin{corrige}[Exercice 9 — Les plants de cacao — Type Probatoire]
\textbf{Barème indicatif (sur 10)} : question 1 : 1 pt ; question 2 : 3 pts ; question 3 : 1 pt ; question 4 : 2 pts ; question 5 : 1,5 pt ; question 6 : 1,5 pt.
\begin{enumerate}
\item Effectif de la classe $[60 ; 70[$ :
\[ 80-(8+16+28+8)=80-60=20 \text{ plants}. \]
\item Fréquences : $\dfrac{8}{80}=10\,\%$ ; $\dfrac{16}{80}=20\,\%$ ; $\dfrac{28}{80}=35\,\%$ ; $\dfrac{20}{80}=25\,\%$ ; $\dfrac{8}{80}=10\,\%$.
\begin{center}
\begin{tabular}{|l|c|c|c|c|c|}
\hline
Taille (cm) & $[30 ; 40[$ & $[40 ; 50[$ & $[50 ; 60[$ & $[60 ; 70[$ & $[70 ; 80[$ \\
\hline
Effectif $n_i$ & 8 & 16 & 28 & 20 & 8 \\
\hline
Fréquence $f_i$ (\%) & 10 & 20 & 35 & 25 & 10 \\
\hline
ECC & 8 & 24 & 52 & 72 & 80 \\
\hline
ECD & 80 & 72 & 56 & 28 & 8 \\
\hline
\end{tabular}
\end{center}
\item La classe modale est $[50 ; 60[$ (effectif $28$). \emph{Interprétation : la plupart des plants de cacao ont, six mois après la mise en pépinière, une taille comprise entre \SI{50}{cm} et \SI{60}{cm}.}
\item Classes de même amplitude ($10$ cm) : hauteurs proportionnelles aux effectifs. Avec \SI{1}{cm} pour $4$ plants, les hauteurs sont : $2$ ; $4$ ; $7$ ; $5$ ; $2$ cm.
\begin{popfigure}
\begin{center}
\begin{tikzpicture}[x=0.5cm,y=0.25cm]
\draw[->,thick] (28,0) -- (82,0) node[below left]{\footnotesize Taille (cm)};
\draw[->,thick] (28,0) -- (28,30) node[above left]{\footnotesize Effectif};
\foreach \x in {30,40,...,80} \draw (\x,0.6) -- (\x,-0.6) node[below]{\footnotesize \x};
\foreach \y in {4,8,...,28} \draw (28.5,\y) -- (27.5,\y) node[left]{\footnotesize \y};
\draw[fill=popblueL,draw=popblue] (30,0) rectangle (40,8);
\draw[fill=popblueL,draw=popblue] (40,0) rectangle (50,16);
\draw[fill=poporangeL,draw=poporange] (50,0) rectangle (60,28);
\draw[fill=popblueL,draw=popblue] (60,0) rectangle (70,20);
\draw[fill=popblueL,draw=popblue] (70,0) rectangle (80,8);
\node at (35,4){\footnotesize 8}; \node at (45,8){\footnotesize 16}; \node at (55,14){\footnotesize 28};
\node at (65,10){\footnotesize 20}; \node at (75,4){\footnotesize 8};
\end{tikzpicture}
\end{center}
\legende{Histogramme des tailles des plants de cacao (échelle : 1 cm pour 4 plants en ordonnée).}
\end{popfigure}
\item \begin{enumerate}
\item « Moins de \SI{60}{cm} » : ECC de la classe $[50 ; 60[$, soit $8+16+28=\textbf{52 plants}$.
\item « Au moins \SI{60}{cm} » : ECD de la classe $[60 ; 70[$, soit $20+8=\textbf{28 plants}$.
\item La moitié de l'effectif total est $\dfrac{80}{2}=40$. Or $52 > 40$ : le technicien a \textbf{raison}, plus de la moitié des plants ($52$ sur $80$, soit $65\,\%$) mesurent moins de \SI{60}{cm}.
\end{enumerate}
\item On joint les points $(30;0)$, $(40;8)$, $(50;24)$, $(60;52)$, $(70;72)$, $(80;80)$. La médiane est l'abscisse du point d'ordonnée $40$ : on lit graphiquement $Me \approx \SI{56}{cm}$.
\begin{popfigure}
\begin{center}
\begin{tikzpicture}[x=0.1cm,y=0.055cm]
\draw[->,thick] (28,0) -- (83,0) node[below left]{\footnotesize Taille (cm)};
\draw[->,thick] (28,0) -- (28,86) node[above left]{\footnotesize ECC};
\foreach \x in {30,40,...,80} \draw (\x,1.2) -- (\x,-1.2) node[below]{\footnotesize \x};
\foreach \y in {10,20,...,80} \draw (28.6,\y) -- (27.4,\y) node[left]{\footnotesize \y};
\draw[very thick,popblue] (30,0) -- (40,8) -- (50,24) -- (60,52) -- (70,72) -- (80,80);
\foreach \p in {(30,0),(40,8),(50,24),(60,52),(70,72),(80,80)} \fill[popblue] \p circle (0.5);
\draw[dashed,poporange] (28,40) -- (55.8,40) -- (55.8,0);
\node[below,poporange] at (55.8,0){\footnotesize $\approx 56$};
\end{tikzpicture}
\end{center}
\legende{Courbe cumulative croissante et lecture de la médiane.}
\end{popfigure}
\end{enumerate}
\textbf{Points de vigilance} : justifier le calcul de l'effectif manquant par $80-60$ ; citer la définition de la classe modale ; pour la question 5.c), comparer explicitement $52$ à la moitié de l'effectif total ($40$) ; pour la médiane graphique, préciser qu'on lit l'abscisse du point d'ordonnée $\dfrac{N}{2}=40$ et tracer les pointillés de lecture.
\end{corrige}

\begin{corrige}[Exercice 10 — Les salaires d'une PME de Douala — Type Probatoire]
\textbf{Barème indicatif (sur 10)} : question 1 : 1,5 pt ; question 2 : 2 pts ; question 3 : 1,5 pt ; question 4 : 2 pts ; question 5 : 2 pts ; question 6 : 1 pt.
\begin{enumerate}
\item Effectif total : $12+18+15+9+6=60$. Fréquences :
\[ \frac{12}{60}=20{,}0\,\% \;;\; \frac{18}{60}=30{,}0\,\% \;;\; \frac{15}{60}=25{,}0\,\% \;;\; \frac{9}{60}=15{,}0\,\% \;;\; \frac{6}{60}=10{,}0\,\%. \]
\item Tableau des cumuls :
\begin{center}
\begin{tabular}{|l|c|c|c|c|c|}
\hline
Salaire (milliers FCFA) & $[50 ; 75[$ & $[75 ; 100[$ & $[100 ; 125[$ & $[125 ; 150[$ & $[150 ; 175[$ \\
\hline
Effectif & 12 & 18 & 15 & 9 & 6 \\
\hline
Fréquence (\%) & 20{,}0 & 30{,}0 & 25{,}0 & 15{,}0 & 10{,}0 \\
\hline
ECC & 12 & 30 & 45 & 54 & 60 \\
\hline
ECD & 60 & 48 & 30 & 15 & 6 \\
\hline
\end{tabular}
\end{center}
\item Classes de même amplitude ($25$ milliers) : hauteurs proportionnelles aux effectifs.
\begin{popfigure}
\begin{center}
\begin{tikzpicture}[x=0.045cm,y=0.22cm]
\draw[->,thick] (45,0) -- (180,0) node[below left]{\footnotesize Salaire (milliers FCFA)};
\draw[->,thick] (45,0) -- (45,20) node[above left]{\footnotesize Effectif};
\foreach \x in {50,75,...,175} \draw (\x,0.4) -- (\x,-0.4) node[below]{\footnotesize \x};
\foreach \y in {3,6,...,18} \draw (46.5,\y) -- (43.5,\y) node[left]{\footnotesize \y};
\draw[fill=popblueL,draw=popblue] (50,0) rectangle (75,12);
\draw[fill=poporangeL,draw=poporange] (75,0) rectangle (100,18);
\draw[fill=popblueL,draw=popblue] (100,0) rectangle (125,15);
\draw[fill=popblueL,draw=popblue] (125,0) rectangle (150,9);
\draw[fill=popblueL,draw=popblue] (150,0) rectangle (175,6);
\node at (62.5,6){\footnotesize 12}; \node at (87.5,9){\footnotesize 18}; \node at (112.5,7.5){\footnotesize 15};
\node at (137.5,4.5){\footnotesize 9}; \node at (162.5,3){\footnotesize 6};
\end{tikzpicture}
\end{center}
\legende{Histogramme des salaires mensuels.}
\end{popfigure}
\item Courbe cumulative croissante : points $(50;0)$, $(75;12)$, $(100;30)$, $(125;45)$, $(150;54)$, $(175;60)$. Courbe cumulative décroissante : on place en abscisse les bornes \emph{inférieures} : points $(50;60)$, $(75;48)$, $(100;30)$, $(125;15)$, $(150;6)$, $(175;0)$.
\begin{popfigure}
\begin{center}
\begin{tikzpicture}[x=0.045cm,y=0.07cm]
\draw[->,thick] (45,0) -- (182,0) node[below left]{\footnotesize Salaire (milliers FCFA)};
\draw[->,thick] (45,0) -- (45,66) node[above left]{\footnotesize Effectif cumulé};
\foreach \x in {50,75,...,175} \draw (\x,1.2) -- (\x,-1.2) node[below]{\footnotesize \x};
\foreach \y in {10,20,...,60} \draw (46.5,\y) -- (43.5,\y) node[left]{\footnotesize \y};
\draw[very thick,popblue] (50,0) -- (75,12) -- (100,30) -- (125,45) -- (150,54) -- (175,60);
\draw[very thick,poporange] (50,60) -- (75,48) -- (100,30) -- (125,15) -- (150,6) -- (175,0);
\foreach \p in {(50,0),(75,12),(100,30),(125,45),(150,54),(175,60)} \fill[popblue] \p circle (0.7);
\foreach \p in {(50,60),(75,48),(125,15),(150,6),(175,0)} \fill[poporange] \p circle (0.7);
\draw[dashed,popdark] (100,30) -- (100,0);
\node[above right,popblue] at (130,50){\footnotesize ECC};
\node[above right,poporange] at (130,13){\footnotesize ECD};
\node[below,popdark] at (100,-3){\footnotesize };
\end{tikzpicture}
\end{center}
\legende{Courbes cumulatives croissante (ECC) et décroissante (ECD).}
\end{popfigure}
\item \begin{enumerate}
\item Les deux courbes se coupent au point d'abscisse $100$ (et d'ordonnée $30=\dfrac{N}{2}$). Le salaire médian est donc $Me \approx \num{100000}$ FCFA : la moitié des employés gagne moins de \num{100000} FCFA, l'autre moitié gagne au moins \num{100000} FCFA.
\item Graphiquement, l'ordonnée de la courbe des ECC en $100$ est $30$, soit $\dfrac{30}{60}=50\,\%$. Vérification par le calcul : $12+18=30$ employés, et $\dfrac{30}{60}\times 100 = 50\,\%$. Donc \textbf{$50\,\%$} des employés gagnent moins de \num{100000} FCFA.
\end{enumerate}
\item « Moins de \num{150000} FCFA » correspond à l'ECC de la classe $[125 ; 150[$, soit $54$ employés. Or $\dfrac{54}{60}\times 100 = 90\,\%$. Comme $90\,\% \geq 75\,\%$, l'affirmation du directeur est \textbf{exacte} (il y a même $90\,\%$ des employés dans ce cas).
\end{enumerate}
\textbf{Points de vigilance} : pour la courbe décroissante, placer les ECD aux bornes \emph{inférieures} des classes ; la médiane graphique se justifie par l'intersection des deux courbes (ou par l'ordonnée $\frac{N}{2}=30$ sur la courbe croissante) ; toute affirmation doit être validée par un calcul de pourcentage explicité.
\end{corrige}

\begin{corrige}[Exercice 11 — Comparaison de deux écoles et rapport — Type Probatoire]
\textbf{Barème indicatif (sur 10)} : question 1 : 2 pts ; question 2 : 2 pts ; question 3 : 1 pt ; question 4 : 1,5 pt ; question 5 : 1,5 pt ; question 6 : 2 pts.
\begin{enumerate}
\item \textbf{Lycée A} ($N=100$) : les fréquences s'obtiennent directement : $10\,\%$ ; $30\,\%$ ; $45\,\%$ ; $15\,\%$.\\
\textbf{Lycée B} ($N=80$) : $\dfrac{4}{80}=5\,\%$ ; $\dfrac{20}{80}=25\,\%$ ; $\dfrac{40}{80}=50\,\%$ ; $\dfrac{16}{80}=20\,\%$.
\begin{center}
\begin{tabular}{|l|c|c|c|c|}
\hline
Note & $[0 ; 5[$ & $[5 ; 10[$ & $[10 ; 15[$ & $[15 ; 20]$ \\
\hline
Fréquence A (\%) & 10 & 30 & 45 & 15 \\
\hline
Fréquence B (\%) & 5 & 25 & 50 & 20 \\
\hline
\end{tabular}
\end{center}
\item Fréquences cumulées croissantes :
\begin{center}
\begin{tabular}{|l|c|c|c|c|}
\hline
Note inférieure à & $5$ & $10$ & $15$ & $20$ \\
\hline
FCC A (\%) & 10 & 40 & 85 & 100 \\
\hline
FCC B (\%) & 5 & 30 & 80 & 100 \\
\hline
\end{tabular}
\end{center}
\item Les deux lycées n'ont pas le même effectif total ($100$ contre $80$) : comparer les effectifs bruts n'a pas de sens, car un même effectif ne représente pas la même proportion. Il faut comparer les \textbf{fréquences} (pourcentages), qui ramènent chaque série à une base commune de $100$.
\item \begin{enumerate}
\item « Au moins $10/20$ » : lycée A : $45+15=60\,\%$ ; lycée B : $50+20=70\,\%$.
\item « Au moins $15/20$ » : lycée A : $15\,\%$ ; lycée B : $20\,\%$.
\end{enumerate}
\item Classes de même amplitude ($5$ points) : hauteurs proportionnelles aux fréquences.
\begin{popfigure}
\begin{center}
\begin{tikzpicture}[x=0.5cm,y=0.09cm]
\draw[->,thick] (0,0) -- (21.5,0) node[below left]{\footnotesize Note};
\draw[->,thick] (0,0) -- (0,55) node[above left]{\footnotesize Fréquence (\%)};
\foreach \x in {0,5,10,15,20} \draw (\x,1) -- (\x,-1) node[below]{\footnotesize \x};
\foreach \y in {10,20,...,50} \draw (0.3,\y) -- (-0.3,\y) node[left]{\footnotesize \y};
\draw[fill=popblueL,draw=popblue] (0,0) rectangle (5,10);
\draw[fill=popblueL,draw=popblue] (5,0) rectangle (10,30);
\draw[fill=poporangeL,draw=poporange] (10,0) rectangle (15,45);
\draw[fill=popblueL,draw=popblue] (15,0) rectangle (20,15);
\node at (2.5,5){\footnotesize 10}; \node at (7.5,15){\footnotesize 30}; \node at (12.5,22.5){\footnotesize 45}; \node at (17.5,7.5){\footnotesize 15};
\end{tikzpicture}
\end{center}
\legende{Histogramme des fréquences — Lycée A.}
\end{popfigure}
\begin{popfigure}
\begin{center}
\begin{tikzpicture}[x=0.5cm,y=0.09cm]
\draw[->,thick] (0,0) -- (21.5,0) node[below left]{\footnotesize Note};
\draw[->,thick] (0,0) -- (0,55) node[above left]{\footnotesize Fréquence (\%)};
\foreach \x in {0,5,10,15,20} \draw (\x,1) -- (\x,-1) node[below]{\footnotesize \x};
\foreach \y in {10,20,...,50} \draw (0.3,\y) -- (-0.3,\y) node[left]{\footnotesize \y};
\draw[fill=popgreenL,draw=popgreen] (0,0) rectangle (5,5);
\draw[fill=popgreenL,draw=popgreen] (5,0) rectangle (10,25);
\draw[fill=poporangeL,draw=poporange] (10,0) rectangle (15,50);
\draw[fill=popgreenL,draw=popgreen] (15,0) rectangle (20,20);
\node at (2.5,2.5){\footnotesize 5}; \node at (7.5,12.5){\footnotesize 25}; \node at (12.5,25){\footnotesize 50}; \node at (17.5,10){\footnotesize 20};
\end{tikzpicture}
\end{center}
\legende{Histogramme des fréquences — Lycée B.}
\end{popfigure}
\item \textbf{Rapport à l'intention du jury} :\\
Les effectifs totaux étant différents ($100$ candidats pour le lycée A, $80$ pour le lycée B), la comparaison s'effectue sur les fréquences. Le lycée B présente de meilleures performances globales : $70\,\%$ de ses candidats ont obtenu au moins $10/20$, contre $60\,\%$ pour le lycée A, et $20\,\%$ ont obtenu au moins $15/20$, contre $15\,\%$. De plus, seulement $5\,\%$ des candidats du lycée B ont une note inférieure à $5/20$, contre $10\,\%$ au lycée A. La classe modale est la même dans les deux lycées ($[10 ; 15[$), mais elle est plus marquée au lycée B ($50\,\%$ contre $45\,\%$). En conclusion, le lycée B a globalement obtenu les meilleurs résultats à ce concours.
\end{enumerate}
\textbf{Points de vigilance} : exiger la justification de la comparaison par les fréquences (effectifs totaux différents) ; « au moins $10$ » signifie $\geq 10$, donc classes $[10;15[$ et $[15;20]$ ; le rapport doit comporter des arguments chiffrés et une conclusion claire.
\end{corrige}

\newpage

\coursec{Fiche bilan}

\begin{popfigure}
\begin{tikzpicture}[
  mindmap,
  grow cyclic,
  every node/.style={concept, execute at begin node=\small\bfseries},
  concept color=popblue,
  level 1/.append style={level distance=4.5cm, sibling angle=72},
  level 2/.append style={level distance=3cm, sibling angle=45},
  branche1/.style={concept color=popblue},
  branche2/.style={concept color=poporange},
  branche3/.style={concept color=popgreen},
  branche4/.style={concept color=poppurple},
  branche5/.style={concept color=poppink},
]
\node[root concept, align=center]{Séries groupées\\en classes}
  child[branche1] { node[align=center]{Regroupement\\en classes}
    child { node {Effectif total $N$} }
    child { node {Amplitude $A$} }
    child { node {Nombre classes $k$} }
    child { node {Largeur $h$} }
    child { node {Bornes $[a_i;b_i[$} }
  }
  child[branche2] { node[align=center]{Tableau\\statistique}
    child { node {Classes $C_i$} }
    child { node {Effectifs $n_i$} }
    child { node {Fréquences $f_i$} }
    child { node {Eff. cumulés $N_i$} }
    child { node {Freq. cumulées $F_i$} }
  }
  child[branche3] { node {Histogramme}
    child { node {Aire = Effectif} }
    child { node {Hauteur $= n_i/h$} }
    child { node {Densité $= f_i/h$} }
    child { node {Classes inégales OK} }
  }
  child[branche4] { node[align=center]{Courbes\\cumulées}
    child { node {Croissante $(x_i; F_i)$} }
    child { node {Décroissante $(x_i; 1-F_{i-1})$} }
    child { node {Lecture $M_e, Q_1, Q_3$} }
  }
  child[branche5] { node {Caractéristiques}
    child { node {Position: $M_e, Q_1, Q_3, Mo$} }
    child { node {Dispersion: $E, EIQ$} }
    child { node {Moyenne $\bar x$ approx.} }
    child { node {Écart-type $\sigma$ approx.} }
  };
\end{tikzpicture}
\legende{Carte mentale : Séries statistiques regroupées en classes}
\end{popfigure}

\begin{bilan}
\courssub{Définitions clés}
\begin{itemize}
  \item \textbf{Population statistique} : ensemble des individus étudiés. Taille $N$ (effectif total).
  \item \textbf{Caractère quantitatif continu} : valeurs numériques dans un intervalle (ex. hauteur, tension, temps).
  \item \textbf{Classe $C_i = [a_i ; b_i[$} : intervalle de regroupement. \cle{Borne inférieure} $a_i$, \cle{borne supérieure} $b_i$. Classes contiguës, exhaustives, mutuellement exclusives.
  \item \textbf{Milieu de classe} $x_i = \dfrac{a_i + b_i}{2}$ (représentant de la classe).
  \item \textbf{Largeur (pas) de classe} $h_i = b_i - a_i$. Souvent constant $h$.
  \item \textbf{Effectif $n_i$} : nombre d'individus dans $C_i$. $\sum n_i = N$.
  \item \textbf{Fréquence $f_i = \dfrac{n_i}{N}$} ($0 \le f_i \le 1$, $\sum f_i = 1$).
  \item \textbf{Effectif cumulé croissant} $N_i = \sum_{j=1}^{i} n_j$. \textbf{Décroissant} $N'_i = \sum_{j=i}^{k} n_j$.
  \item \textbf{Fréquence cumulée croissante} $F_i = \dfrac{N_i}{N}$. \textbf{Décroissante} $F'_i = \dfrac{N'_i}{N}$.
\end{itemize}

\courssub{Formules à connaître par cœur}
\begin{center}
\begin{tabularx}{\linewidth}{|l|X|l|}\hline
\textbf{Notion} & \textbf{Formule / Règle} & \textbf{Contexte} \\\hline
Largeur classe (égale) & $h = \dfrac{\text{Étendue}}{k} = \dfrac{x_{\max}-x_{\min}}{k}$ & Choix $k \approx \sqrt{N}$ (Sturges) \\\hline
Milieu de classe & $x_i = \dfrac{a_i+b_i}{2}$ & Calculs approx. (moyenne, $\sigma$) \\\hline
Fréquence & $f_i = \dfrac{n_i}{N}$ & Tableau, histogramme (densité) \\\hline
Eff. cumulé croissant & $N_i = N_{i-1} + n_i$ ($N_0=0$) & Courbe croissante, médiane \\\hline
Freq. cumulée croissante & $F_i = \dfrac{N_i}{N}$ & Courbe croissante (ogive) \\\hline
Eff. cumulé décroissant & $N'_i = N'_{i+1} + n_i$ ($N'_{k+1}=0$) & Courbe décroissante \\\hline
Densité histogramme & $d_i = \dfrac{n_i}{h_i}$ ou $\dfrac{f_i}{h_i}$ & \textbf{Aire du rectangle} $= n_i$ (ou $f_i$) \\\hline
Médiane (classe $m$) & $M_e = a_m + \dfrac{\frac{N}{2} - N_{m-1}}{n_m} \times h_m$ & $N_{m-1} < \frac{N}{2} \le N_m$ \\\hline
Quartiles ($Q_1, Q_3$) & $Q_\alpha = a_q + \dfrac{\alpha N - N_{q-1}}{n_q} \times h_q$ & $\alpha = 0,25$ ou $0,75$ \\\hline
Moyenne approx. & $\bar x \approx \dfrac{\sum n_i x_i}{N}$ & $x_i$ milieux de classes \\\hline
Variance approx. & $\sigma^2 \approx \dfrac{\sum n_i x_i^2}{N} - \bar x^2$ & $x_i$ milieux de classes \\\hline
\end{tabularx}
\end{center}

\courssub{Méthodes express}
\begin{methode}[Construire les classes]
  \begin{enumerate}
    \item Calculer l'étendue $E = x_{\max} - x_{\min}$.
    \item Choisir $k$ (nombre de classes) : $5 \le k \le 15$, souvent $k \approx \sqrt{N}$.
    \item Calculer $h = E/k$ (arrondir à une valeur « propre » : 1, 2, 5, 10, 0,5...).
    \item Fixer la borne inférieure de la 1ère classe $a_1 \le x_{\min}$ (multiple de $h$).
    \item Définir les classes $[a_1; a_1+h[$, $[a_1+h; a_1+2h[$, ...
  \end{enumerate}
\end{methode}

\begin{methode}[Remplir le tableau complet]
  \begin{enumerate}
    \item Compter $n_i$ pour chaque classe (relevé).
    \item Calculer $f_i = n_i/N$ (souvent en $\%$).
    \item Calculer $N_i$ cumulés croissants (ajouter $n_i$ au précédent).
    \item Calculer $F_i = N_i/N$ cumulés croissants.
    \item (Optionnel) Cumulés décroissants : partir de la fin.
  \end{enumerate}
\end{methode}

\begin{methode}[Tracer l'histogramme (aire = effectif)]
  \begin{enumerate}
    \item Axe des abscisses : bornes des classes (échelle réelle).
    \item Axe des ordonnées : \textbf{densité} $d_i = n_i / h_i$ (ou $f_i/h_i$).
    \item Tracer rectangle sur $[a_i; b_i]$ de hauteur $d_i$. Aire $= d_i \times h_i = n_i$.
    \item \textbf{Attention} : si classes inégales, hauteur $\neq$ effectif !
  \end{enumerate}
\end{methode}

\begin{methode}[Tracer les courbes cumulées (ogives)]
  \begin{enumerate}
    \item Points croissants : $(b_i ; F_i)$ (borne sup, freq. cum. croiss.). Relier par segments.
    \item Points décroissants : $(a_i ; F'_i)$ (borne inf, freq. cum. décroiss.). Relier par segments.
    \item Lecture graphique : $M_e$ à l'abscisse de $y=0,5$ ; $Q_1$ à $y=0,25$ ; $Q_3$ à $y=0,75$.
  \end{enumerate}
\end{methode}
\end{bilan}

\begin{aretenir}[Checklist avant le Probatoire]
\begin{itemize}
  \item $\square$ Savoir définir : population, caractère, modalité, classe, effectif, fréquence, effectif/fréquence cumulés (croissant/décroissant).
  \item $\square$ Savoir choisir le nombre de classes $k$ et calculer la largeur $h$ (arrondi pratique).
  \item $\square$ Construire un tableau complet : Classes $|$ Effectifs $|$ Fréquences $|$ Eff. cum. $|$ Fréq. cum.
  \item $\square$ Calculer les milieux de classe $x_i = (a_i+b_i)/2$.
  \item $\square$ Tracer un histogramme \textbf{aire = effectif} (calculer la densité $n_i/h_i$ si classes inégales).
  \item $\square$ Tracer les deux courbes cumulées (croissante aux bornes sup, décroissante aux bornes inf).
  \item $\square$ Lire graphiquement la médiane $M_e$, $Q_1$, $Q_3$ sur la courbe cumulative croissante.
  \item $\square$ Calculer $M_e$, $Q_1$, $Q_3$ par la formule d'interpolation linéaire (classe médiane/quartile).
  \item $\square$ Calculer la moyenne $\bar x$ et l'écart-type $\sigma$ approximatifs à partir des milieux $x_i$.
  \item $\square$ Interpréter les résultats : position (centre), dispersion (étalement), forme (symétrie/asymétrie).
  \item $\square$ Justifier chaque étape par une phrase (règle de trois, proportionnalité, définition).
  \item $\square$ Vérifier la cohérence : $\sum n_i = N$, $\sum f_i = 1$, $F_{\text{dernière}} = 1$, $N_{\text{dernier}} = N$.
\end{itemize}
\end{aretenir}

\findecours

\end{document}
